\documentclass[12pt]{book}

\usepackage{graphicx}
\usepackage{eso-pic}
\usepackage{tcolorbox}
\usepackage{polyglossia}
\setdefaultlanguage[indentfirst=false,forceheadingpunctuation=false]{russian}
\setotherlanguages{english}

\RequirePackage[a4paper, headsep = 0.5 \headsep, left=2.5cm, right=2.1cm, top=2cm, bottom=2.1cm]{geometry}

\setmainfont{URW Bookman}
\newfontfamily\cyrillicfont{URW Bookman}[Script=Cyrillic, Ligatures=TeX]
\setlength{\parindent}{20pt}

\usepackage{enumitem}
\setlist[itemize, 1]{label =\raisebox{-0.3\height}{\scalebox{1.8}{\textbullet}}}

\title{Принимая отчужденность}
\author{Тодд МакГоуэн}
\date{2019}
\begin{document}

\AddToShipoutPictureBG*{\includegraphics[width=\paperwidth]{title.png}}

\null
    \thispagestyle{empty}%
    \addtocounter{page}{-1}%
    \newpage

\tableofcontents

\chapter{Введение}

''Отчужденность'' звучит непривлекательно. Будто что-то разрушительное, чего лучше избегать. Никто не говорит друзьям ''Давайте отчуждаться'', наверняка это не так приятно, как основание сообщества, романтические отношения, общение с приятелями, пицца, игра в теннис или выпивка. Напротив --- мы хотим помочь подростку избежать отчужденности, а отчужденному постороннему --- почувствовать себя уютней. Левые, согласно Карлу Марксу, понимают политическую деятельность как борьбу с отчуждением. Почти всем отчужденность видится проблемой, которую надо решить.

Приговор отчужденности поспешен и забывает о ее освободительном качестве. Исторически, у отчужденности плохая репутация, но ведь именно отчужденность дает возможность выхода за пределы конкретной ситуации. Следовало бы думать о ней гораздо лучше. Увидеть ее достоинства и постараться усилить ее опыт.

Никто открыто не отстаивал отчужденность как проект, но некоторые принимают ее основой свободы. Немецкий идеализм, кульминацией которого стала философия Гегеля, говорит об отчужденности как условии самосознания. Психоанализ воспринимает отчужденность субъекта как непреклонный факт нашего существования и стремится примирить с ним. Экзистенциализм Сартра и Симоны де Бовуар тоже видит в отчужденности субъекта основу свободы. Франц Фанон объединяет все три течения --- Гегеля, психоанализ и экзистенциализм --- и формулирует антиколониальную теорию, выступая против попыток избегать фундаментальной нам всем отчужденности. Впрочем, эти мыслители никогда не используют отчужденность как основное понятие или лозунг, хотя оно объединяет эти разрозненные попытки говорить о свободе. В этой книге отчужденность выдвигается на первый план.

Отчужденность парадоксальна. Она не значит отклонения от исходной идентичности, от того, что мы когда-то имели но потеряли. Отчужденность неизменно присуща нам. Она не отдаленность от первичного, потому что первична. Мы отчуждены в своей субъектности до самой смерти. Смерть --- единственное лекарство от этой болезни\footnote{В момент смерти наша отчужденность проходит --- мы становимся совершенно чуждыми. Смертью субъект совершенно теряется в инаковости, уже не сохраняя ни малейшей частицы себя.}, так что, наверное, не такая уж это и болезнь. Отчужденность --- не то, о чем можно сожалеть или чего лучше избегать, ведь благодаря ней возможно то, что делает существование достойным продолжения --- свобода, равенство, солидарность, способность выходить за пределы ситуации.

Все творение зависит от отчужденности, она выкорчевывает нас из определенности, освобождает от конкретной ситуации. Каждый рождается в специфичном социуме, в котором на него действуют мощные силы, но не сводится к ним. Первичное отчуждение в языке дистанцирует как от других, так и внутренне: к себе теперь можно относиться как к другому. Мы отчуждены от себя и от других, и что бы мы ни делали, нам никогда не преодолеть эту отчужденность. Она основа нашего существования.

Отчужденность --- это отсутствие самоидентичности. Она отдаляет субъекта от условий, в которых он проявляется. Если я не тождествен себе, если я в противоречии с собой, то внешние силы не могут меня полностью определить. Существо, дистанцированное от самого себя, никогда не сможет быть тем, во что его толкают внешние силы. Этот внутренний раскол значит сопротивление детерминированию биологией или обществом. Именно дистанция от самой себя позволяет субъектности действовать против определяющего. Социальное и натуралистическое утрачивают свою власть на субъектом, потому что его отчужденность делает их незначительными.

Я могу не оставаться до конца своих дней жителем маленького американского городка на Среднем Западе, могу бросить вызов определяющим факторам этой ситуации, и все потому что никогда не отождествлен с положением в котором родился или нахожусь. Это дает мне возможность отнестись к ситуации критически, вплоть до отказа от всех идей и ценностей, которые возложило на меня общество. Я могу быть воспитан с энтузиазмом к капиталистическому строю, но стать марксистом. Если меня заставляют быть женственным, я могу стать мужественным. Полюбив Средний Запад Америки, я, возможно, стану работать в Аргентине. Моя неспособность быть именно тем, кто я есть --- моя не-однозначность --- вот что позволяет мне выйти из исходной ситуации. И даже если я остаюсь в ней и принимаю ее, я все равно отношусь к ней как отчужденный. Можно быть чужим и у себя дома. У меня, как и у всех, есть внутренняя чужеродность, мешающая прямому усвоению диктата общества.

Без этой внутренней дистанции человек не был бы субъектом, не смог бы соотноситься с собой. Субъектность порождена и определена этим расколом. Отделенность от самих себя делает нас нами --- неотделимыми от чуждого в нас\footnote{Славой Жижек в \textit{Добавочном Наслаждении} подчеркивает, что отчуждение конститутивно. Мы не целостны изначально, чтобы \textit{затем} отчуждаться. Мы возвращаемся к изначальному через отчуждение. Поэтому ''вне отчужденности никакой жизни нет, нет никакой позитивной основы отчужденности'' (Slavoj Žižek, \textit{Surplus Enjoyment: A Guide for the Non-Perplexed} (London: Bloomsbury, 2022), 39). Верить, что мы были до и вне отчуждения --- ностальгировать по истоку, которого никогда не существовало.}. Чуждое дает нам нашу уникальность.

Эта внутренняя чуждость и есть то, что Зигмунд Фрейд назовет бессознательным\footnote{Для описания своей концепции Фрейд умышленно выберет \textit{unbewußt}, а не распространенный в то время \textit{bewußtlos}. \textit{Bewußtlos} значил бы ''без сознания'', \textit{unbewußt} же значит противоположное сознанию, чуждое его структуре.}. Бессознательное ставит барьер как между субъектом и ситуацией, так и внутри субъекта. Только отчужденные обладают бессознательным. Вернее, чуждый статус бессознательного значит, что это оно владеет нами. Это иное внутри диктует субъекту его действия, направляет усилия и создает убеждения. Бессознательное ориентирует субъектность, потому что выражает ее самораскол\footnote{В психике Фрейд даст бессознательному структурный приоритет. Ускользая от контроля, оно гораздо больше влияет на нас, чем сознательное мышление: ''Бессознательное --- вот истинная психическая реальность'' (Sigmund Freud, \textit{The Interpretation of Dreams}, trans. James Strachey, in The Complete Psychological Works of Sigmund Freud, vol. 5, ed. James Strachey (London: Hogarth Press, 1953 [1900]), 651). Фрейд написал это на заре своей теории, впоследствии он откажется от некоторых своих идей, но никогда не усомнится в приоритете бессознательного.}.

Отчуждение субъектов от себя видно, когда бессознательное заставляет их идти наперекор своим же интересам. Кто-то объедается пончиками и кексами, вместо того, чтобы есть яблоки и брокколи и предотвратить ишемическую болезнь. Другой вместо комфортного, неконфликтного партнера выбирает сложного, болезненного человека, с которым постоянно ссорится. Продолжает искать и употреблять все новые дурманящие вещества, хотя прекрасно понимает, как они опасны и вредны. Хочет профессионального роста, но раз за разом выступает против начальника. Повсюду субъекта поджидают саморазрушительные, ускользающие от сознательного контроля действия.

Каждый сосуществует с идущей вразрез его желаниям силой. Каждый видит, как сам мешает достижению своих целей, как не может контролировать себя. Отсутствие контроля изнурительно, но также основа способности субъекта становиться иным. Верховенство чуждого --- бессознательного --- открывает возможность свободы. Свобода напрямую зависит от подрыва сознательных желаний бессознательными импульсами.

О бессознательном обычно думают как о принуждающем, забирающем свободу. Мы не обдумываем бессознательные действия и не чувствуем, что они выбраны нами. Но свобода --- не плод сознательных размышлений. Это прерывание решимости, определенности. И именно это делает бессознательное. Оно перебивает социальные и психологические определяющих факторы, открывает разрыв в необусловленную ими альтернативную возможность. Избегая сознательного контроля, бессознательное становится точкой свободы субъекта. Мы свободны благодаря своей раздвоенности.

Самотождественности лишена не только субъектность. Ничто не самотождественно, иначе не смогло бы меняться во времени и перемещаться в пространстве\footnote{Жан-Поль Сартр утверждает, что только субъекты отчуждены, а остальные существа тождественны себе. Он радикально различает субъекта, который для-себя, и остальное бытие, которое в-себе: ''Бытие есть. Бытие --- в-себе. Бытие есть то, что оно есть.'' (Jean-Paul Sartre, \textit{Being and Nothingness}, trans. Hazel E. Barnes (New York: Washington Square Press, 1956 [1943]), 29). Сартр здесь не задумывается, как субъект мог стать отчужденным. Будь бытие лишь тем, что оно есть, субъект не возник бы, не проявился из такого чистого царства самотождественности.}. Это касается и животных и растений, и гор и рек. Будь сущность тем, что она есть, она бы никогда не изменилась. Субъект же превращает свою нетождественность в определяющую его структуру, такой процесс и есть отчуждение.

Первое свидетельство отчуждения --- имя. Ни один ребенок не дает его себе сам\footnote{Мама говорит, что я сам выбрал имя: она положила в шляпу несколько листочков бумаги с именами, и я схватил листочек с именем ''Тодд''. Такой метод не избегает навязывания, ведь все равно контролируется взрослым: он выбирает возможные имена, вырезает и кладет в шляпу листочки с ними. Заключенное в присвоении имени отчуждение неизбежно. Затем мама призналась, что ''Тодд'' было написано на каждом из листочков. Впрочем, признание могло быть следствием тогда уже развившейся у нее болезни Альцгеймера, но так или иначе, оно показывает определяющую роль внешнего в присвоении имени.}. Происходит акт насилия, родители или другие взрослые навязывают ребенку произвольное имя. И даже если он потом воспринимает его своим, оно остается чуждым, сошедшим на него извне. Родители могут называть мальчика Эрнестом и надеется, что он вырастет честным и прямодушным. Могут называть девочку Грейс и думать, что она будет приличной и изящной. Иногда используют имена художников, таких как Лэнгстон или Дэшил, пытаясь ассоциацией приобщить детей к их гениальности. Дают имя, связанное с семейной историей, чтобы оно напоминало ребенку о его происхождении, определяло его поведение. Но и ничего не значащее имя остается чуждым названному означающим. Он не выбирал его, и воспринимает его своим только через отчуждение.

Имя это наша дистанция до себя. В каждом языке есть способ идентифицировать себя именем: я могу сказать ''Меня зовут Тодд'', или ''Ich heiße Todd'', или ''Me llamo Todd''. Всякое такое высказывание провозглашает тождество, и этим выявляет внутреннюю дистанцию между говорящим и именем. Так оно выдает, что я и Тодд --- не одно и то же\footnote{В логике первого порядка утверждение тождества исключает различие, потому что различие значило бы категориальную ошибку: смешение идентификации и предикации. На этом замечании Бертран Рассел строит свою критику Гегеля. Гегель ошибается --- утверждает он --- когда говорит об идентичности в различии, об этой фундаментальной отчужденности субъекта, ведь, например, ''Скотт --- автор ''Уэверли'''' предполагает скорее идентификацию, чем предикацию. Критику этой позиции Рассела и защиту Гегеля см. в \textit{Emancipation After Hegel: Achieving a Contradictory Revolution} (New York: Columbia University Press, 2019).}. Самотождественное существо не смогло бы так сказать.

Даже если я переименую себя, имя останется чуждым. Казалось бы, я называю себя сам, не полагаюсь на данное другими. Разве этого недостаточно для преодоления отчужденности? Нет. Называя себя иначе, я действую исходя их доступных возможностей: выбираю из используемых другими имен или комбинирую их. Выбор имени всегда предполагает долг перед внешним. Мое новое имя не только принадлежит другим, но и выбрано с учетом их впечатления. Откуда его не возьми, оно несет печать самораскола.

Неизбежность отчужденности дает некоторую уверенность. И вместо того, чтобы страдать от саморазделения до самой смерти, мы можем манипулировать им и даже усиливать его. Оставление тщетных попыток преодолеть отчужденность, ее принятие, трансформирует ее. Из бремени она становится освобождением от определенности. Переосмысление отчужденности --- важнейшая политическая задача, особенно в эпоху вездесущих схем его якобы преодоления.

Об отчужденности плохо думают, потому что воображают скрывающуюся за ней завершенность. Где-то в затерянном прошлом, в Эдеме, или в утопическом будущем, нас ждет самотождественность, полнота, совершенство. Она всегда ускользает и тем соблазняет еще больше. Такое лишенное всякого отчуждения состояние очень привлекает нас, но лишь пока остается воображаемым. Малейшая попытка его очертить, конкретизировать, воплотить, выдает его полную несостоятельность\footnote{Никто не считает \textit{Возвращенный рай} лучшей поэмой, чем \textit{Потерянный рай}. Отчужденность приятна, потому что ее неприятно пытаться преодолевать. Мильтон невольно демонстирует это разницей в качестве поэм и --- более того --- изображает неотчужденных Бога, Христа, ангелов в \textit{Потерянном раю} совсем не такими привлекательными, как отчужденного Сатану.}.

Нам следует мыслить отчужденность вне ассоциаций с полнотой или самотождественностью. Мы отчуждены, но не от чего-то могущего дополнить нас. Мы и не теряли ничего такого. Помыслив нехватку вне связи с целостностью, мы поменяем политическую значимость отчужденности. Отчужденность не определяется через не-отчужденность.

Отчужденность --- источник всех человеческих страданий, но оно и источник всякого удовлетворения. Отчужденность предоставляет пространство действию вопреки, без него субъект неотличим от сил, породивших его. Заниматься сексом, есть торт или обсуждать немецкий идеализм --- ничего из этого не возможно без отчужденности, потому что без нее нет субъекта.

Если что и отличает говорящих существ от прочих, так это их осознание отсутствия самоидентичности. Через язык отчужденность осознается субъектами неизменно присущей им. Язык создает дистанцию, и это меняет отношение существа к отчужденности. Субъекты означающего могут осознать и принять себя отчужденными, могут найти удовлетворение в этом отсутствии самотождественности, и только отказ рассматривать отчужденность конститутивной составляющей мешает этому.

Чаще мы думаем о своих экзистенциальных и политических проектах как о преодолевающих отчужденность. За редкими исключениями, мы стремимся дать отчужденности бой, искоренить ее. Психотерапевты воображают, что помогают пациентам обрести самоидентификацию. Революционеры стремятся построить общество, в котором жили бы в гармонии с другими. Экоактивисты работают на формой существования, которая была бы в гармонии с природой\footnote{Впрочем, не всегда. Например, эколог Мюррей Букчин, будучи гегельянцем, говорит, что для экологического проекта необходимо принять отчужденность. См. Murray Bookchin, \textit{The Philosophy of Social Ecology: Essays on Dialectical Naturalism} (Chico: AK Press, 2022 [1990]).}. Разительное разнообразие проектов не затрагивает их сути --- избавления от чуждого. Отчужденность непрестанно видится тем, с чем следует бороться, хотя сама борьба возможна только благодаря ней.

Эта книга предлагает совершенно иной взгляд на отчужденность: усилия в ее преодолении значат не радикальный ответ ситуации, а неспособность увидеть конституирующую роль отчужденности. Нам следует оставить попытки преодолеть ее. Отчужденность должна стать нашей экзистенциальной и политической программой. Отчуждение --- это освобождение.

\chapter{Беспокойный субъект}

\section{Продажная идентичность}

Субъект не то же, что присвоенная ему символическая идентичность. Невозможно понять отчуждение, не видя этого базового различия --- между субъектностью и идентичностью. Идентичность складывается из социальных позиций: работа, семейный статус, религиозные пристрастия, политические предпочтения, национальность и т.д.. Как бы легко я себя в них не видел, я никогда не отождествляюсь с ними полностью. Этот недостаток и есть субъектность. Ни одна из идентичностей не может полностью совпасть с ней, как не подгоняй.

Субъектность возникает в отчуждении, в разрыве между субъектом и принимаемой им символической идентичностью. Субъект не способен стать символической идентичностью, потому что соотнесение с ней значит дистанцию от нее. Всегда есть разрыв между мной как субъектом и мной как личностью. Само усилие, порыв соответствовать идентичности выражает расхождение с ней\footnote{Жан-Поль Сартр показывает это в своем рассказе об официанте, ведущем себя так, будто он официант. Пытаясь полностью воплотить символическую идентичность, он неизбежно терпит неудачу: ''Но посмотрите на этого официанта. Его движения быстрые и точные. Чуть слишком точные и чуть слишком быстрые. Он идет к посетителям чуть слишком быстрым шагом, наклоняется вперед чуть слишком нетерпеливо, голос и глаза его выражают чуть слишком заботливый интерес к заказу клиента'' (Jean-Paul Sartre, Being and Nothingness, trans. Hazel E. Barnes (New York: Washington Square Press, 1956 [1943]), 59). Для Сартра вера в то, что человек и есть символическая идентичность --- самообман, и даже обманщик не может преодолеть разрыв между ними. Слишком высокая точность выдает несовершенство.}. Ведь сводись субъектность к символической идентичности, самоидентификация и не понадобилась.

Субъектность неопределима. Кто мы как субъекты --- вопрос без ответа, хотя мы и пытаемся ответить на него: мужчина, женщина, китаец, итальянец, индус, мусульманин, юрист, дантист. Когда я чувствую, что теряю контроль, плотность существующего субъекта, я напоминаю себе, что я американец, или что я сын, или что я профессор --- напоминаю себе, что я есть. Но никакие усилия не снимают дистанцию между субъектностью и идентичностью\footnote{Чем больше прикладывается усилий соответствовать, тем явственней несоответствие. Усилие значит отсутствие того, на что направлено. Если бы я на самом деле соответствовал своей идентичности профессора, мне не пришлось бы стараться одеваться как профессор, говорить как профессор и ходить как профессор.}. Все попытки соответствовать проваливаются\footnote{Субъект не моряк Попай и не может сказать ''Я таков, каков я есть'' (''I Yam what I yam and that's all what I yam'').}.

Символическая идентичность и маскирует, и маркирует вопрос. Путаница неизбежна, потому что субъектности без присвоенной ей символической идентичности не бывает. Идентичность неизбежна, только она ориентирует нас в социальной ситуации.

И в то же время идентичность --- бегство от субъектности. Мы ищем передышку в борьбе за безопасность, надеемся укрыться в идентичности, но она всегда не то, что мы искали. Отношение к ней определяет наше политическое бытие.

Мы можем противостоять символической идентичности отчужденной субъектностью. Кажется, что это бесполезно, но это политический акт --- первый и самый важный, потому что позволяет действовать за пределами символического.

Погружаясь в символическую идентичность, мы избегаем этого противостояния. Избегаем политических вопросов, на которые, кажется, отвечает идентичность. Символическая идентичность значит избегание проблемы свободы. Этот момент подчеркивает Аленка Зупанчич в \textit{Что есть секс?}: ''(освободительная) политика начинается с ''утраты идентичности'', утраты, в которой нет ничего прискорбного''\footnote{Alenka Zupančič, What IS Sex? (Cambridge: MIT Press, 2017), 36}. Идентичность деполитизирует. Всякая идентичность конформна.

Какой бы трансгрессивной не казалась идентичность, субъект, когда ее принимает --- или пытается принять --- подчиняется силам установившей ее социальной власти. Бессознательное желание субъекта всегда проскакивает мимо требований социального авторитета, и в поисках одобрения субъект моделирует себя идентификацией по требуемому образцу. Бессознательное действует вне авторитета, а символическая идентичность ищет его признания. Субъект принимает символическую идентичность, чтобы ее признал социальный авторитет, поэтому она сразу уже им ограничена. И диссонанс между идентичностями не выводит за эти границы, в основе каждой из них все еще лежит социальное признание. Сколь бы радикальной не воображалась идентичность, через нее в конечном счете ищется чье-то одобрение.

Политическое различие между субъектностью и символической идентичностью разоблачается всякий раз, когда возникает скользкая, сложная социальная проблема. Вопрос про отношение к смертной казни, например, провоцирует положиться на символическую идентичность, дать заготовленный ответ, не обдумав все варианты. Принять точку зрения Католицизма и ответить как католик --- что против смертной казни. Упомянуть, что родители воспитали меня в строгом соответствии с законом Десяти Заповедей, так что смертная казнь --- выражение справедливости. Решить, что я либерал, и что, значит, смертная казнь --- слишком суровое наказание. В каждом варианте некоторая символическая идентичность помогает отставить бремя субъектности. Идентичность предлагает способы решения экзистенциальной проблемы субъектности. В любой непонятной ситуации я идентифицируюсь с определенной символической позицией, и получаю ясный ответ.

Переход от субъектности к идентичности закрывает вопросы, на которые нет готовых ответов. Идентичность ориентирует и показывает, как реагировать, поэтому мы ищем в ней убежище. Но как субъект, я должен ответить на экзистенциальные вопросы вне ее руководства. Соответствие ответа некоторой идентичности не оправдывает его.

Символическая идентичность манит в бегство от субъектности\footnote{Принятая идентичность сама может стать в тягость, если я не привык ей соответствовать.}. Поиск идентичности, которая бы полностью нам соответствовала --- это попытка вылечить отчужденность. Субъект видится сводимым к идентичности, когда она затмевает субъектность. Субъект всегда в разладе с собой, но идентичность --- даже если находится в постоянном движении --- создает мираж целостности. Мы можем оставаться в рамках вопроса о субъектности только признавая всякую идентичность неудачной. Такое признание возможно, потому что субъект никогда не соответствует закрепленной за ним идентичности. Контраст с идентичностью выводит субъекта на первый план.

Речь вносит дистанцию в говорящего. Субъектность и есть эта дистанция, эта неспособность совершенно быть собой. Отношение субъекта к самому себе опосредуется означающим, внешним. Означающее и образует форму символической идентичности, и выявляет раскол, заставляющий субъекта искать ответ на вопрос, чем он является. Всякое означающее диалектично: подчинившись ему, существо обретает способность мыслить отрицание себя. В \textit{Прочти мое желание} Джоан Копжек утверждает: ''обозначение неизбежно порождает сомнение, возможность собственного отрицания; научает думать об аннигиляции, полномасштабном разрушении всей нашей означаемой реальности''\footnote{Joan Copjec, \textit{Read My Desire: Lacan Against the Historicists} (Cambridge: MIT Press, 1994), 54–55.}. Без возможности отрицания --- возможности отринуть не только присвоенное значение, но целый означаемый универсум --- нет субъектности. Для отчужденного все данное уязвимо сомнению и может быть отвергнуто, даже утверждать свое означающее он может только через нехватку самотождественности.

Идентифицируя себя как немец, субъект показывает свою нетождественность с немецкостью. Но любая, даже самая прямая идентификация с немецкостью, только лишний раз демонстрирует ее чуждость субъекту. Если бы немецкость действительно не отличалась от него, акт идентификации был излишним и --- более того --- невозможным. Утверждение повторения свидетельствует о различии. Обычно мы не осознаем, как самоидентификация означающими самодистанцирует. Мы используем чужие слова, чтобы описать себя, никто не изобретает означающие. Акт утверждения совпадения --- признак расхождения. Это печать отчужденности.

Марин Ле Пен, лидер французской партии ''Национальное объединение'' заинтересована соответствовать тому, что значит быть француженкой. Для нее Франция --- католическая страна, и должна быть экономически независимой. Традиции Франции, включая язык, отделяют ее от остальной Европы и от мусульманских стран. Но как бы Ле Пен не заботилась о своей французской идентичности, она отличается от нее. Сама ее артикуляция этой идентичности вносит зазор. Никто не привержен традициям Франции так, как Ле Пен, но даже ей приходится трудиться, чтобы укоренить себя в своей французскости.

Самоописание раскрывает самоудаленность. Если бы субъект был идентичен французскости, никакое самоописание не потребовалось бы. Означающее идентичности указывает, что субъект не тот, кем себя называет, что должен отыгрывать свою идентичность, а не просто быть ею. Идентичность субъекта всегда вне его самого. Повторяющимися актами идентификации он принимает и поддерживает ее. Никто не получает идентичность бесплатно.

Множество научно-фантастических фильмов фокусируются на различии между субъектом и его идентичностью. В \textit{Бегущем по лезвию} (Ридли Скотт, 1981), \textit{Вспомнить все} (Пол Верховен, 1990) и \textit{Темный город} (Алекс Пройас, 1998) мы видим, как эта дистанция проявляется в повествовании\footnote{В \textit{Бегущем по лезвию} Рик Декард (Харрисон Форд) --- охотник за головами --- преследует беглых репликантов, искусственно созданных существ, которым для приручения вживили воспоминания. Но понимая, что его собственная идентичность столь же ложна, как идентичности репликантов, Декард освобождается от нее. Он, может быть, тоже репликант, но даже если не так --- в символической идентичности нет надежности. Фильм заканчивается тем, что Декард вступает в романтические отношения с женщиной, зная, что она репликант.}. Персонажи обнаруживают, что самое сокровенное --- воспоминания --- навязано им внешней [злой] силой. Символическая идентичность оказывается чужда субъектности. Дуг Куэйд (Арнольд Шварцнегер) в ''Вспомнить все'' узнает, что раньше был Карлом Хаузером (также Арнольд Шварцнегер) и создал свою нынешнюю личность, чтобы внедриться в восстание на планете Марс и уничтожить его. Все, что он знает о себе --- о своей личности --- ложь, плод манипуляций Хаузера. Личность оказывается внешней манипуляцией. В конце фильма Куэйд вырывается из-под контроля Хаузера, но никакой личности, кроме имплантированной Хаузером у него не остается. Не остается символической идентичности, на которую можно всецело положиться. В Куэйде проявляется то, что касается всякого субъекта, и потому этот персонаж столь притягателен.

Конечно, никто не внедряет в нас фальшивые воспоминания, но, подобно Куэйду, мы не можем полностью отождествиться со своей идентичностью\footnote{Эти фильмы почти всегда показывают отчужденность субъекта через паранойю: внешнее [зло] навязывает ложную идентичность, и тем самым создает раскол в субъекте. Но раскол возникает гораздо раньше и естественней --- на этапе включения субъекта в сферу символического.}. Подобно нему, мы сталкиваемся со своей субъектностью, идущей вразрез с идентичностью. Разве что Куэйд, в отличие от зрителей, не может просто дистанцироваться от своей идентичности, а должен отвергнуть ее как выдумку. Куэйд показывает отчужденную субъектность без сколь-нибудь надежной символической идентичности. Он парадигмальный субъект, совершенно отчужденный, без возможности уткнуться в символическую идентичность. Если мы не как Дуг Куэйд, то тем хуже для нас\footnote{Куэйд, когда встречает лидера восстания Марса --- мутанта Куато, способного читать мысли --- спрашивает, кто он без истинных воспоминаний. На что Куато, как прилежный ученик Сартра, отвечает: ''Ты то, что ты делаешь. Определяют действия, а не память''. Куэйд удовлетворяется таким ответом, и продолжает играть важную роль в революции.}.

Внешний статус идентичности --- отсутствие непосредственной самоидентичности --- порождает бессознательное, действующее вне контроля субъекта. Бессознательное проистекает из разрыва между субъектом и его идентичностью --- между тем, кто он, и тем, что он. Раскол между субъектностью и символической идентичностью соответствует расколу между бессознательным и сознанием. Не всякая символическая идентичность сознательна, но всякая создана в попытке уйти от бессознательного, от чуждого в субъектности.

\section{Аппетит к разрушению}

Есть разрыв между тем, что имеем в виду и тем, что говорим, между намерением и действием, сознанием и бессознательным. Мы --- отчужденные, потому что сказанное нами выражает скорее наше желание, чем то, что мы имели в виду. Наши действия показывают, кто мы, какими бы ни были наши намерения. Бессознательное не тайно, мы всегда демонстрируем его другим своей речью и действиями. И потому другие знают его лучше, чем мы сами: видят его проявление там, где мы видеть не можем. Сами же мы сталкиваемся со своей субъектностью как с чуждым в своей речи и поступках, как с тем, чего не хотели говорить и делать.

Бессознательное желание --- не просто неосознанное желание. Сознание не может принять его как свое, потому что оно нарушает его логику. Это и есть фундаментальный раскол между сознательными и бессознательными желаниями, и этот раскол не преодолеть никакими размышлениями. Даже те, кто совершенно осознает его, все равно поддаются бессознательному желанию, подрывающему сознательные. Субъект постоянно в противоречии с собой. В психоанализе раскол проявляется в негативной терапевтической реакции.

Безуспешность терапии Фрейд объясняет бессознательным. Пациент хочет поправиться, но бессознательное удовлетворяется расстройством и борется за его поддержание. Согласно Фрейду, ''удовлетворение бессознательного чувства вины, пожалуй, самая большая выгода субъекта от болезни, и проявляется в силах, стремящихся сохранить это состояние''\footnote{Sigmund Freud, “The Economic Problem of Masochism” (1924), trans. Joan Riviere, in \textit{The Standard Edition of the Complete Psychological Works of Sigmund Freud}, vol. 19, ed. James Strachey (London: Hogarth Press, 1961), 166.}. Неудачи психоаналитического подхода Фрейда рассказывают о бессознательном гораздо больше, чем его успехи. Пациент, близкий к преодолению своего расстройства, будто удваивает усилия остаться с ним, хотя сознательно хочет выздороветь. Негативная терапевтическая реакция свидетельствует о первичности и превосходстве бессознательного над сознательным в субъекте. Это превосходство и есть отчужденность.

Сознательно субъекты желают преследовать свои интересы и способствовать своему благу. Ставят цели, достижение которых сделает счастливыми. Я хочу жить в комфортной квартире, работать на интересной работе, смотреть увлекательные телепередачи и есть изысканные десерты и уверен, что буду тогда счастлив. Совершенно точно я не хочу быть несчастным, даже если делаю что-то отвратительное --- ем сырые устрицы или занимаюсь борьбой в грязи.

Блез Паскаль был великим теоретиком структуры сознательных желаний. Он отлично понимал важность желания счастья в сознательной жизни субъекта: ''Все без исключения люди стремятся быть счастливыми. Какие бы средства они не использовали --- все стремятся к этой цели. Это мотив каждого поступка каждого человека, даже того, кто решает повеситься''\footnote{Blaise Pascal, \textit{Pensées}, ed. and trans. Roger Ariew (Indianapolis: Hackett, 2005 [1669]), 181.}. Говоря о самоубийце, Паскаль стремится подчеркнуть совершенную повсеместность стремления к счастью: его поиск неизбежен, мы всегда преследуем свои интересы, это наше фундаментальное свойство.

Если ограничиваться сознательным, то анализ Паскаля вполне убедителен: будучи сознательным существом, субъект ищет счастья. Все, что делает сознательное существо --- корыстно, даже самоуничтожение, потому что структура сознания организована вокруг воли субъекта, а не его желания. Бессознательное желание остается слепым пятном в теории Паскаля, и потому он как мыслитель не котируется в наше время. Субъект сознательно желает быть счастливым, но не может избежать бессознательного желания поддерживать в себе отчужденность, что требует продолжающихся актов саморазрушения. Если бы мы были теми, кто мы есть, Паскаль был бы прав. Но разлад в субъектности приводит к действиям, уводящим от счастья. Субъект отчужден от себя и желает сохранить это любой ценой\footnote{Субъект готов платить любую цену, лишь бы сохранить отчужденность, ведь чем больше цена --- тем надежней отчуждение. Ущерб символической идентичности идет на пользу бессознательному желанию, оно процветает в саморазрушающих действиях. Бессознательное питается неудачами сознательного.}.

Уникальность субъекта не в сознательном. Можно представить --- и в научно-фантастических фильмах это показано --- клонирование сознания. Но соотнесенность субъекта с сознательным повторить не получится. Эта соотнесенность и есть бессознательное, основа единичности субъекта. Оно отдаляет нас от сознательного, и благодаря дистанции мы, субъекты, отдельны. Бессознательное не позволяет нам быть теми, кто мы есть, и потому мы выделяемся из природы и культуры. Бессознательное --- добавочное к символическому для получения субъектности\footnote{Сознание отвергает то, что кажется ему отталкивающим, что не соответствует его ценностям, бессознательное же не отвергает ничего. Об этом пишет Фрейд в своем кратком эссе ''Отрицание'': ''В бессознательном нет места ''нет'''' (Sigmund Freud,
''Negation'' (1925), trans. Joan Riviere, in \textit{The Standard Edition of the Complete Psychological Works of Sigmund Freud}, vol. 19, ed. James Strachey (London: Hogarth Press, 1961), 239). Всеприятие заставляет бессознательное принимать и то, что символическая структура скрывает.}. Отчужденность коренится в бессознательном и --- в отличие от сознательных желаний --- несводима ко всему, что на нас влияет.

Субъекты сознательно хотят быть счастливыми, но подсознательно находят удовлетворение в подрыве этой своей воли. Удовлетворение заключается в предотвращении достижения счастья, счастье подобно миражу, что обнажило бы его достижение. Счастье не может быть сейчас. Достигнутость значит, что достигнуто не то, что желалось. Поэтому счастье существует только как обещание счастья. Бессознательное мешает его достижению, тем самым сохраняя его в единственно возможном виде --- в виде мечты.

Изначального удовлетворяющего объекта нет, поэтому бессознательное удовлетворяется созданием и поддержанием отчужденности. Субъект находит удовлетворяющий объект только через его утрату --- утрата позволяет желать. Желание следует рассматривать не как естественное явление, а как структуру, возникающую из взаимосвязи телесных потребностей и социальных требований. Бессознательное желание существует, подрывая сознательное стремление к счастью, постоянно создавая препятствия на пути исполнения сознательных желаний. Бессознательное удовлетворяется в противодействии сознательному.

К каждому высказыванию субъекта бессознательное добавляет значение, которое он не имел в виду, поэтому субъект, будучи пронизан бессознательным, неспособен на лишь констатацию факта. Его бессознательное желание определяет каждое высказывание, и даже самый, казалось бы, невинный комментарий несет его клеймо, которое и заставило субъекта говорить. Когда я говорю соседу, что зловещие тучи предвещают дождь, я не просто предсказываю погоду. Причина этого высказывания кроется за его абстрактным значением, подрывая его невинность. Возможно, я хотел испортить соседу настроение, зная, что он собирался на пикник, и завидуя ему. Или, может, я хочу укрепить связь между нами, показав нас как товарищей по несчастью. Но скорее всего этот банальный комментарий --- способ пережить травму от встречи с соседом. Говоря бессмысленные банальности, я показываю, что мне нечего противопоставить отчуждающей силе его субъектности. Неважно, какой из этих вариантов верен. Бессознательное желание определяет даже самые безобидные замечания, определяет, вообще, все высказывания субъекта. В речи желание неизбежно искажается.

Искажение проявляется и в интерпретации высказываний других. Когда друг говорит мне, что у меня хорошая квартира, я не принимаю это как лишь констатацию факта. Я задаюсь вопросом, почему он это сказал, и отвечаю согласно причине. Возможно, он собирается занять у меня денег, или считает, что обычно моя квартира выглядит грязной, или даже не любит опрятных и собирается разорвать всякие отношения со мной. Желание если и явлено, то во множестве возможностей, а принять утверждение за чистую монету нельзя. Даже не имеющее сознательного манипулятивного намерения, прямое заявление о чистоте моей квартиры несет багаж бессознательного желания. Чтобы говорить, субъект должен хотеть говорить --- и даже молчание подчинено желанию, ведь так же происходит на фоне речи. Если вместо комплимента мой друг промолчит, я все равно буду думать --- теперь над смыслом его молчания. Будучи отчужденным, субъект не может говорить, не извращая сказанное бессознательным. И даже первые, еще не учитывавшие бессознательное, попытки построения теории субъектности укореняли субъектность в отчужденности.

\section{Почему Кант не может быть собой?}

Происхождение термина \textit{субъект} выявляет его внутреннюю связь с отчужденностью. Говоря о субъекте, мыслители всегда говорят об отчужденном субъекте, даже если прямо не указывают это. И они никогда не говорят о субъекте вне его самораскола. \textit{Субъект} о всегда внутренне дистанцированном от себя, в отличие от \textit{я}, \textit{индивида}, \textit{личности} или \textit{человеческого существа} --- эти термины понятней и несут меньше теоретического жаргона, но не подразумевают отчужденность. Конечно, можно уточнить, что речь о раздвоенном я или о раздвоенной личности, но \textit{субъект} уже предполагает это. Не бывает неотчужденного субъекта, всякий субъект --- субъект отчуждения. Это очевидно в том, как термин был первый раз использован в сегодняшнем смысле философом.

Иммануил Кант не изобрел термин \textit{субъект} (\textit{Subjekt} по-немецки), но в \textit{Критике чистого разума} и последующих работах популяризировал его использование в новом смысле. Канта можно по праву считать первым философом субъектности. Начиная с него, \textit{субъект} значит уже не подчиненного --- а именно в этом смысле термин использовали предшественники Канта, например, Барух Спиноза и Готфрид Лейбниц\footnote{Фридрих Ницше будет считать это грубой ошибкой Канта. Согласно Ницше, субъект --- фикция, не реальная сущность, а философов путает грамматика, и они считают субъекта чем-то действительно существующим (в английском и немецком языках \textit{субъект} и \textit{подлежащее} обозначаются одним словом). В \textit{По ту сторону добра и зла} он задает риторический вопрос: ''Не следовало бы нам относиться к субъекту так же, как относимся к сказуемому и дополнению? Не пора ли философам быть выше веры в грамматику?'' (Friedrich Nietzsche, \textit{Beyond Good and Evil: Prelude to a Philosophy of the Future}, trans. Judith Norman (Cambridge: Cambridge University Press, 2002 [1885]), 35). Но Ницше не поясняет, как его концепция сверхчеловека избегает такой грамматической ловушки. Разве нельзя точно так же свести мысль Ницше к грамматике и так отвергнуть?}. После Канта \textit{субъект} значит сущность, негативно соотнесенную с самой собой, не являющуюся тем, чем является\footnote{Это касается не всех --- Луи Альтюссер, например, критикует концепцию субъекта, говоря, что она способствует иллюзии самоидентичности и самообладания, и основывая на этом свою критику идеологической инерпелляции --- процесса, порождающего, согласно Альтюссеру, субъектность.}.

Кант использует концепцию субъектности для отделения сущности, называемой субъектом, от остального мира. Поэтому для Канта субъектность --- всегда отчужденная субъектность, хотя он, конечно, не пишет это явно и не использует термин \textit{отчужденность}\footnote{Кант популяризировал термин \textit{субъект} в философском смысле, Гегель же сделал то же с термином \textit{отчуждение} (Entfremdung). Идею отчуждения обычно связывают с Марксом, но у Гегеля, в отличие от Маркса, она занимает центральное место.}. Согласно Канту, субъект никогда не чувствует себя в мире как дома, потому что конституирует свой мир посредством акта трансцендентального воображения --- без этого не было бы никакого мира, только скопище бессвязных входных данных. Субъект не существует в своем мире, но конституирует его. Субъектность --- условие возможности опыта. Кант, впрочем, не отрицает внешний мир, и старательно --- особенно во втором издании \textit{Критики чистого разума} --- сепарируется от Джорджа Беркли\footnote{Кант добавляет во второе издание \textit{Критики чистого разума} раздел \textit{Опровержение идеализма} исключительно чтобы читатели не спутали его позицию с позицией Беркли.}, но утверждает, что мир получается связным только благодаря трансцендентальному акту субъекта.

Термин \textit{субъект} возникает в начале \textit{Критики чистого разума}, где Кант описывает, как чувственность структурирует наши интуиции. Чувственность --- одна из трех способностей наряду с рассудком и разумом, посредством которых субъект придает воспринимаемому им миру связность. Использование термина \textit{субъект} указывает, что описываемая сущность должна действовать сама на себя, чтобы получить осмысленный опыт. Для Канта субъект --- отделенное от себя и потому способное на себя действовать, существо.

Будучи источником пространства и времени, субъект (в понимании Канта) сам не может быть в мире ни пространственно, ни темпорально. Мир сам по себе не пространственен и не темпорален, эти формы --- вклад субъекта. Согласно Канту, ''Время --- ... лишь субъективное условие нашего созерцания (которое всегда чувственно, т.е. возможно только благодаря воздействию объектов на нас), и само по себе, вне субъекта --- ничто. И все же оно необходимо объективно в отношении всякого явления --- всякой вещи, могущей предстать перед нами в опыте''\footnote{Immanuel Kant, \textit{Critique of Pure Reason}, trans. Paul Guyer and Allen W. Wood (Cambridge: Cambridge University Press, 1998 [1781]), 164.}. Время как форма возникает от субъекта, но ее нельзя отбросить: время структурирует опыт, и поэтому Кант настаивает на его объективности. Но формы пространства и времени также показывают вневременность и внепространственность самого субъекта. Субъектность порождает темпоральность только потому что сама ей не подвержена.

Кантианский субъект есть свое же разделение. Он осознает себя потому что от себя дистанцирован, не тождественен тому, что он есть. Самораскол субъекта особенно заметен в отношении субъектности к свободе и детерминированности. Кант утверждает --- и пытается доказать --- что субъект не свободен \textit{или} детерминирован --- он свободен \textit{и} детерминирован. И дело не в том, что одни его действия свободны, а другие детерминированы --- каждое действие и свободно, и детерминировано. С теоретической точки зрения субъект детерминирован. Но с практической --- как существо в действии --- свободен. И, согласно Канту, оба подхода верны\footnote{Кристина Корсгаард в своем эпохальном исследовании Канта \textit{На пути к царству целей} пишет: ''Возможны два применения разума: теоретическое и практическое. Мы рассматриваем себя как феномены, когда, применяя его теоретически, описываем и объясняем свое поведение; и как ноумены --- когда практически, решая, что делать. Эти две точки зрения не пересекаются, поскольку их цели --- объяснение и решение --- исключают друг друга'' (Christine M. Korsgaard, \textit{Creating the Kingdom of Ends} (Cambridge: Cambridge University Press, 1996), 204). Такая концепция субъектности --- разделенной на теоретическое и практическое --- концепция субъектности как внутренне отчужденной.}.

Только благодаря фундаментальной отчужденности субъект может быть сразу и свободен и детерминирован. Кантианский субъект столь отделен от самого себя, что его можно одинаково успешно мыслить как теоретически, так и практически --- но не теоретически и практически сразу: теория и практика в философии Канта строго несовместимы. Позднее немецкие идеалисты попытаются объединить эти два полюса, что, согласно Канту, рискует исключить фундаментальную для субъекта отчужденность. Фихте, например, будет искать теоретическую перспективу внутри практической, но субъект у него получится уже не столь разделенным, хотя это и исключит для него некоторые проблемы, свойственные философии Канта.

Само разделение философии Канта на теоретическую --- представленную в \textit{Критике чистого разума} --- и практическую --- наиболее известное изложение которой дано в \textit{Критике практического разума} --- символизирует отчужденность субъекта. Эти два направления не просто сосуществуют в его философии, но и чужды и противоречат друг другу. Для теоретического Канта практический Кант --- выход за пределы мыслимого, но для практического Канта теоретический слеп к проблеме свободы субъекта. Нельзя предпочесть одно другому, отчужденность столь важна для Канта, что сама структура его философии проникается ей. По мнению Канта, отчужденность --- не то, что субъект может преодолеть. Потому что это определение субъектности.

\section{Сломанные животные}

Нельзя понять, как работает политика, не принимая во внимание отчужденность. Мы являемся политическими существами не из-за партийной принадлежности, а из-за предвзятого отношения к общественному порядку. Быть политизированным значит не быть нейтральным. Всякий субъект политизирован из-за своей искаженности означающим, своей фундаментальной отчужденности. Отчужденность конституирует политизированность, и нам не отказаться от нее, не убежать от нее в ситуации, кажущиеся деполитизированными. Отчужденность необходима, ни один субъект не может быть сведен к биологии, составляющей его материальную основу\footnote{Согласно Джорджио Агамбену, нацистские лагеря смерти были призваны лишить заключенных в них всякой политической силы, сделать их лишь биологическими существами. В \textit{Останках Освенцима} он пишет о лагере как о месте биовласти, обгладывающей субъекта до его биологических функций, оставляющей от него лишь животное. Он пишет: ''биовласть действует через производство не жизни или смерти, а изменчивого и практически нескончаемого выживания'' (Giorgio Agamben, \textit{Remnants of Auschwitz: The Witness and the Archive}, trans. Daniel Heller-Roazen (New York: Zone Books, 1999 [1998]), 155). Но субъект не может просто выжить --- как бы жестоко на него не действовали репрессивные структуры власти, он остается отчужденным --- а значит и политическим --- существом.}. И дело не в том, что субъектность никак не связана с человеческим животным --- искажение означающим, порождающее субъект, столь фундаментально меняет животное, что должно приниматься отправной точкой в анализе того, что мы делаем. Несмотря на статус биологических существ, мы --- благодаря отчужденности --- можем действовать неестественно. Неестественное склонности в нас превосходят физиологические императивы.

Представляя субъектов естественными и/или культурными сущностями, мы забываем про их отчужденность. Субъективность, понятая как продукт природных или культурных сил, лишена политического бытия --- несводимости к тому, что ее определяет. Натурализация или культурализация субъектности не лишает субъектов политизированности, но делает ее незаметной. Будучи натурализованы или культурализованы, политические решения рассматриваются уже как только выражение врожденного личного интереса, а не как действия свободного субъекта. Но отчужденный субъект --- не корыстное существо, не реагирует так, как ему выгодней, согласно его планам. Субъектности свойственны бессознательные попытки саботировать собственные интересы, находить и создавать себе трудности.

Субъектность не делает человека хозяином остального мира\footnote{Декарт надеялся, что научные исследования, в отличие от чисто спекулятивной философии, ''сделают нас как бы господами и хозяевами природы'' (René Descartes, \textit{Discourse on the Method of Rightly Conducting One’s Reason and Seeking the Truth in the Sciences}, trans. Robert Stoothoff, in The Philosophical Writings of Descartes, vol. 1 (Cambridge: Cambridge University Press, 1985 [1637]), 142–143). Многие комментаторы указывали на опасность такого отношения к природе, Декарт здесь изменяет собственному пониманию субъекта как отдельного от мира и от самого себя, пониманию, которое развивает на протяжении всего \textit{Рассуждения о методе}. Так, ранее в трактате он говорит о равенстве, основываясь на внутренней отчужденности субъекта от места его происхождения.}. Напротив, она отчуждает его от мира, делает его господство над миром невозможным. Из-за отчужденности отношение субъектности к миру всегда на расстоянии и всегда искажено\footnote{Смотреть сквозь символическое --- еще один способ избежать отчужденности. Его в \textit{Реальных Обманах} рассматривает Дженнифер Фридлендер: ''Как субъекты, мы должны бороться с фундаментальной ложью Символического'' (Jennifer Friedlander, \textit{Real Deceptions: The Contemporary Reinvention of Realism} (Oxford: Oxford University Press, 2017), 123). Попытка выйти из блуждания в символической структуре --- это отказ принимать отчужденность. Необходимость фундаментального обмана значит необходимость отчужденности.}.

Субъект в своей основе --- естественное существо, но совершенно извращенное субъектностью. Оно уже неестественно или даже де-естественно. Образ укорененной в природе субъектности на протяжении веков вводил в заблуждение мыслителей. Аристотель, пытаясь определить политическое бытие человечества, обращается к миру природы и нашему к нему отношению. В начале ''Политики'' он дает самое известное определение роли политики в человеческом существовании: ''Человек по своей природе --- политическое животное''\footnote{Aristotle, Politics, trans. B. Jowett, in \textit{The Complete Works of Aristotle}, vol. 2, ed. Jonathan Barnes (Princeton: Princeton University Press, 1984), 1278b.}. Аристотель ошибочно считает источником политизированности наше положение \textit{в} природе, а не отчужденность от нее. Он натурализует политическое бытие говорящего субъекта, что притупляет противоречия, возникающие в обществе и внутри самой субъектности, противоречия, которые сохранялись в человеческой исключительности Платоном\footnote{Сторонники освобождения животных питают слабость к Аристотелю, ведь это принятие преемственности делает его их потенциальным союзником, в отличие от Платона, постулирующего фундаментальный разрыв между животным и человеком. И хотя Аристотель утверждает четкую иерархию среди животных и людей, он все же наделяет животных душой.}.

Политика, согласно Аристотелю --- не борьба противоположностей, а объединение ради общего блага. Он не видит антагонизма между господином и рабом или между мужчиной и женщиной, его политическая философия допускает беспрепятственное угнетение. Понимание человека как нечта естественного, как продолжения природы, прямо связано с присвоением рабам и женщинам статуса второстепенных\footnote{Платон, в отличие от Аристотеля, видит идеальное общество обществом без рабов и ущемления женщин. Он отвергает демократию, считая философов правящим классом в таком обществе, обществе равенства, которое Аристотель не может даже представить.}.

Преемственность человека с миром природы позволяет Аристотелю натурализовать --- и тем самым оправдать --- неравенство внутри политики. Ведь если возможно --- и естественно дозволено --- неравенство в природе, то возможно и дозволено неравенство в человеческом мире. Юридический закон, согласно Аристотелю, действует как естественный закон, описывая уже существующие отношения. Не требует от субъектов действовать против самих себя.

Аристотель так легко принимает неравенство, потому что не осознает отчужденность субъекта от мира природы. Попытка укоренить субъектность в природном всегда ведет к оправданию социального неравенства. Если мы просто природные существа, то равенство невозможно: Флоренс бегает быстрее, чем Инге, Боб умнее Карла, а Карим выше Кальвина. В естественных отношениях доминирует явное неравенство, но, к счастью, мы никогда не остаемся в лишь естественных отношениях друг с другом, чего Аристотель и не учитывает.

Понятие равенства требует говорить о денатурализованном, отчужденном субъекте. Только если политизированность субъекта обеспечена его разрывом с миром природы --- только тогда возможно равенство субъектов. Политическую борьбу не следует рассматривать как лишь поле различных, конкурирующих друг с другом, за продвижение собственных интересов, животных. Именно благодаря денатурализованности субъект способен на антагонистическую борьбу во имя свободы и равенства. То, как мы поймем субъектность, определяет, как мы станем конституировать политическое поле и политические возможности, так что вопрос этот --- не чисто теоретический.

Разрыв между субъектом и природой оставляет между ними очевидную связь. Субъект обладает природным происхождением и животным телом. Он все еще смертен, уязвим к болезням и телесным повреждениям. Животное, став субъектом, остается животным. Но животности в чистом виде для него уже нет, все ее аспекты извращены субъектностью.

Отчужденность дистанцирует субъекта от животной сущности, но она же конституирует его как сущность иного типа, которой бы иначе не было. Субъектность --- не результат эволюции, а разрыв, внесенный означающим в сферу животности. Возможно, означающее возникло в ходе эволюции, но из-за него были нарушено императивы естественного отбора и даже выживания. Разрыв с означающим --- не триумф культуры над природой, потому что субъектность зависла между ними\footnote{Психоанализ имеет дело не с человеческим животным или культурной идентичностью, а с заключенным в них субъектом. Фрейд не вполне понимает природу собственной теории, когда надеется на ее подтверждение биологией. Никакое научное открытие не способно подтвердить достоверность психоанализа, надежда на обратное значит --- пусть даже невольный --- отказ проекта.}.

\section{Извращенные удовольствия}

Отчужденность является в неумеренности для компенсации недостатка самотождественности. Иначе говоря, несостоятельность идентичности приводит к свеврхидентификации. Перед лицом отчужденности субъект становится фигурой избытка, что открывает ему новые пути к удовлетворению. Он не может найти по-настоящему сытую пищу и ест слишком много. Не может заниматься исключительно репродуктивным сексом и изобретает множество извращенных сексуальных актов. Еда и секс для субъекта отделяются от животных функций. Они, конечно, продолжают удовлетворять биологические потребности, но биология теперь вторична и не играет роли в удовлетворении.

Биологическая потребность может даже мешать удовлетворению отчужденного желания субъекта. Меня не привлекает то, что лучше всего удовлетворяло бы мои биологические потребности. Зато шоколадные пончики вызывают такое желание, что устоять можно только приложив немалые усилия. Даже если этот выбор эволюционно обусловлен, получаемое удовлетворение свидетельствует о разрыве между желанием и биологией: вредное воздействие пищи на организм делает ее вкусной. Наши желания идут против биологических потребностей\footnote{Некоторым удается питаться здоровой пищей, избегать сладкого и жирного, но делают они это не естественным путем. Здоровое питание --- всегда ответ на возможность нездорового питания. Выбирая шпинат вместо жареного печенья Oreo, моя свекровь находит удовлетворение в \textit{отказе} от вредного, а не в выборе полезного. Для отчужденного субъекта нет неотчужденного отношения к еде.}.

В сфере сексуальности раскол между желанием и биологией еще очевидней. Женщины стали рожать в более позднем возрасте, и теперь пары рассчитывают, когда будет овуляция, и занимаются сексом в это время, чтобы максимизировать шансы на зачатие ребенка. Будь удовлетворение от секса удовлетворением от самовоспроизводства, такие точно рассчитанные акты были бы самыми приятными. Но все наоборот --- секс ради зачатия становится выполнением обременительного долга, что лишает удовольствия. Это все равно что питаться только сырой капустой.

Если отслаивание инстинкта от чисто биологической функции --- извращение, то все субъекты --- извращенцы в еде и сексе\footnote{Фрейд так и определяет сексуальные извращения, но тут же настаивает, что из-за распространенности они не являются чем-то ненормальным: ''Перверсии --- сексуальные действия, которые либо (а) анатомически выходят за пределы предназначенных для сексуального акта областей тела, либо (б) фиксируются на непосредственных отношениях с сексуальным объектом, откладывая конечную сексуальную цель'' Sigmund Freud, \textit{Three Essays on the Three of Sexuality}, in \textit{The Standard Edition of the Complete Psychological Works of Sigmund Freud}, vol. 7, trans. and ed. James Strachey (London: Hogarth Press, 1953 [1905]), 149–150.}. Для говорящего субъекта удовлетворение от еды и секса возможно только благодаря их дистанцированности от соответствующих биологических потребностей. Переедая, он ест не то, что способствовало бы его выживанию --- напротив, потребляет вредные продукты: печенье, пирожные и мороженое. Отдаваясь сексуальным страстям, он изобретает и практикуют предельно отдаляющее от воспроизводства: мастурбацию, минет, куннилингус, анальный секс и пр.. Никто не объедается брокколи и никто не доводит себя до инфаркта репродуктивным сексом.

Переедание и сексуальную зависимость обычно пытаются объяснить биологически. Ожирение рассматривают не как следствие отчужденности субъекта, а как развитие эволюционного несоответствия: мы развили привычки, но ситуаций, в которых эти привычки помогали, больше не возникает, так что теперь они только вредят. Как пишет биолог-эволюционист Дэниел Либерман: ''Миллионы лет эволюции успешней были те, кто искал и много ел богатую простыми углеводами пищу''\footnote{31 Daniel Lieberman, \textint{The Story of the Human Body: Evolution, Heath, and Disease} (New York: Pantheon, 2013), 16.}. Эволюция подготовила человека-животного должным образом, но культурные изменения оказались быстрее естественного отбора и привели к широкому распространению ожирения в современном обществе. Ученые и большая часть общественности склонны принимать такую интерпретацию и рассматривать людей как животных, а не как [отчужденных] субъектов.

Излишества в сексуальной сфере принято объяснять аналогично таковым в кулинарной. Самцы изменяют, потому что их \textit{гены хотят} распространяться как можно шире. Самки выбирают моногамию, потому что их \textit{гены видят} в этом лучшую стратегию --- помогающую выжить их потомству. В книге ''Как работает разум'' Стивен Пинкер объясняет мужскую полигамию адаптацией: ''Чем с большим количеством женщин мужчина займется сексом, тем больше потомства он оставит --- много не мало. Из-за этого мужчины склонны стремиться к случайным половым связям, власти и богатству''\footnote{Steven Pinker, \textit{How the Mind Works} (New York: Norton, 1997), 473. Пинкер утверждает, что такая адаптация свойственна только мужчинам, а женщины едва ли [генетически] заинтересованы иметь много партнеров. Эволюционный психолог Дэвид Бусс считает иначе: ''Экономика рынка спаривания такова, что женщины в целом заинтересованы не только в надежном партнере, но и в лучших генах для своего потомства, даже если эти гены может предоставить только кратковременный партнер --- в идеале они стремятся при этом сохранить отношения с надежным. Генетически успешный мужчина часто готов на кратковременную, не обременяющую связь с не столь успешной женщиной'' --- David Buss, \textit{The Evolution of Desire: Strategies of Human Mating}, rev. ed. (New York: Basic Books, 2003 [1994]), 235–236. Но между этими аргументами можно выбирать только на политических основаниях. Современный кантианец может даже говорить здесь об антиномиях натуралистического разума.}. Естественный отбор настолько всеобъемлющая концепция, что может объяснить даже самые экстремальные половые акты --- участие в них свидетельствует о приспособленности. Нет извращения столь извращенного, чтобы его нельзя было объяснить адаптацией.

Некоторые мыслители-натуралисты приложили удивительно много усилий к объяснению изнасилований адаптацией. Рэнди Торнхилл и Крейг Палмер подвергались критике со стороны коллег за изучение причин изнасилований как явления. Согласно результатам их исследования ''Естественная история изнасилования: биологические основы сексуального принуждения'' допускается две такие причины. Торнхилл считает, что изнасилование --- результат адаптации, Палмер же --- что это побочный продукт адаптационного процесса. Торнхилл а Палмер старательно подчеркивают, что через свое исследование пытаются найти более эффективные механизмы предотвращения изнасилований, но будучи понято как следствие биологического наследования, изнасилование будто оправдывается\footnote{See — or not — Randy Thornhill and Craig T. Palmer, \textit{A Natural History of Rape: Biological Bases of Sexual Coercion} (Cambridge: MIT Press, 2000).}.

Торнхилл и Палмер довели натуралистическое объяснение до конца, обнажив его ограничения: если все эксцессы субъекта --- лишь развитие естественных склонностей, то это же верно и для изнасилования. Поиск натуралистического решения проблемы изнасилований политически нецелесообразен, но напрямую следует из рассмотрения людей как животных, а не как субъектов. Все действия [неотчужденных] человеко-животных --- даже самые чрезмерные --- естественны и, значит, укладываются в адаптационные рамки. Адаптационные объяснения очень популярны, так что если бы не Торнхилл и Палмер, это сделал бы кто-то другой. Натуралистический подход не может признать, что жестокость изнасилования --- продукт извращенной субъектности, результат того, что субъект удовлетворяется унижением другого. Согласно Торнхиллу и Палмеру, ни одно извращение не выходит за пределы естественного. Но вопрос изнасилования --- не самый проблемный для натуралистов.

Согласно натуралистическому подходу, самый отъявленный сексуальный извращенец --- не калечащий свой пенис мазохист и не пытающий свою жертву садист, а уже лишившийся репродуктивных способностей, но все еще живой, человек. Как отмечает антрополог Питер Грей в своем обзоре эволюции человеческой сексуальности: ''Жизнь людей за пределами репродуктивных способностей --- эволюционная загадка''\footnote{Peter B. Gray, “Evolution and Human Sexuality,” \textit{Yearbook of Physical Anthropology} 152 (2013): 20.}. Согласно натуралистической теории, продолжительность нашей жизни как животных избыточна, потому что превышает оптимальную для воспроизводства своих генов. Более того, люди продолжают заниматься сексом еще долгое время после потери репродуктивного потенциала --- они не посвящают себя целиком социальной функции бабушек и дедушек.

Распространенность половых инфекций в пенсионных сообществах идет вразрез с представлением о субъекте как естественном существе, ведь ни особь, ни вид не получают пользы от такой сексуальной активности --- особи женского пола уже не способны к размножению. Да, большинство мужчин в таком возрасте --- в отличие от женщин --- все еще способны к размножению, но в целом описанное явление противоречит духу натурализма, показывая чрезмерность и бесцельность для естественного отбора. Люди продолжают заниматься сексом в старости именно из-за отчужденности, до крайностей их доводят бессознательные желания, не соответствующие потребностям.

Конечно, все эти вызванные субъектностью извращения можно представить как развитие или ответвление естественного отбора человека-животного. И даже абсурдность секса в старости для человека-животного можно отрицать, приводя аргументы, которые некоторые примут. Но эти аргументы неадекватны тому, что пытаются описать, потому что предполагают, что люди действуют из личных интересов, а не в ответ на свою отчужденность как субъектов.

И такие необъяснимые самовоспроизведением извращения субъектности --- не исключения, обнаруженные у нескольких изолированных субъектов с какой-нибудь ошибкой при кодировании генов. Они проявляются у каждого субъекта, и более того --- дают им удовлетворение. Избыточности субъектности избыточны не как аппендикс --- они дают повод существовать дальше. Без извращенных эксцессов самовоспроизведение не заинтересовало бы ни одного субъекта, у них просто не было бы желания это делать. Они продолжают существовать, удовлетворяясь извращенным нарушением императива самовоспроизводства. Их самовоспроизведение держится на отказе следовать генетическому импульсу самовоспроизводиться.

Если бы субъекты могли просто преследовать свои интересы, то принадлежали бы миру, не были бы отчужденными. Субъект не принадлежит миру, он постоянно отклоняется от преследования блага --- своего или мира. Субъект не принадлежит миру из-за своей извращенности, и эта извращенность --- несмотря на все производимые ей ужасы --- ключевое в субъекте, именно она делает его тем, кто он есть.

Натуралистический подход к извращениям субъектности не только не адекватен, но и опасен. Согласно ему, эти чрезмерности никогда не чрезмерны --- это нарушения, естественные, как и мир, в пределах которого остаются. Не придавая этим излишествам должной радикальности, такой подход оставляет субъекта частью мира, хотя именно денатурализованная --- самоотчужденная --- субъектность порождает их. Будучи отчужденным, субъект более не ограничен принадлежностью миру.

Отчужденность ответственна за ужасы и разрушения, но ответом на него должно быть принятие. Без отчужденности субъекту не хватило бы уникальности, придающей ценность его существованию. Понимая извращенность субъекта как естественное отклонение, мы нормализуем ее. Но отчужденность является одновременно и проблемой, и решением. Мы должны видеть в отчужденности субъекта основу его этического существования.\footnote{Понимание отчужденности фундаментальной для себя как субъекта --- это то, что Гегель в \textit{Феноменологии Духа} называет абсолютным знанием. Достигнув его, можно чувствовать себя как дома в том, что чуждо. Согласно Гегелю, дух, который достиг абсолютного знания, ''в \textit{своей} инаковости как таковой тождественен с собой'' (G. W. F. Hegel, \textit{Phenomenology of Spirit}, trans. Terry Pinkard (Cambridge: Cambridge University Press, 2018 [1807]), 454). Быть как дома в своей инаковости --- значит обитать в конститутивной субъектности отчужденности, принять неуместность как свое единственное возможное отношение к месту. Привыкнуть к непривычности. Только так можно сохранить проблему субъектности и не стремиться к ее решению --- такое стремление только усугубило бы ситуацию.} Отчужденность ужасает если пытаться ее преодолеть, но общество модерна дает субъекту новые возможности поддерживать дистанцию от самого себя.

\chapter{Неприкаянная современность}

\section{К бесконечности и дальше}

Обозначение порождает отчужденность. Она не связано с историческим явлением или эпохой, как если бы люди ходили по улицам в счастливой неотчужденности, когда вдруг появились капиталисты и навязали им отчуждающую систему. Отчужденность неизбежна для говорящего, потому что присуще обозначению как таковому, а не какому-либо социальному порядку. Но несмотря на повсеместность, последствия отчужденности стали очевидны только с научными революциями модерна, и потому отчужденность часто ассоциируют с эпохой модерна, когда неприкаянность субъекта выходит на первый план. Но такое развитие событий должно встречать с радостью.

Гелиоцентризм, гравитация и относительность --- все эти открытия сделали вклад в демонстрацию отчужденности субъекта. Субъекту не отведено конкретное место.\footnote{В своей знаменитой истории современной науки Э.А.Бертт отмечает, что система Галилия вытесняет человека с прежней позиции. Согласно Берту, в этой системе ''мир природы изображен обширной, автономной математической машиной из движущейся в пространстве и времени материи, человек же с его целями и чувствами был исторгнут из нее как ничего не значащий зритель, как полуреальный побочный эффект великой математической драмы вовне'' (E. A. Burtt, \textit{The Metaphysical Foundations of Modern Science} (Mineola: Dover, 2003 [1924]), 104). Современная наука, стремясь понять вселенную, оставляет субъектность в стороне.} Когда нет ни центра Вселенной, ни целенаправленной, связывающей всё силы, то субъект не может понять свое положение в порядке вещей. Он обнаруживает свою неприкаянность. В традиционном мире субъекты могли легко убедить себя, что находятся на своем месте. Но оказалось, что составляющие Вселенную процессы существуют вне зависимости от субъектов. Субъект ни центр, ни даже периферия --- он избыток. Субъект оказывается неуместен и отчужден, как, впрочем, и все остальное.

Субъект должен посмотреть в лицо своей неприкаянности. Нигде он не будет чувствовать себя как дома. Это и разрушительный удар, и освобождающее откровение. Субъектность всегда была онтологически неуместной, но только теперь это стало заметно. Субъект всегда был неприкаян, но только в эпоху модерна его дикость выходит на передний план благодаря произошедшим открытиям.

Понимание субъектности с эпохой модерна подверглось хоть и не онтологической, но экзистенциальной трансформации. Опыт субъектности начал меняться с пришествием гелиоцентризма, эта революция вызвала оцепенение у религиозной власти. Бросив вызов геоцентризму в \textit{De Revolutionibus Orbium Coelestium} в 1543г., Коперник --- даже если не стремился к этому --- опрокинул традиционную позицию субъекта в мире.\footnote{Изъятие Земли из центра мира --- решающее событие современности. Стивен Вайнберг совершенно верно указывает: ''Какой бы ни была научная революция, она началась именно с Коперника'' (Steven Weinberg, \textit{To Explain the World: The Discovery of Modern Science} (New York: HarperCollins, 2015), 147). И хотя Коперник не пострадал за свои убеждения как затем Галилео Галилей и Джордано Бруно, его прорыв был основополагающим для эпохи модерна.}

Но смещение Земли и, значит, человечества, с центра творения само по себе еще не источник отчужденности. Чувствовать себя как дома можно и в центре, и на периферии.\footnote{Согласно Аристотелю, окружающие Землю небесные сферы и стоящий за ними перводвигатель являются сферами совершенства, и об этом свидетельствует их круговое движение. Подлунный же статус Земли указывает на ее несовершенство. Для христианских мыслителей-аристотеликов, позиция Земли в мире соответствует низкому статусу человечества в сравнении с остальным творением.} Перифирийность еще не значит неприкаянность. Перифирия все еще место и связана с центром, вокруг которого и находится, так что удар гелиоцентрической гипотезы не так уж силен. Отчужденность субъекта вскроется позже, когда другими учеными будут указаны последствия гелиоцентрической системы за пределами выводов Коперника.

Джордано Бруно был первым кто отринул веру в конечный мир и постулировал открытую вселенную. В \textit{Причина, Принцип и Единство} он выдвинул гипотезу о бесконечном пространстве, противоречащую как геоцентризму Аристотеля, так и гелиоцентризму Коперника. Бруно утверждал: ''Вселенная... бесконечна и безгранична и, следовательно, неподвижна''.\footnote{Giordano Bruno, \textit{Cause, Principle, and Unity}, trans. and ed. Robert De Lucca (Cambridge: Cambridge University Press, 2004 [1584]), 87.}. Утверждая это, он отказывается от своей принадлежности к конкретному месту, потому что отсутствие предела космоса делает невозможным конкретную локализацию себя относительно него. В бесконечной Вселенной отчужденность субъекта становится очевидней. Вклад Бруно в выявление отчужденности оказывается даже больше, чем таковой Коперника, и другие мыслители вскоре станут делать еще больше выводов, подчеркивающих неприкаянность субъекта.\footnote{Артур Лавджой говорит о Бруно как о сыгравшем более значимую роль в осознании человечеством отчужденности субъекта, нежели Коперник: ''Джордано Бруно --- основной представитель доктрины децентрализованной, бесконечно большой и населенной Вселенной --- он не только проповедовал ее по всей Европе с пылом евангелиста, но и первым подробно изложил ее основания, добившись тем признания широкой общественности'' (Arthur O. Lovejoy, \textit{The Great Chain of Being: A Study of the History of an Idea} (Cambridge: Harvard University Press, 1936)). Приверженность бесконечному нарушает привычный способ думать о человечестве как об имеющем определенное место. В бесконечности идея центра и расстояния до него теряет смысл.}

В современной Вселенной нет конкретных мест. Все вещи локализованы только по отношению друг к другу, без привязки к фону. Это повлияло на анализ созвездий, проводимый Иоганном Кеплером. В \textit{Из закрытого мира --- к бесконечной Вселенной} Александр Койре показывает, что теоретизация на основе беконечных пространств исключает всякую выделенную позицию, и даже звезды больше не имеют места: ''для конечного мира разумно выбирать конкретный паттерн из звезд, но в бесконечном мире принцип достаточного основания не позволяет это геометрически мыслящему Богу Кеплера''.\footnote{Alexandre Koyré, \textit{From the Closed World to the Infinite Universe} (Baltimore: Johns Hopkins University Press, 1957), 78.} Бесконечная вселенная не способна вместить конечные места, в ней нет однозначной точки отсчета. Ни у звезд, ни у Земли, ни у населяющих ее существ --- ни у чего нет абсолютного места, места во Вселенной.

И хотя кажется комфортным и даже приятным иметь свое место, свой дом, такое место указывало бы на несвободу. Когда у всего есть Место, нет никакой свободы.\footnote{Даже Спартак поднимает восстание не во имя универсальной свободы --- столь ценным было Место в римском мире. Но уже на современном Гаити универсальная свобода смогла стать локомотивом восстания.} Свобода --- это способность оставить назначенное внешним Место.\footnote{В древнем Индуизме была концепция свободы, но то была не свобода действия --- а свобода от действия. Согласно Бхагаватгите ''тот, кто свободен от эгоистичных привязанностей, овладевший собой и своими страстями, тот достигает высшего совершенства свободы от действий'' (\textit{Bhagavad Gita}, trans. Eknath Easwaran (Tomales: Nilgiri Press, 1985), 18:49). Тысячилетия спустя после написания Бхагавадгиты Махатма Ганди положит свободу от действия в основание политической свободы. Ганди теоретизировал свое восстание через Бхагавадгиту, но важнее, что оно бросило вызов ограничениям социального Места.} Те, кто получает выгоду от социальной иерархии, заверяют, что все должно оставаться как есть.\footnote{В ''Теории субъекта'' Ален Бадью определит правящий класс через его приверженность стабильности Места: ''Правящий класс --- хранитель этого Места'' (Alain Badiou, \textit{Theory of the Subject}, trans. Bruno Bosteels (New York: Continuum, 2009 [1982]), 184). Эпоха модерна станет эпохой революционных взрывов именно потому что бросит вызов Месту.} Нарушая имеющееся распределение мест, эпоха модерна открывает окно для свободы.

Отсутствие стабильных мест становится основой социальных потрясений. Все великие эмансипационные движения современности --- от восстания рабов на Гаити до суфражистских движений по всему миру --- берут неприкаянность своей отправной точкой. Лидеры этих восстаний не смотрят на звезды в мыслях о бесконечности Вселенной, но схватывают отсутствие Места через социальный контекст. Они восстают против засилья Места.

Франц Фанон --- знаменитый теоретик отказа от Места, мыслитель, наиболее точно выразивший угнетение им. Он признает высвобождаемое неприкаянностью фундаментальное равенство. Выступая против расистского общества \textit{Черная Кожа, Белые Маски}, он призывает к разрыву с прошлым: ''Я не имею права пускать корни. Я не имею права допускать ни малейшей частички бытия в свое существование. Я не имею права увязать в решениях прошлого''.\footnote{Frantz Fanon, Black Skin, White Masks, trans. Richard Philcox (New York: Grove Press, 2008 [1952]), 204–205.} Концепция политики Фанона требует отзака от любых социальных или историческсих удобств. Это основа революционных действий.

Чтобы принять свою неприкаянность, не нужно эмигрировать с родной земли или даже покидать родной город. Неприкаянность не значит кочевой образ жизни. Кант был таким, хотя никогда не уезжал далеко от своего родного Кенигсберга. Дело в отказе от концепции Места, принадлежности к чему-то стабильному. Это необходимая основа эмансипации --- признание неприкаянности вместо стремления к воображаемому дому.\footnote{В книге Дэвида Гребера и Дэвида Венгроу ''Рассвет всего'' превосходно показано, что эпоха модерна была не просто европейским проектом. Великие идеи современности --- свобода, равенство и солидарность --- возникли в диалоге с другими, зачастую более просвещенными в этих вопросах, народами. И хотя Гребер и Венгроу не связывают эти ценности с отчужденностью и вообще не признают за ним никаких достоинств, их книга подобна глотку свежего воздуха среди эволюционных историй человечества. См. David Graeber and David Wengrow, \textit{The Dawn of Everything}: A New History of Humanity (New York: Farrar, Straus and Giroux, 2021)}

\section{Расставание с собой}

Эмансипация стала очевидно возможна только в рамках современной науки. Обнаруженная ей неприкаянность намекает на такое смещение и внутри самих людей. Это и есть отчужденная субъектность. Конечно, внутреннее смещение было и до современной науки, но стало заметно только теперь. И это радостное событие, не повод для сожалений. Отчужденность отделяет говорящего от самого себя и позволяет ему действовать вопреки внешним факторам, которые иначе определяли бы его. Но отчужденность по-настоящему освобождает только когда мы осознаем ее.

Исторически сложилось так, что общества предлагали говорящим субъектам образы, соответствующие традициям этих обществ.\footnote{Марксу удастся очень точно выразить должное отношение к традиции в \textit{Восемнадцатое брюмера Луи Бонапарта}. В начале этой работы он заявляет: ''Традиции ушедших поколений подобно кошмару тяготеют над умами живущих'' (Karl Marx, \textit{The Eighteenth Brumaire of Louis Bonaparte, in Political Writings, Volume 2: Surveys from Exile} (London: Verso, 2010 [1852]), 146). Политическое видение Маркса справедливо предполагает необходимый отход от традиций. Он признает, что традиция --- всегда несвобода, какой бы комфортной она нам не казалась.} Такая, характерная для домодернового мира, принадлежность наверняка давала людям некоторое спокойствие о своем статусе в нем, но лишала их уникальности, свойственной говорящему субъекту. Эта уникальность неотделима от статуса субъекта как отчужденного и от своего социального окружения, и от самого себя. Только дистанцируясь от себя, субъект в наибольшей степени может быть собой. Он становится целым только отказываясь от соблазна принадлежать. 

Разрушение иллюзии принадлежности --- великое, хоть и часто удручающе неадекватное --- достижение современности. Оно позволило субъекту ощутить определяющую его отчужденность, увидить расстояние между ним и его идентичностью. Современная наука выкорчовывает субъекта, и в наиболее радикальной форме последствия такого изгнания раскрываются современным искусством. Трагедия показывает отчужденного субъекта тем, кто способен бросить вызов не только своему социальному положению, но и всякому социальному позиционированию.

Шекспир напишет три свои самые значимые трагедии в начале эпохи модерна. Первая из них --- \textit{Гамлет} --- будет поставлена в 1600 году, в день, знаменующий собой рассвет современности. Трагедии \textit{Отелло} и \textit{Король Лир} также посвящены теме отчужденности. Их персонажи способны на трагическое величие потому что не вписываются в окружающий мир или не могут достичь внутренней гармонии. Даже дьявольская злоба некоторых из них проистекает из фундаментального внутреннего противоречия, теперь --- в эпоху модерна --- очевидного. Кардинальная смена смысла трагедии видна в сравнении \textit{Гамлета} с \textit{Царем Эдипом} или \textit{Отелло} с \textit{Аяксом}. В героях Шекспира выступает раскол, неведомый целеустремленным героям Софокла. Во вселенной Шекспира нельзя уже просто исполнять свой долг перед богами, как Антигона. Должно подвергать сомнению всякое должное вне зависимости от стоящего за ним авторитета. В мире Софокла есть явные противостояния, в мире Шекспира есть внутренние противоречия. Он показывает, что отчужденность не нужно преодолевать, что это основа свободы.

В \textit{Гамлете} подход Шекспира в отчужденности виден лучше всего. Гамлет --- фигура неуверенности и самокритики, и эта отчужденность дистанцирует его от диктата общественного порядка. В начале пьесы он получает приказ от умершего отца --- высшей фигуры символической власти.\footnote{При первом же появлении Гамлета видно как он далек от своего дяди Клавдия, правящего в Дании. Пока остальные празднуют победу, он остается в стороне, настаивая на дистанции от Клавдия и своей матери, вышедшей за него замуж.} Но Гамлет не приступает к выполнению приказа убить Клавдия, он усомнился, он думает, как следует повиноваться если власть Клавдия законна, и каков его статус как сына короля. Антигона знала, что должна сделать --- похоронить своего брата Полиника, не взирая на запрет Креонта под страхом смертной казни, --- и она просто делает это.\footnote{Непоколебимость Антигоны очевидна с первой сцены одноименной трагедии: она не принимает вопросы своей сестры Исмены во внимание, и, вкратце изложив ситуацию, категорично заявляет: ''Теперь ты можешь показать, кто ты: хорошая сестра или трусиха, опозорившая наших храбрых предков'' (Sophocles, \textit{Antigone}, in \textit{The Theban Plays of Sophocles}, trans. David R. Slavitt (New Haven: Yale University Press, 2007), 3). Гамлет же не только не отмахивается от вопросов других, но и сам же засыпает себя ими --- что немыслимо для античного героя вроде Антигоны.} Гамлет же ведет себя отстраненно по отношению к своему долгу и самому себе.

Из-за отчужденности Гамлет медлит с убийством Клавдия. Но это бездействие, эта неуверенность в себе и самокритика --- формы действия. Сколь бы убедительны не были объяснения этого ''бездействия'', они упускают отчужденность субъекта.\footnote{Гипотеза, что Гамлет не действует, а медлит, наиболее убедительно объяснена Зигмундом Фрейдом в \textit{Толковании сновидений} и затем более полно изложена в \textit{Гамлет и Эдип} Эрнеста Джонса (New York: Norton, 1976 [1949]): Гамлет подсознательно хочет сделать то же, что и Клавдий --- убить своего отца и заняться сексом со своей матерью.} Его действие здесь --- не через уверенность в себе, а через самоистязание, отдаляющее субъекта от ситуации и от самого себя.\footnote{В эссе о действиях Гамлета Уолтер Дэвис пишет, что вся пьеса состоит из нападок Гамлета на Клавдия (и всякого другого персонажа) с целью поставить их перед лицом травмы субъектности: ''Шекспир поместил в монолог [когда Гамлет воздерживается от убийства Клавдия...], чтобы даже неопытный читатель видел, что Гамлет делал все это время, занимаясь психотеррором и убивая людей так, как это делает его истинный преемник Яго, отравляя их души и наблюдая что с ними происходит'' (Walter A. Davis, ''Beyond Humanism and Postmodernism: A \textit{Hamlet} for the 21st) Century,'' in Shakespeare After 9/11: How a Social Trauma Reshapes Interpretation, eds. Matthew Biberman and Julia Reinhard Lupton (Lewiston: Edward Mellen Press, 2011), 280} Гамлет --- субъект эпохи модерна, потому что его действия направлены на подрыв фигуры символического авторитета и своих же убеждений.

В начале пьесы Гамлет выражает как свою собственную, так и общую отчужденность от мира. Это видно в странной темпоральности, когда Гамлет заявляет: ''Время вышло из-под контроля. О, проклятая злоба, я был рожден, чтобы все исправить''.\footnote{William Shakespeare, \textit{The Tragedy of Hamlet, Prince of Denmark, in The Riverside Shakespeare}, 2nd ed., ed. G. Blakemore Evans (Boston: Houghton Mifflin, 1997 [1600]), act 1, scene 5, lines 188–189.} В некоторым смысле здесь речь об устроенном Клавдием беспорядке: мир не в ладах конкретно с Гамлетом. Но также Гамлет говорит здесь от лица субъекта модерна вообще --- время нарушено для каждого, неприкаянность и отчужденность становятся очевидны всем, не только потерявшим своего отца. И хотя Гамлет говорит, что нужно все исправить, его действия указывают обратное --- он не верит в возвращение чувства Места. За всю пьесу Гамлет ни разу не отказывается от вопрошания --- акта, определяющего субъектность. Он пренебрегает авторитетом традиции --- авторитетом, в котором ищут убежище бегущие от отчужденности.

Как отец Гамлета приказывает Гамлету, так и традиция приказывает субъекту. Но субъект модерна, в отличие от субъекта традиции, может отреагировать сомнением вместо послушания. Гамлет не доверяет образу отцовского авторитета, и постепенно приходит к вопросам о значении самого существования. Пренебрежение авторитетом традиции провоцирует каскад сомнений, делающий Гамлета образцовым субъектом. Он чувствует, что не способен должным образом ответить на требование отца, потому что уже не вполне принадлежит его миру. Его вопрошание свидетельствует об отчужденности от мира, в котором он существует.

Неиссякаемые вопросы Гамлета определяют его неприятие отцовского авторитета. Вместо повиновения он задает вопросы, что указывает на его модерновость. Гамлет никогда заново не открывает для себя традицию. Традиция отводит ему четко определенное место, но он чужд ей. И его вопрошание не мешает ему действовать: когда Гамлет убьет Клавдия, Шекспир не представит это исполнением предначертания отца, призрак отца Гамлета нигде не появится ни до, ни во время, ни после смерти Клавдия. Убийство Клавдия здесь поступок самого Гамлета, его сомнения отделили этот поступок от изначально руководившей им власти. Он действовал, не полагаясь на авторитет --- и потому же это не бунт. Шекспир понимает, что отчужденность требует взятия ответственности за свои действия. Гамлет не строит из сомнения свою символическую идентичность, и отсутствие отца Гамлета во время убийства ясно показывает: Шекспир никогда не отказывался от заложенного с самого начала пьесы фундаментального раскола.\footnote{Можно представить альтернативный финал \textit{Гамлета}, в котором сразу после смерти Клавдия призрак отца появляется на сцене с довольным выражением лица. Наверное так было бы, пиши \textit{Гамлета} Джордж Лукас, подобно тому как в финале фильма Ричарда Маркуанда \textit{Возвращение джедая} (1983) появляются призраки Оби-Вана Кеноби (Алек Гиннесс), Йоды (Фрэнк Оз) и Энакина Скайуокера (Себастьян Шоу).} Мы никогда не вернемся из отчужденности модерна к авторитету традиции.

Популярность \textit{Гамлета} обусловлена его образцовой модерновостью. Во все времена люди пытаются вернуться к символической идентичности, но в \textit{Гамлете} Шекспир показывает невозможность найти в ней утешение. Попытки всегда заканчиваются неудачей, что демонстрирует поведение Гамлета по отношению к другим героям пьесы. Отказавшись полагаться на символическую идентичность в своих действиях, Гамлет предлагает парадигму модерна, обнажающую несостоятельность любых подобных инвестиций. \textit{Гамлет} показывает, почему у субъекта эпохи модерна не выйдет облечь себя в символическую идентичность.

В \textit{Отелло} есть две фигуры отчужденной субъектности --- Отелло и Яго. Они не остаются там, куда их помещает социальный порядок, не вписываются в рамки определенной символической идентичности. Несмотря на кажущуюся противоположность --- Отелло честно защищает установленный порядок, тогда как Яго пытается его разрушить --- оба они используют символическое для оспаривания своего Места, в пренебрежении им их траектории пересекаются.

Будучи военоначальником, Отелло поддерживает структуру венецианского общества. Отстаивая его интересы, он все чаще посещает дома его элит, включая дом Брабанцио. Там между ним и дочерью Брабанцио возникает роман, бросающий вызов расистским наклонностям общества Отелло: будучи мавром, он не видится Брабанцио подходящим зятем. Роман идет вразрез с интересами Отелло --- теперь он противоречит обществу, которое сам же защищал. Любовь к Дездемоне усиливает его отчужденность от общества и самого себя, что постепенно доведет его до саморазрушения, когда он убьет возлюбленную за воображаемую неверность.

К такому акту Отелло подведет главный злодей пьесы Яго, и в этом он будет пользоваться отчужденностью Отелло. Он убедит его, что у Дездемоны есть роман с Майклом Кассио. Зная, что не вписывается в общественный порядок, Отелло уязвим к такой манипуляции. Так осознание Яго отчужденности Отелло дает ему преимущество: Отелло никогда не подозревает Яго в двуличии. Понимая, как отчужденность структурирует отношения между людьми, Яго оказывается способен разрушить жизнь Отелло и Дездемоны.

Яго привлекателен как персонаж потому что разбирается в тонкостях обмана. Он посеял в душе Отелло сомнения о Дездемоне, в то же время заявляя, что ничего подозрительного не происходит. Такой двойной жест превосходно срабатывает на Отелло, наивном касательно эффектов обозначения. Яго говорит: ''Свои самые черные грехи дьяволы начинают с райских представлений, как это делаю и я сейчас''.\footnote{William Shakespeare, \textit{The Tragedy of Othello, the Moor of Venice}, in \textit{The Riverside Shakespeare}, 2nd ed., ed. G. Blakemore Evans (Boston: Houghton Mifflin, 1997 [1604]), act 2, scene 3, lines 351–353.} Отелло не заметит двуличия Яго, пока не убьет Дездемону за то, чего она не совершала. Но в его последней речи видно, что он осознал свою до этого ускользавшую отчужденность.

Перед смертью Отелло станет относиться к себе отстраненно и даже с отвращением, потому что это субъект, который позволил Яго обмануть его и предал любовь к Дездемоне --- и потому заслуживает только презрения. Отелло убивает его, убивая себя, при этом провозглашая: ''Прибавьте к сказанному: как-то раз/ В Алеппо турок бил венецианца/ И поносил сенат. Я подошел,/ За горло взял обрезанца-собаку/ И заколол. Вот так''.\footnote{Shakespeare, \textit{Othello}, act 5, scene 2, lines 355–356.} Это момент самопреодоления Отелло, трансценденции, возможной только через отчужденность. Расправляясь с самим собой, Отелло показывает, что осознает последствия саморазрушения как никогда раньше. В конце пьесы он наконец принимает отчужденность. В финальной характеристике себя Отелло показывает, что отчужденность --- основа современной субъектности.

Яго, в отличие от Отелло, с самого начала пьесы четко осознает отчужденность. Он знает, что обозначение неизбежно искажает говоримое, что действия всегда воспринимаются неправильно и что нельзя преодолеть этот разрыв. Он использует это свое понимание во зло. Это зло проистикает не из отказа принять отчужденность, а из интеграции низбежности отчужденности в его концепцию субъектности. Яго олицетворяет совершенную открытость эпохи модерна. Принятие отчужденности делает возможным не только освобождение, но и беспочвенное, невыразимое зло.

Зло Яго не банально.\footnote{Ханна Арендт, как известно, называет зло Адольфа Эйхмана банальным в своей книге \textit{Эйзман в Иерусалиме}. Она, конечно, ошибается, принимая за чистую монету утверждение Эйхмана что он был всего лишь партийным функционером и не питал враждебности к евреям, но в ее настаивании на банальности его злодеяний можно видеть попытку помешать ему достичь статуса Яго или Вотрена (из романа Оноре де Бальзака \textit{Отец Горио}). Арендт пишет: ''Именно явное легкомыслие --- нечто, отнюдь не тождественное глупости,--- предрасположило его к становлению одним из величайших преступников того периода'' (Hannah Arendt, \textit{Eichmann in Jerusalem: A Report on the Banality of Evil} (New York: Penguin, 2006 [1963]). 287–288). Для Арендт отнести Эйхмана к воплощенному дьяволскому злу --- значит приписать нацизму трансцендентность, которой он достичь не может.} Он делает его не для какой-то более масштабной цели. Он --- воплощение дьявольского зла, совершающий зло ради зла, не ради удовольствия или каких-то своих интересов. Кант бы сказал, что Яго обладает ''абсолютно злой волей'' и делает ''сопротивление закону'' конечным для своих действий.\footnote{Immanuel Kant, \textit{Religion within the Boundaries of Mere Reason}, in \textit{Religion within the Boundaries of Mere Reason and Other Writings}, trans. and ed. Allen Wood and George di Giovanni (New York: Cambridge University Press, 1996 [1793]), 82.} Такому субъекту как Яго просто нравится быть злым.

Кант хоть и рассматривает такое зло как теоретическую возможность, но быстро отвергает ее. Он не верит, что субъект способен желать зла ради зла как такового. Согласно Канту, возможно радикальное зло --- попытки творить добро из неверных соображений, но дьявольского зла быть не может потому что нельзя пытаться не творить добро вообще. Очевидно, Канту не доставало чтения ''Отелло'' и встречи с Ганнибалом Лектером. Шекспир предлагает убедительный портрет человека злой воли. Это возможно благодаря отчужденности субъекта. Кант же эту возможность упускает, и потому не видит, как дьявольское зло способно пролить свет на моральные действия и политическую эмансипацию.\footnote{Аленка Зупанчич утверждает, что Кант отверг такое зло ради защиты своей версии морали: согласно Зупанчич, моральный акт формально неотличим от дьявольского зла. В книге \textit{Этика реального} она напишет: ''Следуя Канту и в то же время против него, стоит указать на неотличимость дьявольского зла, высшего зла, от высшего блага и что оба они --- определения совершенного поступка. Иначе говоря, на уровне структуры этического поступка нет разницы между добром и злом --- зло формально неотличимо от добра.'' (Alenka Zupančič, \textit{Ethics of the Real: Kant, Lacan} (London: Verso, 2000), 92)}

Достигая вершин зла, Яго раскрывает ограниченность своей позиции по сравнению с таковой эмансипации: в действиях Яго отсутствует свобода. Ему нужно, чтобы Кассио и Отелло оказались в положении противостоящих ему врагов. Яго нужны враги, чтобы подрывать его авторитет, что противопоставляет его дейстельность эмансипации, которая не требует врагов. Эмансипация начинается с всеобщей отчужденности и видит свою разделенность в разделенности другого, и потому ей не нужны враги. Зло Яго не может зайти так далеко, и потому остается несвободным. Он не достигнет высот Корделии из \textit{Короля Лира}.

У Шекспира именно Лир, а не Гамлет --- олицетворение нерешительности. В начале пьесы он выражает желание выйти за пределы отчужденной субъективности и насладиться уединением. Но для субъекта не существует уединения. Как бы он не пытался уйти от проблем существования (и интриг королевства, в случае Лира), он неизбежно оказывается вовлечен в них. Внутренний раскол субъекта приводит к его неумолимой вовлеченности в мир.\footnote{То, что хочет Лир --- ''стряхнуть все заботы и дела с нашего века'' --- невозможно для говорящего субъекта (William Shakespeare, \textit{The Tragedy of King Lear, in The Riverside Shakespeare}, 2nd ed., ed. G. Blakemore Evans (Boston: Houghton Mifflin, 1997 [1606]), act 1, scene 1, line 39). Субъект не может освободиться от забот потому что он всегда вне себя, всегда в мире.} \textit{Король Лир} --- пьеса о невозможности избежать своей отчужденности, невозможности изолироваться.

В первом акте Лир отказывается от управления королевством, передавая власть своим дочерям. Решая, как разделить власть, он просит каждую дочь сказать ему, как сильно она его любит. Начинается состязание в лести, но исход был предрешен: у Лира есть фаворитка --- Корделия --- которой он планирует отдать наибольшую долю. Но все пойдет не так, как он ожидает.

Лир не понимает, что и он, и его собеседники --- субъекты, отчужденные от того, что говорят. В ответ на требование выразить любовь, две его дочери --- Гонерилья и Регана --- льстят ему, а Корделия --- которая искренне любит Лира и понимает, что субъекты не способны выразить мысли напрямую, по команде --- разочаровывает Лира, в ее ответе нет риторических излияний ее сестер.

Корделия, в отличие от Лира, отлично осознает отчужденность субъектности. Она демонстрирует свою любовь к Лиру именно тем, что не превращает ее в спектакль. Косвенность необходима, потому что субъектность отчуждена от значения. Она отвечает отцу: ''Что сказать Корделии? Любовь, и молчать''.\footnote{Shakespeare, \textit{King Lear}, act 1, scene 1, line 62.} Лир упрекает Корделию в недостатке выразительности, и она еще дальше отходит от прямого ответа: ''Я так несчастна, что не могу выразить свое сердце словами. Я люблю ваше величество согласно моим обязанностим --- ни больше, ни меньше''\footnote{Shakespeare, \textit{King Lear}, act 1, scene 1, lines 91–93.} Отчужденность Корделии --- и ее понимание этой отчужденности --- мешает ей изливать душу и льстить как это делают ее сестры.

Отказ Корделии предать свою отчужденность --- и вообразить себя идентичной символическому --- делает ее главной героиней \textit{Короля Лира}. Она отказывается вести себя как низведенная до положения лишь дочери, как ее сестры. Отчужденность мешает ей играть ту роль, которую требует от нее отец. Пьеса заканчивается примирением Лира с любимой дочерью, но они умирают сразу после. Нежелание Лира признать свою отчужденность обрекает его на изгнание. На примере Лира пьеса показывает ущерб, который приносят попытки бежать от отчужденности. Этическая сущность Корделии проистекает из ее стойкого принятия своей отчужденной субъектности.

В каждой из этих трагедий видна несводимость субъекта к условиях, в которых он появился и находится, их герои не вписываются в свои ситуации --- Гамлет сомневается в авторитете умершего отца, Лир не может убежать от проблем --- герои Шекспира демонстрируют отчужденность субъекта, его дистанцию от самого себя и от общества. Трагедии Шекспира указывают на эмансипацию, подчеркивая неизбежность отчужденности.

\section{''Я мыслю там, где меня нет''}

Рене Декарт сделал в философии то, что Шекспир сделал на сцене. Он не использовал термины \textit{субъект} и \textit{отчужденность}, но фактически был первым философом, изучавшим отчужденность субъекта --- и, значит, первым философом модерна. Для Декарта он значит отказ от авторитета [философской] традиции. Это отчуждение не мешает мышлению, а позволяет ему ступить на новый путь.\footnote{Именно явное отчуждение от авторитета традиции делает Декарта первым философом модерна. Мешель де Монтень --- его главный конкурент за этот титул --- не указывал на этот разрыв.}

Декарт, как известно, начинает свой проект с опыта и вопросов о нем, а не с дискуссий, как в схоластической традиции. В \textit{Рассуждениях о методе} он подчеркивает свою отчужденность от традиции, заявляя, что думает, не полагаясь на какие-либо философские предпосылки: ''Ибо мне казалось, что гораздо больше истины в рассуждениях, которые человек делает по поводу касающихся его вопросов, чем в тех, которые какой-нибудь ученый делает в своем исследовании о чем-то спекулятивном''.\footnote{René Descartes, \textit{Discourse on the Method of Rightly Conducting One’s Reason and Seeking the Truth in the Sciences}, trans. Robert Stoothoff, in \textit{The Philosophical Writings of Descartes}, vol. 1 (Cambridge: Cambridge University Press, 1985 [1637]), 115.} Не изучая, но размышляя, Декарт занимает негативную позицию к общепринятой мудрости и развивает свою философию, относясь к традиции как к чему-то чуждому.

Метод Декарта предполагает переход от установившейся уверенности к фундаментальному сомнению. Подобно Гамлету, Декарт посвящает себя проекту радикального сомнения, отчуждаясь от всякой определенности традиции. Декарт отказывается от чтения текстов философов и формулирует свою философию на основе неприятия всего, что в его обществе считается истинным. Это философское выражение субъектности, проявляющееся в радикальном отказе от условий, которые могли бы конституировать субъекта. Быть субъектом, в понимании Декарта, значит отбросить лестницу, на которой стоишь, и так остаться как бы подвешенным в воздухе --- ничем не удерживаемым и не защищенным от падения. Субъект отчужден от всего.

Отчуждение субъекта распространяется и на него самого. Проект радикального сомнения, который Декарт опишет в \textit{Размышлениях о первой философии},--- это борьба с самообманом, и она требует отчужденности от себя: ''Несколько лет назад я был поражен, сколько лжи принимал за истину в детстве, и сколь сомнительно то, что я на ней построил. Я понял, что придет пора все разрушить и начать с основ''.\footnote{René Descartes, \textit{Meditations on First Philosophy}, trans. John Cottingham (Cambridge: Cambridge University Press, 1986 [1641]), 12.} Мышление Декарта продиктовано ложью, но он способен остановиться и оспорить ее. Это возможно только потому, что он как субъект не полностью определяется своей ситуацией.

Радикальная отчужденность Декарта обыгрывается в известном анекдоте о том, как он выпивал в баре. После нескольких кружек к нему подошел бармен и спросил: ''Не хотите ли еще?'', на что Декарт ответил: ''Не думаю'',--- и исчез. Декарт связывает свое существование с мыслящим субъектом, и обыгрывая  эту связь, анекдот раскрывает разрыв между субъектом и его Местом. Даже выпивая в баре, Декарт --- как мыслящий субъект --- всегда где-то \textit{не здесь}. Он всегда может если и не исчезнуть, то абстрагироваться от ситуации.

Свобода субъекта от ситуации --- его ключевая особенность. Она основана на его отчужденности от любого Места. Декарт видит прямую связь между свободой субъекта от власти и его равенством с прочими субъектами --- где бы они сами не находились. Свобода от контекста в равной степени присуща всем субъектам. Иерархия строится на контексте, но субъектность --- как ее понимает Декарт --- это отчужденность от своего контекста. Быть субъектом --- не принадлежать какому-либо Месту, нигде не чувствовать себя как дома.

Теоретическая несводимость к какому-либо Месту делает Декарта образцовым философом модерна. Отношение к Месту определяет его отношение к разным народам: Декарт считает приобретаемые отношения и привычки полностью зависимыми от места его рождения и взросления. Поэтому установки и привычки человека не имеют ничего общего с его субъектностью, составляя лишь символическую идентичность. Если бы немец рос в Китае, он бы не вел себя как немец. Для Декарта это следствие абсолютной свободы субъекта --- отношения и привычки внутренне с ним не связаны. Субъект способен ''отказываться от всех мнений, которые до сих пор принимал'', и так демонстрирует свое негативное отношение к породившему его контексту.\footnote{Descartes, Discourse on the Method, 118.} Декарта традиционно видят фигурой западного высокомерия и безразличия к инаковости, но его концепция отчужденности субъекта позволила ему сформулировать раннюю версию мультикультурного эгалитаризма.

Пожалуй, сильнее всего Декарт отошел от традиции, когда решил писать на родном языке, а не на латыни --- языке всей схоластической философии. Само существование оригинальной французской версии его текста знаменует его утверждение субъектности через неприятие традиционного авторитета. Крупная философская работа на французском языке --- почти беспрецедентный поступок в то время, показывающий намерение Декарта не ограничиваться высокообразованным классом, но обращаться ко всей грамотной публике. Как субъект отчужден от породившего его контекста, так и этот текст отчужден от стоящей за ним схоластической традиции.

Стремление Декарта к отчужденности проявится не только в его работах: он будет много странствовать, впервые покинет Францию в 22 года и проведет за границей более половины взрослой жизни. Он откликнется на приглашение королевы Кристины и неохотно посетит Швецию, когда ему будет чуть за 50, и умрет там в 1650г.. Отсутствие жилья и холод в Швеции в конечном счете погубят Декарта, но они же будут подталкивать его к размышлениям. В этом они схожи с Шекспиром.

\section{Предатели модерна}

Прорыв к модерну встречал серьезное сопротивление. Сопротивлялись в основном сторонники традиции --- те, кто отстаивает идентичность, укорененную в определенном Месте. Они не рассмотривают отчужденность как всегда уже существовавшую --- вселенная модерна, согласно ним, навязывает людям отчужденность и вытесняет их из Мест, в которых они могли бы жить по-настоящему. Иногда они реагируют террором, иногда --- установлением теократических режимов, иногда даже проектами геноцида. Они продвигают возвращение к традиции как единственный способ противостоять ужасам современности.

Эти призывы почти всегда неискренни. Сторонники традиционных ценностей стремятся к ним не жертвуя технологическими достижениями современности. Они, значит, не способны диалектически осмыслить эпоху, против которой выступают, не видят, что пороки современности тесно связаны с ее достижениями. Потому они и сохраняют спокойствие в своем реакционном неприятии модерна.

Реакционные критики модерна делают отчужденность основанием для своих нападок. Так, философия Мартина Хайдеггера --- призыв к восстановлению того, что было уничтожено современностью. Хайдеггер убедительно излагает традиционализм, основывая его на специфике Места. Модерн же делает все Места одинаковыми, однородными, так что нельзя больше жить в своем Месте. Хайдеггер связывает жилище или нахождение дома с формой бытия, которую следует принять смертным: наша привязанность к Месту соответствует нашей конечности, и потому ценна. В эпоху модерна сохранять такую привязанность становится все труднее. Современность --- уверяет Хайдеггер --- обманывает нас, заставляя погружаться в бесконечность и закрывать глаза на определяющую наше бытие конечность. Потерять свое Место --- значит потерять связь с конечностью.

Согласно Хайдеггеру, наиболее распространенная форма отказа от конечности --- тенденция считать себя субъектами. Субъект --- тот, кто не имеет Места. Кант использовал термин \textit{субъект}, чтобы подчеркнуть нашу разобщенность по отношению к миру, а Хайдеггер отвергает его, чтобы подчеркнуть нашу фундаментальную принадлежность к миру. Многие комментаторы отмечают явное неприятие Хайдеггером термина \textit{субъект}, но никто из них не указывает, насколько это важно для его антимодернистской философии и для его неприятия наследия немецкого идеализма, взявшего субъектность отправной точкой.\footnote{Хайдеггер отвергает субъекта, о котором говорил Кант, но главный его оппонент --- Гегель, совершенно беззастенчиво превозносящий освобождающую силу субъектности. Если Кант ввел этот термин в философию, то Гегель превратил его в первое слово в философском лексиконе модерна.} Хайдеггер стремиться избавиться от субъекта, потому что, как он чувствует, от него несет модерном.

Хайдеггер говорит не о субъекте, а о \textit{Dasein}, обозначая так существ, способных к самоанализу. Он отказывается говорить о субъектности потому же, почему Кант и Гегель о ней говорят: это предполагало бы существование нечта отдельного от своего мира. Dasein же значит бытие в мире, бытие, неотделимое от своего мира. Хайдеггер использует два разных значения термина Dasein: обычно он значит ''существование'', но две части слова --- da и sein --- значат ''здесь'' и ''бытие''. Мыслить в терминах Dasein --- значит рассматривать сущность как помещенную в мир, как существующую через свое пребывание в мире.

Наиболее содержательно Хайдеггер критикует понятие субъекта в \textit{Бытии и времени}, где противопоставляет два подхода к существованию: отношение к объектам как к подручным (\textit{zuhanden}) --- данных лишь в их использовании нами, и как к наличным (\textit{vorhanden}) --- данных как теоретические сущности, отделенные от использования. В \textit{Бытии и времени} Хайдеггер отвергает приоритет наличного над подручным, говоря, что теоретический подход к миру вводит в заблуждение.

Замена субъекта на Dasein делает его подручной сущностью. В \textit{Бытии и времени} Хайдеггер напишет: ''Онтологически определять ''я'' как ''субъект'' --- значит рассматривать его как наличное''\footnote{Martin Heidegger, \textit{Being and Time}, trans. John Macquarrie and Edward Robinson (San Francisco: HarperCollins, 1962 [1927]), 367.}. Согласно Хайдеггеру, это худшее, что можно сделать: человек так превращает вовлеченное присутствие в мире в нечто от мира удаленное и никак в нем не присутствующее. Субъективизм --- отказ от взаимодействия с миром. Это оплот философов, испугавшихся существования.

Хайдеггер избегает говорить не только о субъекте, но и об отчужденности. Казалось бы, если Dasein --- сущность, заброшенная в мир, то можно говорить о его отчужденности больше и положительней. Ведь заброшенность не дает Dasein овладеть миром, объектами в нем или самим собой: Dasein никогда не над миром, никогда не за его пределами, всегда смешан с ним.

В критике модерна Хайдеггер сосредотачивается на иллюзии господства от идеи субъективности. В эпоху модерна мы теряем связь со своей конечностью и верим, что все может нам подчиниться. Но, согласно Хайдеггеру, мы не хозяева даже самим себе. В \textit{Канте и проблемах метафизики} он напишет: ''В зависимости от того, чем он не является, человек не хозяин того, чем является''.\footnote{Martin Heidegger, \textit{Kant and the Problem of Metaphysics}, 5th ed., trans. Richard Taft (Bloomington: Indiana University Press, 1997 [1929]), 160.} Невозможность владеть самим собой согласуется с отчужденностью, но Хайдеггер не пойдет в этом направлении. Он будет критиковать иллюзию господства, но ошибочно свяжет ее с субъектностью и модерновым вытеснением человеческого. Согласно Хайдеггеру, нашей нехватке господства соответствует пребывание в Месте, а не отчужденность.

Хайдеггер будет упомянать отчужденность только чтобы критиковать ее. Отчужденность для него --- не освобождение от Места, а прискорбная удаленность от него. Задача мышления, он убежден, в облегчении возвращения на Место, в преодолении отчужденности модерна.

В \textit{Бытии и времени} Хайдеггер связывает отчужденность с неаутентичностью, в которую впадает Dasein. Это потеря главного для нас: ''когда Dasein, убаюканный ''пониманием'' всего, сравнивает себя со всем, он движется к отчуждению [\textit{Entfremdung}], где от него скрыта наибольшая возможность-к-Бытию. Исчезающее бытие-в-мире не только заманчиво и успокаивает, но и отчуждает''.\footnote{Heidegger, \textit{Being and Time}, 222.} На языке Хайдеггера, погружение в спокойствие повседневности --- это отчуждение. Как и многие левые, он будет критиковать отчуждение модерна, но винить в нем не капиталистическую экономику, а современную науку и современность вообще.

Философский проект Хайдеггера нацелен на переосмысление конечности и Места, на противостоянии отчужденности. Возвращение к Месту становится ключом к восстановлению понимания бытия в мире. Из-за однообразия зданий в модерне в них невозможно жить, потому что так нельзя определить свое Место: мы лишаемся его специфики, здания во Фрайбурге выглядят так же, как в Джакарте. Люди в них живут и не могут в них жить. В них они забывают свое Место. Эти здания --- чужеродные сосуды, совершенно неадекватные Dasein и его бытию-в-мире.\footnote{В своем глубоком изложении философского проекта Хайдеггера Джефф Малпас напишет: ''Человек --- это, по сути, бытие-в-Месте так же, как и бытие-в-мире'' (Jeff Malpas, \textit{Heidegger and the Thinking of Place: Explorations in the Topology of Being} (Cambridge: MIT Press, 2012), 63). Малпас так свяжет Место с миром.} Модерн не дает нам Места для бытия.

Многие антимодернисты плохо разбираются в философии Хайдеггера, а между тем он был приверженцем позиции, отражающей основу антимодернизма. Внимание Хайдеггера привлек нацизм, обещавший возвращение к \textit{Blut und Boden} (''крови и почве''). И хотя философия Хайдеггера была безразлична к крови, почва была в ней ключом к ощущению Места. Бытие Dasein всегда укоренено в почве, и этой укорененности угрожает модерн. Нацизм же обещает покончить с извращением, запутыванием нашего бытия-в-мире и создать Место для жизни.\footnote{Хайдеггер объясняет привлекательность нацизма для себя в \textit{Введении в метафизику}, где указывает на ''внутреннюю истину и величие этого движения'' в противостоянии ''столкновению глобальной технологии и современного человечества'' (Martin Heidegger, \textit{Introduction to Metaphysics}, 2nd ed., trans. Gregory Fried and Richard Polt (New Haven: Yale University Press, 2014 [1953]), 222). Хайдеггер видит в нацизме надежду на борьбу с модерном за восстановление ощущения Места. Национализм партии, пропаганда арийской идентичности и безжалостность к врагам --- все это является частью сделки.}

Мало кто из нацистов читал Хайдеггера, и ему не удалось построить карьеру нацистского философа. Но им двигал тот же антимодернистский импульс, что и нацизмом и другими реакционными политическими философиями. Хайдеггер --- великий философ антимодерна --- пытается утвердить Место в мире вопреки срыву Вселенной с него в эпоху модерна.\footnote{Хайдеггер --- один из основоположников экзистенциализма, но его антимодернизм отделяет его от Жан-Поля Сартра и Симоны де Бовуар. Сартр и Бовуар принимали модерн, и настолько, что называли себя попутчиками коммунизма --- позиции, которую Хайдеггер будет отвергать и называть симптомом модерновой неприкаянности.}

Модерн порождает реакционеров, отвергающих проблеск отчужденности. Они отказываются от нее и обещают своим приверженцам ее преодолеть. Но отчужденность --- не плод революции модерна, она всегда уже присуща субъекту, и потому реакционное возвращение к истокам всегда заканчивается неудачей. Возвращаться некуда, и не потому, что модерн разрушил Место, а потому что никакого Места не было.

\section{Возвращение к нигде}

Попытавшись найти Место, к которому могли бы вернуться, мы неизбежно попадемся на удочку реакционной программы --- продвигающей нацию, общину или религиозную группу как то самое Место, где мы сможем чувствовать себя как дома. Эмансипация же всегда идет в обратном направлении. Есть только визжащая тишина свободы, и проблема человека --- не про обитание в конкретном Месте, а про неизбежность отчужденности от всякого такого места.

Этот контраст хорошо виден в двух научно-фантастических работах. Первая --- обновленный телесериал \textit{Звездный крейсер Галактика}, шедший с 2004 по 2009 год и широко известный благодаря политической остроте.\footnote{Оригинальная версия \textit{Звездного крейсера Галактика} выходила всего один сезон (1978-1979). В новом сериале большинство персонажей остались в основном теми же, но с некоторыми изменениями. Впрочем, не будем об оригинале.} В сериале показано противостояние между людьми и сайлонами --- группой машин, стремящихся уничтожить людей и утвердиться в качестве основной силы в своей части галактики. В начале сериала сайлоны внезапно атакуют и уничтожают почти всех людей, остается лишь несколько кораблей и одноименный крейсер \textit{Галактика}. Крейсер уводит уцелевшие корабли подальше от разрушенных колоний, продолжая отбиваться от сайлонов. Выжившие люди ищут то, что называют потерянной колонией, Землю, где надеются обрести дом.\footnote{После террорестических актов, произошедших в США 11 сентября сериал будет представлять своих героев --- людей --- как повстанцев, прибегающих к терактам и даже самоубийственным бомбардировкам в плену у сайлонов.}

В последнем --- четвертом --- сезоне, после многих испытаний и лишений, люди, наконец, найдут Землю. Но вместо дома их встретят безжизненные радиоактивные руины. Неужели людям придется научиться жить без дома? Слава Богу, нет, ведь это только середина сезона. Поиски продолжаются, и люди в итоге оказываются на гораздо более пригодной планете, которую и назовут Землей. Здесь они наконец чувствуют себя как дома, обживаются и называют себя предками людей, населяющих Землю сегодня. Сериал заканчивается прибытием домой. И пусть это даже не то место, куда изначально направлялись, это Место. В нем люди теперь сыграют основополагающую роль в развитии разумной жизни. \textit{Звездный крейсер Галактика} покажет, что нужно идти на все ради места, где сможешь чувствовать себя как дома.\footnote{\textit{Звездный крейсер Галактика} --- один из немногих сериалов, испорченных концовкой. Если бы он закончился на первой --- пустынной и непригодной --- Земле, он бы стал одним из лучших когда-либо снятых телесериалов.}

Акцент на поиске дома делает \textit{Звездный крейсер Галактика} антимодернистским произведением. Если бы он закончился на прибытии к разрушенной Земле, это создало бы ощущение отсутствия дома, показало необходимость принять отчужденность. Но заключительные эпизоды развернули сериал в сторону политической позиции Мертина Хайдеггера.

Второе произведение --- фильм \textit{Темный город} (Алекс Пройас, 1998) --- не стало таким популярным. Фильм вышел на экраны в 1998 году и остался почти незамеченным. Он рассказывал о группе инопланетян (известных как \textit{Незнакомцы}), манипулирующих жителями городами, переделывая его и меняя личности жителей каждую ночь: можно лечь спать бедным и одиноким в убогой квартире, а проснуться богатым с семьей в роскошном особняке. Мизансцена фильма ясно покажет, что личность --- продукт внешних манипуляций. Незнакомцы так ищут человеческую душу, то, что делало бы человека уникальным. Один персонаж --- Джон Мердок --- заметит манипуляции, проснувшись, когда они происходят. Окажется, что он тоже может манипулировать реальностью, а значит, может бросить вызов Незнакомцам и их власти над городом. Но на протяжении большей части фильма он даже не пытается это сделать --- вместо этого он фантазирует, как сбежит из города в Шелл-Бич, который считает домом своего детства. Эта фантазия делает его послушным поданным.

В конце фильма Мердоку все уже удастся сбежать из города и отправиться туда, где, он уверен, находится Шелл-Бич. Он сломает стену в конце туннеля, и обнаружит, что никакого Шелл-Бич нет: через дыру в стене видно лишь бескрайнее пространство. Фильм показывает дом Мердока вымышленным, и события в нем не привязаны к какому-либо месту --- даже к Солнечной системе.

Мердок, впрочем, не отчается, а восстанет против Незнакомцев, победит и перестроит мир фильма. \textit{Темный город} показывает, что дома нет и единственный выход --- рассмотреть эту отчужденность как свободу. Мердок не способен вернуть жителей города в их дома, потому что они навсегда утеряны. Они обречены провести жизнь, дрейфуя в космосе. Но это же --- условие их свободы. Как только Мердок отказался от возвращения домой, он обрел возможность освободить и себя, и остальных, кто находится под властью Незнакомцев. Победив Незнакомцев, люди остаются один на один со своей отчужденностью. У них нет надежной символической идентичности, нет того, на что можно опереться. Таково освобождение от внешнего контроля. Фильм показывает, что свобода приходит через саму отчужденность, а не через ее преодоление.

Мораль \textit{Темного города}, в отличие от \textit{Звездного крейсера Галактика}, в том, что неуместность неизбежна. Отчужденность от места становится основой для создания эгалитарного общества без какой-либо деспотичной иерархии. Мы отчуждены в свободные и равноправные отношения, но упускаем эту возможность, когда берем обязательство преодолеть отчужденность и найти место, где чувствовали бы себя как дома.

Модерн выводит проблему субъекта на первый план, и потому его всячески сторонятся антимодернистские традиционалисты и просто те, кто против видеть в отчужденность неизбежной. Такое затворничество лишают возможностей, предложенных вселенной модерна. Свобода --- синоним принятия отчужденности.

\chapter{Природа угнетения}

\section{Незримые деления}

Современность порождает столько ужасного, что едва ли выйдет хвалить модерн, не сделав оговорок. Расистское рабство, этнические чистки, капиталистическая эксплуатация, атомная бомба.\footnote{Юрген Хабермас уверен, что все зло проекта модерна --- от неспособности его завершить. В \textit{Modernity and Postmodernity} он рассматривает подходы к критике модерна --- от традиционного консервативного до неоконсервативного и левой софистики. Он утверждает: ''Я думаю, что вместо отказа от модерна и его проекта как от безнадежного, нам следует учиться на ошибках тех экстравагантных программ, которые пытались свести модерн на нет'' (Jürgen Habermas, “Modernity versus Postmodernity,” trans. Seyla Benhabib, New German Critique 22 (1981): 11). Хабермас, конечно, прав, отвергая то, что называет постмодернизмом (под ним он подразумевает левых философов, берущих пример с Фридриха Ницше и Мартина Хайдеггера), но он недооценивает, насколько ужасы модерна действительно модерну и принадлежат. Сталинизм, например, нельзя списать лишь на победу антимодернистских сил --- он возник в попытке модернизации любой ценой. Впрочем, Хабермас находится на верном пути, указывая причиной всех этих ужасов незавершенность проекта модерна.} Некоторые из ужасов современности, конечно, результат антимодернистских реакций, но и они все же происходили в эту эпоху. Импульс модерна хотя бы отчасти виновен в них.

Стоит ли модерн всего того, что за него приходится платить? Не следует ли, учитывая дороговизну, разработать программу против него?\footnote{Бруно Латур выступает против модерна, не занимая антимодернистскую позицию, а отказываясь даже признавать его. Для Латура модерн --- лишь еще одна традиция. См. Bruno
Latour, We Have Never Been Modern, trans. Catherine Porter (Cambridge: Harvard University Press, 1993 [1991]).} Модерн уже нанес огромный ущерб и представляет самую серьезную угрозу жизни на планете. Отношение к отчужденности поможет разобраться в проблемах, которые вносит модерн, и понять, стоит ли поддерживать его проект.

Может показаться, что все эти ужасы --- от отчужденности, столь явной в модерне. И большинство теоретиков так и считают: отчужденность в модерне ведет к неизбежному разрушению природного мира, позволяет солдатам безнаказанно уничтожать этнические группы, поощряет владельцев бизнеса выжимать из своих работников все, на что они способны, разрешает пилоту бомбардировщика мгновенно убить около ста тысяч человек без малейшего сожаления\footnote{Именно так обстоит дело с Полом Тиббетсом --- пилотом, первым сбросившим атомную бомбу. Подробный отчет о психологических причинах его состояния (в том числе от самого Тиббетса) см. в Walter A. Davis, The Holocaust Memorial: A Play about Hiroshima (Bloomington: 1st Books, 2000).} --- и это лишь некоторые достижения модерна. Изучение истории почти любого сделает постмодернистом --- так же как посещение скотобойни сделает веганом даже заядлого мясоеда. Кажется, отчужденность в модерне виновна во всех его грехах, но при ближайшем рассмотрении оказывается, что все совсем иначе.

Ужасы модерна не могли возникнуть из порожденной им отчужденности, потому что отчужденность была всегда. И даже не считая себя отчужденными, люди всегда ими были. Современность лишь заставляет осознать отчужденность субъекта, показывает его неуместность, она вскрывает отчужденность, а не создает ее. И этот процесс освобождает, а не угнетает --- без него субъект оставался бы пленником традиции.

В эпоху модерна все идет наперекосяк не из-за отчужденности, а от попыток убежать от нее. Коллективный отказ принимать отчужденность как неотъемлимую черту субъектности приводит ко всем страшным грехам, которые видела современность. Угнетение возникает в порыве создать образ Места и избежать травмы отчужденности. Каждый акт и каждая структура угнетения коренятся в попытке избежать столкновения с фундаментальной отчужденностью. Мы не боремся с последствиями отчужденности, потому что отрицаем ее. Субъекты модерна открывают ящик Пандоры в попытках укрыться в нем. Необходим не отказ от модерна, а еще большее принятие того, что он вскрывает.

В ответ на откровение модерна --- обнажение отчужденности --- люди бежали от него к невиданным формам варварства.\footnote{Утверждая в \textit{The Better Angels of Our Nature}, что мир модерна куда менее жесток, чем предшествующий ему, Стивен Пинкер не учитывает массовое насилие. Конечно, европейский город в 2020г гораздо безопасней того же города в 1520г, но такой уровень безопасности обеспечен невидимым простому горожанину структурным насилием: добычей полезных ископаемых в ужасающих условиях в Конго или где-нибудь еще, опосредованных (''прокси-'') войнах по всему миру и тщательно патрулируемой границе --- и это лишь самые заметные его проявления. Современность скрывает насилие, но не уменьшает его. Его плоды используются для ухода от травмы отчужденности, от отсутствия Места.} Модерн только умножил ужасы прошлого, потому что не замечать отчужденность стало трудней. Насилие капитализма, расизма и сексизма, не говоря уже о современных войнах и разрушении планеты --- все это ответ на явность отчужденности в модерне.\footnote{Современность дает больше возможностей психологического избегания отчужденности: человек может совершать насильственные действия, отрицая их реальность. Зигмунд Фрейд назовет это фетишистским отрицанием. Конечно, такое отрицание работало и до модерна, но в нем оно становится невероятно популярным.} Это насилие превосходит то, что было до модерна, потому что угроза потери мнимого Места велика как никогда. Насилие используется для удержания людей в определенном символическом положении и так воспроизведении эффекта Места. Но в модерне отчужденность слишком очевидна чтобы ее возможно было скрыть.

Освобождающий характер отчужденности остался перекрыт порывами связать ее с угнетением. В \textit{Диалектике просвещения} Макс Хоркхаймер и Теодор Адорно свяжут варварство модерна с самоотчуждением: ''Люди стали столь отчуждены от самих же себя и природы, что теперь знают только, как эксплуатировать друг друга и причинять друг другу вред''\footnote{Max Horkheimer and Theodor W. Adorno, Dialectic of Enlightenment: Philosophical Fragments, trans. Edmund Jephcott (Stanford: Stanford University Press, 2002 [1947]), 211.}. Они не видят в отчужденности свободы, для них наша дистанцированность от себя и от других лишь дарует способность к жестокости.\footnote{Ханна Арендт, хоть и яростно противопоставляет себя Адорно как в теоретическом, так и в личном плане, все же разделяет его критику отчужденности. Это согласие показывает, как далеко зашел критический консенсус. В \text{Origins of Totalitarianism} она напишет: ''К тоталитарному господству в нетоталитарном мире людей толкает одиночество --- когда-то переживаемое лишь в маргинальных социальных условиях вроде старости, теперь ставшее повседневным для целых масс. Безжалостный процесс, в который тоталитаризм вовлекает и организует массы --- бегство от этой реальности'' (Hannah Arendt, The Origins of Totalitarianism (New York: Harcourt, 1968 [1951]), 478). То есть отчужденность порождает одиночество, а одиночество провоцирует тоталитаризм.}

Адорно и Хоркхаймер не понимают, что разрушение самоидентификации --- не то, что следует преодолеть --- это основа нашей свободы, нашего равенства, нашей солидарности и нашего творчества. Без отчужденности мы подчинены социальному контексту, заключены в заданную извне идентичность --- полностью определены внешним. Отчужденность как освободительная сила делает возможной политизацию. Принятие отчужденности стоит рассмотреть как привлекательную альтернативу животной определенности естественного существа, а не обрекать себя на сизифов труд бегства от него.

Потребуется переосмыслить угнетение. Его, как не странно, нельзя приравнять переживанию отчужденности, потому что последняя --- универсальна. Угнетение не обрекает на отчужденность и не подталкивает к ней, напротив --- дарует как угнетающим, так и угнетаемым их символическую идентичность, воспроизводит Место. Репрессивная система привлекает угнитателя тем, что обеспечивает каждому конкретное социальное положение, как если бы это положение и было его Местом. Угнетатель и угнетаемый становятся не более чем символическими идентичностями. Цель угнетения --- превратить субъектов в самоидентичные сущности. Но члены общества угнетения не становятся волшебным образом такими самоидентичными существами: даже в самом деспотичном обществе никто не может избежать субъектности и ее синонима --- отчужденности.\footnote{Иногда репрессивные силы открыто выражают желание низвести угнетаемых субъектов до уровня животных. Но даже самая репрессивная система не способна покончить с расколом субъекта, с бессознательным угнетаемого. Репрессивные системы никогда не имели успеха. Даже у угнетателя есть подсознание --- даже он желает чего-то \textit{другого}.} Быть угнетенным --- значит быть подталкиваемым к самоидентификации без явного пути выхода из нее. Угнетение работает на сокрытие самораскола, хотя само опирается на него.

Большинство хорошо знают, как выглядит угнетение: социальный порядок, обрекающий массу людей на нищету, на приниженное по сравнению с другими положение, остракизм. Угнетать --- значит классифицировать как нечто целое, не позволять существовать как противоречивому субъекту с желаниями, расходящимися с его собственными.\footnote{Как Национальное Объединение относится к иммигрантам-мусульманам во Франции? С точки зрения этой политической партии, они не могут быть отчужденными (хотя легко представить, насколько отчуждающа жизнь мусульманина в современной Франции, учитывая предрассудки, с которыми ему приходится столкнуться), напротив --- их желания понятны: они отвергают французскую светскость, отрицают равноправие женщин, хотят навязать нации законы шариата и уничтожить католическое наследие страны. Эта репрессивная точка зрения не допускает отчужденности иммигрантов-мусульман, их противоречивых представлений о светскости, принудительном равенстве и Франции. Угнетатели видят в своих жертвах людей, совершенно отождествляющих себя со своей идентичностью. Иммигранты --- это просто иммигранты, у них нет иных желаний, чем желания иммигранта, считают они.} Угнетатель действует так, будто он и другой --- самотождествевнные существа.

Один из признаков угнетения --- сведение группы к коллективной идентичности с определенными желаниями и способами существовать. Все субъекты сводятся к последовательной символической идентичности. Идентичность подавляет субъектность, и потому так часто угнетенных идентифицируют коллективно: ''От них плохо пахнет'', ''Они ленивы'' или даже ''Они разрушают нашу нацию''. Коллективизируя других, человек отказывается признавать их самоотчужденность.

Важно, что коллективная характеристика может быть и положительной: ''Все матери любят своих детей''. Кажется, что здесь нет ничего сексистского, ведь высказывание однозначно позитивно --- но оно приписывает всем родившим женщинам единую символическую идентичность. Женщины, согласно этому высказыванию --- не субъекты, а цельные существа, ведь будь матери субъектами, некоторые любили бы своих детей, но некоторые --- нет, и даже те, что любили, делали бы это не всегда. Субъект не идентичен самому себе, а его желания противоречивы, но силы угнетения пытаются лишить своих жертв этой отчужденности. Угнетение направлено на овеществление.

Модерн вскрывает отчужденность, но субъект модерна обычно отступает от этой явленности, прячется в символической идентичности. Угнетение в современном мире проявляется через это отступление --- оно стремится скрыть отчужденность, убедить субъектов, что для них возможна цельность, возможно Место. Принятие всеобщей отчужденности направлено против угнетения.\footnote{Иначе говоря, борьба за аутентичное субъектности существование и борьба за равенство --- это, по сути, одно и то же. И отрицать отчужденность субъекта --- значит отказываться от борьбы за равенство.}

Может показаться, что свойственная субъекту отчужденность и отчужденность в капиталистическом обществе --- разные. Но социальные формы отчуждения всегда лишь манифестации фундаментальной отчужденности субъекта. Различные виды отчуждения обнаруживают дистанцию в субъекте --- между ним и его идентичностью. При капитализме, мое отчуждение от продукта моего труда напоминает мне о моей фундаментальной отчужденности как субъекта --- и тем не угнетает, а освобождает. Проблемно представление об отчужденности в капитализме как о проблеме --- именно оно провоцирует угнетение.

\section{Капиталистическое будущее}

В капиталистическом обществе субъекты испытывают отчуждение: они работают и производят товары, которые приносят прибыль другим, оставляя им только нехватку. Они постоянно видят себя в стороне от удоволствия, которое, они уверены, получают другие, и надеются найти его в том, чего у них еще нет. Вместо накопления богатства, большинство накапливают долги. Капитализм так сильно идеологизирует людей не отчужденностью, а обещанием лекарства от нее. В капиталистическом обществе каждая сделка символизирует прерывание отчужденности.\footnote{В заключительных строках своего первого романа \textit{Сестра Кэрри} Теодор Драйзер опишет положение капиталистической субъектности. Когда события романа подойдут к концу, Кэрри Мибер обнаружит, что все еще стремится к чему-то большему даже поднявшись по социальной лестнице. Рассказчик комментирует положение Кэрри: ''В своем кресле-качалке, у окна, мечтая, ты будешь тосковать в одиночестве. Будешь мечтать о счастье, которого, возможно, никогда не испытаешь'' (Theodore Dreiser, Sister Carrie (New York: Bantam, 1982 [1900]), 400)} Капитализм являет отчужденность, но в то же время обещает лекарство от нее. Он представляет ее проблемой, проблемой решаемой, а не самой структурой субъектности. Это делает капитализм привлекательным даже для тех, кто не извлекает из него материальной выгоды.\footnote{Адам Смит утверждает, что капитализм примиряет индивида с общественным порядком: индивиды преследуют личные интересы, которые в конечном счете совпадают с интересами общества. Смит пишет: ''Каждый прилагает усилия, чтобы найти наиболее выгодное применение имеющемуся у него капиталу. Конечно, он ищет выгодное себе, а не обществу. Но этот поиск естественным образом --- даже с необходимостью --- заставляет его предпочесть ту работу, на которой он будет наиболее выгоден обществу'' (Adam Smith, An Inquiry into the Nature and Causes of the Wealth of Nations (Hamburg: Management Laboratory Press, 2008 [1776]), 343–344). Ясней всего отчужденность субъекта видна в его несоответствии общественному порядку, в котором он живет --- но, согласно Смиту, это не верно для капитализма. И эта теория просачивается в практику, капитализм скрывает нашу отчужденность, убеждая, что наши желания соответствуют тому, чего требует социальный порядок.}

Капитализм привлекает людей как желающих существ. У него есть либидозное измерение, где он обещает преодолеть отчужденность. Товарная форма определяет капиталистическое общество и указывает на приближающийся конец отчужденности. Глядя на товар, я вижу возможность достичь изобилия. Товар оторван от процесса труда и так выдает себя за чистый избыток, не запятнанный отчужденностью. Покупая, я преодолеваю отчужденность, ощущая этот избыток --- так уверяет капиталистическое общество. И это верно не только для покупателя, но и для капиталиста, который продал товар. Люди продают и покупают, чтобы приблизиться к совершенному избытку, уйти из определяющей субъектность отчужденности. Капитализм не создает отчужденность, но всеми силами убеждает нас, что ее можно избежать.

Вот простой товар --- автомобиль. Он создает ценность для продающих его капиталистов: они так увеличивают свой капитал, стремясь к пусть даже недостижимому, бесконечному накоплению. Автомобиль ценен и для покупателей: они видят в нем обещание бесконечного удовольствия, которое он может принести. Рекламодатели убеждают покупать все новые --- пусть и ненужные --- автомобили, обещая исцеление от отчужденности, которое имеющийся автомобиль уже не может дать. Покупка никогда не лишает отчужденности, но притягательная сила товара для капиталистического субъекта не слабеет --- столь огромна над людьми власть возможности избежать отчуждения. И производитель, и потребитель вкладывают свои средства в товар ради будущего вне отчужденности.

Поиск лекарства от отчужденности ежегодно порождает массу новых товаров. Большинство этих товаров не найдут широкого применения и перестанут производиться, но само их количество и разнообразие свидетельствуют о стремлении капитализма предложить таблетку от отчужденности. Никому не нужен Mountain Dew без коффеина и сахара, но само его существование --- как и любого другого товара --- подогревает мечту о неотчужденном существовании. В \textit{Очерке критики политической экономии} Маркс показывает, что капитализм должен создовать иллюзорные потребности: ''Производство поставляет не только материал для удовлетворения потребностей, но и саму потребность''\footnote{Karl Marx, Grundrisse, trans. Martin Nicolaus (New York: Penguin, 1993 [1857–1858]), 92.}. Эта потребность основана на фантазии, которую помогает создавать каждый из миллионов товаров. Капитализм заваливает субъектов десятками тысяч новых продуктов каждый год, чтобы исправно поддерживать иллюзию возможности избежать отчужденности. Без шквала новизны мы быстро осознаем свою отчужденность и порвем с фантазиями об избавлении от нее.

Акцент на отчужденности как конституирующей субъектность позволяет переосмыслить критику капиталистического общества. Капитализм не создает отчужденность, а мешает осознать ее неизбежность. Структура капитализма заставлят думать об отчужденности как о проблеме, более того --- решаемой проблеме, и потенциальное решение полагается существующем в товаре.

Капитализм окружает нас товарами и нам самим придает товарную форму. Он неустанно предлагает решение проблемы отчужденности в товаре. ''Проблема'' никуда не уходит, но каждая неудача все глубже погружает субъекта в логику капитализма: чем больше мне не удается достичь неотчужденности, тем больше я к ней стремлюсь.

Но есть еще один, более пагубный, эффект. Когда я выбираю один товар, а не другой, другой не пропадает: он остается и определяет мое отношение к выбранному, он все еще таит в себе потенциал излечить отчужденность --- потенциал гораздо больший, чем в выбранном товаре. Этот потенциал соответствует моему отношению к не-выбранному товару как к воображаемому объекту: если бы я его купил, магия рассеялась, так что он притягателен только оставаясь не-выбранным.

Оставленный товар сохраняет ту полноту, которую теряет купленный, и так я застреваю в постоянном стремлении к тому, что не выбрал. Я всегда выбираю неправильно. Я выбираю мятные леденцы \textit{Junior Mints} вместо драже \textit{Milk Duds}, но вкус леденцов напоминает мне, что я не ощущаю вкус драже. Замечая свою отчужденность, я думаю: ''Хм, эти мятные леденцы на самом деле не такие уж и вкусные'', и спешу переключить внимание на драже. Это психологическая ловушка, которую для меня создает товарная форма. Капиталистический субъект всегда упускает что-то лучшее, что-то \textit{все еще} обещающее вылечить отчужденность. При капитализме мы часто сталкиваемся с отчужденностью, но так же часто тут же думаем о товаре как о желанном избавлении от нее.

Товарная форма провоцирует думать о своих возможностях и забыть о своих ограничениях. Для товарной формы, единственная приемлемая граница --- субъектность другого, а не наша собственная. Это идея глубоко ошибочна. Рассматривая субъектность как внутренне неограниченную, цельную, мы интерпретируем каждый внешний барьер как посягательство на наши возможности. И потому кажется, что мы подвергаемся нападкам со всех сторон и никак не можем реализоваться в полной мере. Мы не признаем, что отчужденность фундаментальна для субъектности, капитализм убеждает нас, что она лишь временна и мимолетна.

Фантазия о неотчужденности ведет нас к повсеместному принятию системы массового неравенства, социальной изоляции и посягательств на свободу. Капитализм компенсирует это фантазией. Никто никогда не получал от капитализма обещанной им выгоды, но его притяжение не имеет себе равных.\footnote{Сериал \textit{Наследники} ясно показывает, насколько неудовлетворенными остаются те, кто выигрывает в капиталистической игре. Ни один здравомыслящий зритель не захочет быть на месте кого-либо из семьи Рой или людей из их социально-экономической когорты. Будучи победителями в игре капитализма, все участники сериала должны взять на себя обязательство преодолеть отчужденность, чего бы это ни стоило их отношениям.} Он соблазняет и казалось бы преуспевающих, и многих из тех, кто терпит неудачу.

В основе капитализма лежит фантазия о преодолении отчужденности через накопление. Хищничество капиталистов прямо следует из центральной роли этой фантазии и без нее рассеялось бы. Капитализм --- это практика, но бездумное накопление, необходимое для его поддержания, невозможно без перспективы исцеления раны отчужденности. Нет капиталистического мира, в котором отчужденность принимается как неизменный факт. В этом и заключено угнетающее воздействие капитализма. Единственный способ с ним бороться --- настаивать на истинности являемого в моментах отчужденности, оставаться разочарованным в \textit{Junior Mints} без надежды найти что-то другое в \textit{Milk Duds}.

\section{Гонка ко дну}

Расизм --- изобретение модерна. Конечно, ненависть между народами возникла задолго до него, до модерна субъекты идентифицировали себя как представителей разных народов и объединялись в разные группы, но не было понятия расовых различий. Изобретение расы сделало возможным рабство, колониализм, систематическое угнетение, повсеместное воровство и разграбление огромных территорий. Оно лежит в основе самых страшных преступлений современности.\footnote{Влияние расизма распространяется на все аспекты современности. Невозможно отделить, например, развитие капитализма от расизма, поскольку рабство сыграло решающую роль в становлении капитализма. Даже если капитализм в конечном счете обходится без рабства, ограничиваясь оплачиваемым трудом, необходимое для его развития накопление немыслимо без этого начального периода рабского труда.}

Расизм позволяет субъектам нарушать моральные принципы: все моральные философы современности осуждают покупку и продажу людей как материальных товаров, и все же, находясь под влиянием расизма, рабовладельцы и работорговцы не прекращают свою деятельность, продолжая считать себя хорошими, нравственными людьми. Они просто не верят, что имеют дело с субъектами, когда покупают и продают рабов. В расистской социальной структуре расовые различия значат отсутствие субъектности. Современность провоцирует видеть во всех других равных, но концепция расовых различий перебивает этот порыв. Отстаиваемые современным обществом, права человека распространяются только на принадлежащих этому обществу.\footnote{Несмотря на серьезные разногласия между Кантом и утилитаристами, обе стороны отвергают мораль рабовладельческого общества. Согласно Канту, рабовладение нарушает моральный закон, сводя субъекта к средству для достижения цели. Согласно же Джону Стюарту Миллю, рабовладение несовместимо со стремлением к максимально возможному счастью для всех.} Остальные --- не более чем символическая идентичность, не имеющая статуса в общественном порядке.

Европейский расизм неравномерно распределяет влияние модерна. Позволяет одним считать себя субъектами модерна, видя других в домодерновом статусе. В \textit{The Wretched of the Earth} Франц Фанон рассматривает очернение представителей другой расы как оборотную сторону модерна: ''И эта Европа, никогда не перестававшая говорить о человеке, никогда не перестававшая заявлять, что человек ее единственная забота --- эта же Европа теперь знает цену страданий, которую человечество заплатило за каждую из своих духовных побед''\footnote{Frantz Fanon, The Wretched of the Earth, trans. Richard Philcox (New York: Grove Press, 2004 [1961]), 236.}. Ценности модерна, отмечает Фанон, применимы не ко всем и не в равной степени. Расизм позволяет вовлеченным в него видеть в представителе другой расы самоидентичное, животное существо, которое можно терроризировать без угрызений совести.

Согласно расистской точке зрения, субъектностью не обладают и сами расисты --- и это привлекает их в расизме больше всего. Расизм избегает рассмотрения как самоотчужденного и нападающего расиста, и защищающегося. Есть лишь две несопоставимые расовые идентичности, противостоящие друг другу в устойчивой иерархии. Расисты знают свою символическую идентичность, свое Место --- потому что знают идентичность другой расы.\footnote{В своем \textit{Ориентализме} Эдвард Вади Саид объясняет, почему основной жест ориентализма (сродни любому расизму) заключается в лишении другого какой-либо расколотости: ''Восточный человек дан как фиксированный, стабильный, нуждающийся в исследовании, даже в знании о самом себе. И никакая диалектика ни желательна, ни допустима'' (Edward Said, Orientalism (New York: Vintage, 1978), 308). Ориентализм создает устойчивую оппозицию, в которой расовый другой существует в конкретном символическом положении. У другого не должно быть никакого неосознанного желания, ведь вспышка такого желания свидетельствовала бы о провале расистского проекта.} С расой человек обретает чувство безопасности, надежности своей символической идентичности. Расизм --- система распределения надежных позиций в символическом мире. Она позволяет своим участникам не замечать отчужденности себя и другого.

Расизм возвышает одну символическую идентичность за счет принижения другой. Вовлеченные в расистскую структуру обретают чувство собственного превосходства, но --- что гораздо важнее --- они обретают представление о том, кто они есть. Символически идентифицируя представителей другой расы, они утверждаются в том, кто они, утверждаясь в том, кто не они. Без еврея у арийца не было бы идентичности. Без чернокожести белокожесть не имела бы никакой ценности. Под влиянием расизма все эти идентичности кажутся обоснованными и защищенными от травмы отчужденности.

Философская основа расизма --- отрицание собственной отчужденности и отчужденности других. Это попытка отвергнуть открытие модерна, найти убежище в расовой идентичности. Эта идентичность призвана заменить субъектность. На одной стороне она идеализируется, на другой --- принижается.\footnote{Как отмечает Ахилл Мбембе: ''Понятие расы позволило представить неевропейские человеческие группы попавшими в ловушку низшей формы бытия. Они были искаженным отражением идеального человека, отделены от него непреодолимой временной пропастью'' (Achille Mbembe, Critique of Black Reason, trans. Laurent Dubois (Durham: Duke University Press, 2017 [2013]), 17)} Это разделение между не столько субъектом и другим, сколько одной символической идентичности и другой --- каждая из них служит способом не замечать отчужденность. Легко поверить в преодоление отчужденности, веря в подлинную реальность этих идентичностей.

Первоочередная задача расистского проекта --- изобретение расы. Прежде чем создать угнетаемую расу, нужно создать несколько различных рас. Раса не имеет биологического фундамента и у людей нет причин мыслить в терминах расовых различий, но такое мышление предлагает им путь в обход вопроса об их субъектности. Для этого нужно воспринимать расы как естественно существующие, ведь если видеть их результатом расистского акта, этот акт не сработает. Он должен быть незаметен. Именно потому Барбара и Карен Филдз работают над разоблачением расистского акта, скрытого в каждом упоминании расы.\footnote{Сегодня стало модным говорить о расах, дистанцируясь от существования биологических рас. Каждый признает, что раса --- культурная конструкция, а не биологический факт. Сестры Филдс справедливо отвергают такую позицию: ''Раса как культура --- всего лишь вежливое название для биологической расы, потому что невозможно всерьез постулировать культурную однородность среди тех, чьей расовой однородностью занимались ученые-биологи'' (Karen E. Fields and Barbara J. Fields, Racecraft: The Soul of Inequality in American Life (London: Verso, 2012), 156). Ссылаться на расу как на культурный факт --- значит участвовать в мистификации расизма в той же мере, что и ссылаясь на нее как на факт биологический.}

В \textit{Расовстве} сестры Филдс показывают, что важнейшим шаг для расистского проекта --- скрыть себя под маской расы. Ведь если раса действительно существует, то расизм --- ее неизбежное проявление. Расистский проект прикрывает расизм расовыми различиями, сводит возникшее в ответ на отчужденность действие к вопросу идентичности. Как пишут сестры Филдс: ''Расово-расистский маневр превращает расистские действия (оскорбительное обращение) в расовые (врожденные различия)''\footnote{Fields and Fields, Racecraft, 149.}. Расизм --- акт группы субъектов, социальный заказ, но акт \textit{замаскированный} под реакцию на якобы существовавшую до него идентичность. Так получается, будто расовая идентичность провоцирует действия субъекта-расиста, хотя на самом деле причинно-следственная связь работает с точностью до наоборот: расовой идентичности нет до акта расизма. За каждым формированием расовой идентичности стоит субъект-расист и его акт.

Борьба с расизмом, значит, требует перевода символической идентичности обратно в действия [отчужденных] субъектов. Там, где нам показывают расу, мы должны видеть создавший это расовое различие расистский акт. Раз расизм предлагает субъектам способ избежать принятия отчужденности, то вызов ему заключается в непосредственном ее принятии.

Об этом и пишет Пол Гилрой в \textit{Против Расы}. Видя, какие ужасы расизма в модерне, он утверждает, что сам модерн принес только освобождение: ''Привлекательность авторитарного иррационализма стала частью того, что значит быть человеком модерна. Это связано с мечтами о просвещении и автономии как вечной альтернативе''\footnote{Paul Gilroy, Against Race: Imagining Political Culture Beyond the Color Line (Cambridge: Harvard University Press, 2000), 237.}. Борьба с расизмом модерна (это, впрочем, единственная форма расизма) предполагает универсализацию модерна. Покуда есть хотя бы одна группа субъектов вне модерна, расизм будет работать.

Исключение некоторых субъектов из модерна обрекает на антимодерновое существование не только их, но и тех, кто их так исключает в попытке предаться варварскому обращению с представителями другой расы. Расизм поощряет анти-модерн по обе стороны баррикад. Без расизма модерн становится радикальным, универсальным и освобождающим.

\section{Место женщины}

Сексизм, подобно расизму, пытается облечь свои цели и участников в конкретную символическую идентичность. В сексизме никто не отчужден, ведь противопоставление мужчины и женщины дает две устойчивые идентичности. Обе стороны занимают конкретное место, но место женщин обыкновенно ограничивается домашним хозяйством.\footnote{Представление о работающей матери показывает степень ассоциации женщины с домом: внешняя занятость остается лишь дополнением. Остающийся дома для заботы о детях, мужчина, напротив, воспринимается --- в том числе самим собой --- менее мужественным, как если бы был не на своем месте.}

Согласно сексистской идеологии, мужчина выходит в мир и взаимодействует с другими, тогда как женщина занимается домашними делами. Мужчина узнает себя в участии в номосе, женщина же остается неотчужденной в оикосе: ''Место женщины --- дома''. Сексизм не приводит к отчуждению женщин, а относится к ним как к целостным существам, неподверженным отчужденности.\footnote{Один из главных художественных проектов Тони Моррисон посвящен тому, как патриарханльные фигуры препятствуют столкновению женщин с их отчужденностью. Такая работа ограничивает возможности женщин к восстанию. Роман \textit{Рай} начинается с нападения сексистов на группу женщин, отказывающихся оставаться на своем Месте. Эти пять женщин живут в монастыре, предоставляя убежище местным женщинам, столкнувшимся с домашним насилием, проводят аборты и всячески помогают в борьбе с патриарханльным контролем. Мужчины из городка Руби убивают их за осознание своей отчужденности и подталкивание к этому осознанию других женщин. Выстраивая роман вокруг этого акта насилия, Моррисон показывает главный страх сексистов --- отчужденность женщин.}

Феминистский анализ сексизма показывает роль отчужденности в угнетении. Пример угнетения можно увидеть в романе Симоны де Бовуар ''Второй пол''. Бовуар описывает угнетение женщин на протяжении всей аграрной истории, низведение женского начала до уровня самоидентичности. Женщина не способна пережить отчуждение от самой себя, усвоить универсальное. Бовуар рассказывает о неспособности женщин в сексистском обществе осознать отчужденность от своей конкретности: ''женщина... утопает в имманентности; но лишь потому, что была заточена в ней''\footnote{Simone de Beauvoir, The Second Sex, trans. Constance Borde and Sheila Malovany-Chevallier (New York: Vintage, 2011 [1949]), 643.}. Женщина становится лишь тем, кто она есть, неспособной осознать отчужденность от своей идентичности. Бовуар описывает женственность как сведение женской субъектности к статусу конкретной идентичности. Она критикует сексистское общество за отказ женщинам в принятии самораскола: сексизм мешает женщинам увидеть себя отчужденными субъектами.\footnote{Бовуар не отождествляет сексуальное угнетение с биологической детерминацией. Она настаивает, что женщиной не рождаются, а, скорее, становятся, добавляя, что никакое биологическое, психическое или экономическое предназначение не определяет роль женщины в обществе (Beauvoir, The Second Sex, 283).}

Бовуар рассказывает, как сексизм на протяжении всей своей долгой истории отождествлял женщину с ее телом. Быть женщиной --- значит быть намертво привязанной к телесной идентичности. Прежде семейной жизни женщина отождествляет себя со своим телом. Такая идентификация, согласно Бовуар, не позволяет ей заниматься спортом, тяжелой работой и интеллектуальной деятельностью\footnote{Как интеллектуальная деятельность связана с телесностью, остается загадкой (пр. переводчика).}. Будучи заключена в своем теле, женщина не может стать священником или политическим лидером. Процессы женского организма --- менструация, беременность, менопауза --- становятся символами, от которых женщина не может избавиться, и такая телесность --- и, что важнее, воплощенность --- часто используется как оправдание для принижения женского начала.

Впрочем, восхваление женского начала ничуть не улучшило бы ситуацию, ведь может быть столь же угнетающе, как и принижение. Бетти Фридан в книге \textit{Загадка женственности} уверяет, что прославление женского начала эффективно только если считать женственность неотъемлемой частью личности. В этой книге она пытается напомнить женщинам об их отчужденности, не поощряя их преодолевать ее. Освобождение возможно только через отчужденность.

Согласно Фридан, указание на ценность женственности становится основным способом угнетения женщин: ''высшей ценностью и единственным обязательством для женщин утверждается реализация их собственной женственности, и что якобы величайшей ошибкой западной культуры на протяжении большей части ее истории была недооценка этой женственности''\footnote{Betty Friedan, The Feminine Mystique (New York: Norton, 2013 [1963]), 35–36.}. Фридан реагирует на попытку возвести женщин на пьедестал. Сексизм направлен на убеждение женщин в первоочередной ценности их идентичности как женщин: кем они являются как женщины --- вот что выводится на первый план. Женщины должны принять Место в обществе, а не отчужденность.\footnote{Большая часть \textit{Загадки женственности} посвящена обвинению Фрейда в продвижении этой загадки для уверения женщин, что их проблемы имеют скорее психические, нежели социальные причины: ''Загадка женственности черпает свою силу из идей Фрейда, это он стремился заставить женщин и тех, кто их изучал, неправильно истолковывать разочарования своих матерей, обиды и неадекватность своих отцов, братьев и мужей, а также собственные эмоции и доступные выборы в жизни. Эта фрейдистская идея, будучи подана как очевидный факт, сегодня завладела умами почти всех американских женщин'' (Friedan, Feminine Mystique, 110). Конечно, психоаналитическая практика, особенно в США, склонна держаться вводящего в заблуждение типа анализа, который и критикует Фридан, но Фрейд так же выделялся своей критикой маскулинности как самозванства. Фаллос он всегда считал чем-то мошенническим. Подробней о таком понимании психоанализа см. Alenka Zupančič, The Odd One In: On Comedy (Cambridge: MIT Press, 2008).}

Опасение, что женщины обнаружат свою отчужденность --- главная тема крайне недооцененного фильма Оливии Уайлд \textit{Не беспокойся, дорогая} (2022). Этот феминистский шедевр показывает, как далеко готовы зайти мужчины ради убеждения женщин в отсутствии у них отчужденности, и как легко женщины поддаются соблазну отождествления с символической идентичностью. Мужчины даже согласны мириться с ужасной экономической ситуацией, пока способны поддерживать иллюзию безопасной идентификации у объектов своего вожделения.

\textit{Не беспокойся, дорогая} изображает замкнутый мир в городке Виктори, штат Калифорния, где женщины заботятся о хозяйстве, а их мужья каждый день работают в штаб-квартире ''Виктори''. Идиллический городок --- это сексистская утопия, в которой женщины остаются на своем Месте пока их мужчины работают ради поддержания этого Места. Но оказывается, что эта фантазия на самом деле лишь симуляция, созданная Фрэнком, убедившим остальных мужчин профинансировать ее.

По мере того, как главная героиня Элис Чемберс постепенно приходит в себя и понимает, что живет в симуляции, зрители также понимают, что мужчины финансируют участие себя и своих супруг за счет работы в реальном мире. Каждое утро, выходя на работу, они фактически покидают симуляцию и отправляются работать за ее пределами. Элис --- отнюдь не скучающая домохозяйка в изолированном городке, а врач отделения неотложной помощи. Из воспоминаний об этом периоде жизни Элис становится ясно, что она страдает от отчуждения, вызванного этой работой. Она устает от долгих часов изнурительного труда, но в то же время испытывает удовлетворение, которого ей так недостает в симуляции.

Фильм противопоставляет уют символической идентичности Элис в ''Виктори'' и отчужденность ее жизни врача за его пределами. Симуляция предлагает образ целостности, но отчужденное существование в реальной жизни кажется лучше --- и не только потому, что оно настоящее, но и потому, что само страдание от отчужденности приносит удовлетворение, в отличие от укорененности в Месте.

Стабильная идентичность привлекательна, но угнетение вызывает отторжение, потому что и отчужденность не лишена привлекательности. Отчужденность дает независимость от определения внешним. Она освобождает от конкретной символической идентичности, и тогда можно действовать вопреки соответствующим ей ожиданиям. Борьба против угнетения --- борьба за принятие отчужденности от репрессивного навязывания идентичности. Стремиться принять свою отчужденную субъектность --- значит выступать против всякого угнетения.

\chapter{Мечта о преодолении}

\section{Капитальное наказание}

Весь тот негатив в литературе, с которым сталкивается отчужденность, объясняется не только отсутствием хорошего агента или публициста --- он напрямую связан с успехом Карла Маркса, его историческим триумфом над Гегелем --- главным философом проекта эмансипации. Даже когда марксистские режимы рухнули по всему миру, критики капиталистического общества не отступили, продолжая клеймить капиталистическое общество отчуждающим, и повторять эту мантру в моменты кризиса капиталистической экономики.\footnote{Большинство современных марксистов-доктринеров вообще никогда не упоминают об отчуждении. Дэвид Харви, например, рассматривает проблему самовоспроизводства капитализма в терминах не отчуждения, а экономического роста. В \textit{Загадке капитала и кризисе капитализма} он напишет: ''Какие бы инновации или сдвиги не происходили, выживание капитализма в долгосрочной перспективе зависит от способности достичь совокупного роста 3\%'' (David Harvey, The Enigma of Capital and the Crisis of Capitalism (Oxford: Oxford University Press, 2010), 130). Согласно марксистам вроде Харви, отчуждение не имеет значения, когда речь заходит о борьбе с капитализмом --- фундаментальными являются экономический рост, долг и финансы.} Вопреки Марксу, Гегель утверждал, что отчужденность непреодолима и, значит, универсальна. Капитализм хоть и отчуждает, но не имеет монополии на отчуждение. Но никто не обращается к гегелевской философии всеобщей отчужденности как к альтернативе Марксу, хотя он и предлагает способ осмысления освободительного политического проекта, признающего несводимость субъекта. Следует пересмотреть взгляды Гегеля и Маркса на отчужденность.

С точки зрения Маркса, гегелевскую систему следует поставить обратно с головы на ноги: он ошибается, считая, что миром движут идеи, а не материальные условия. Да, противоречия --- двигатель истории, но противоречия материальные, не идеальные. Гегелевская диалектика движется в совершенно неверном направлении --- от идеи к материи. Маркс стремится внести коррективы, подчеркнуть приоритет материальных отношений над идеями.

В постфактуме ко второму изданию ''Капитала'', опубликованному в 1873г, Маркс представит вышеизложенный аргумент. Критики восхищены: философию Гегеля пора вернуть на землю, раскрыть экономическую подоплеку, которая на самом деле управляет нашим существованием.\footnote{Пожалуй, ни один последователь Маркса не воспринимает критику гегелевского идеализма острее, чем Теодор Адорно --- несмотря на его симпатию к философии Гегеля. В цикле лекций ''История и свобода'' Адорно утверждает, что Гегель ''хочет охватить вообще все --- даже противоречивое. Он придерживается точки зрения универсума, склонен идеологически и в духе конформизма утверждать, что универсум прав'' (Theodor W. Adorno, History and Freedom, trans. Rodney Livingstone (Malden, MA: Polity, 2006 [1964–1965]), 65). Занимая сторону универсального, Адорно отказывается от примата материального, говоря о нем как о лишь остатке.} Маркс предложит сочетать диалектический метод с эмпирической реальностью: ''Мистификация, которой подвергается диалектика в руках Гегеля, никоим образом не мешает ему быть первым всесторонне и сознательно представившим ее общие формы движения. У него диалектика, но вверх ногами. Ее нужно перевернуть, чтобы отделить рациональное зерно от мистических плевел''\footnote{Karl Marx, “Postface to the Second Edition,” in Capital: A Critique of Political Economy, Volume I, trans. Ben Fowkes (New York: Penguin, 1976 [1867]), 103. Как выразились Карл Маркс и Фридрих Энгельс в \textit{Немецкой идеологии}: ''Не сознание определяет бытие, но бытие определяет сознание'' (Karl Marx and Friedrich Engels, The German Ideology (Moscow: Progress Publishers, 1976 [1845–1846]), 42). Это и есть переворот Гегеля обратно с головы на ноги.}. Маркс возьмет эту инверсию за основу своей трансформации гегелевской системы. Маркс уверен, что такая радикальная коррекция ориентирует философскую систему Гегеля на политическое действие.\footnote{Ряд марксистов противопоставляют философию Гегеля и политическую программу Маркса. Среди них как негативно (Этьен Балибар), так и позитивно (Франк Руда) настроенные к Гегелю.} Так философия, прямо заявлявшая, что не может изменить мир, политизируется в программу, единственной целью которой будет такое изменение.\footnote{В своей единственной опубликованной работе, посвященной политике, Гегель в предисловии прямо отказывается предлагать теорию должной организации общества. В \textit{Философии права} он пишет: ''Этот трактат ... в той мере, в какой он касается политической науки, должен быть лишь попыткой понять и изобразить государство как по сути рациональную сущность. Как философское произведение, оно должно максимально дистанцироваться от обязательства построить государство определенным должным образом --- его инструкции направлены не на указание государству, каким должно быть, но на демонстрацию, как следует воспринимать его как этическую вселенную'' (G. W. F. Hegel, Elements of the Philosophy of Right, trans. H. B. Nisbet, ed. Allen W. Wood (Cambridge: Cambridge University Press, 1991 [1821]), 21). Гегель явно отвергает какое-либо политическое значение философии, но придает большое политическое значение тому, как мы распознаем государство --- что, он считает, является задачей философии.}

Но Маркс упускает, как отходит от Гегеля в куда более важном вопросе --- об отчуждении. Маркс напишет, что капитализм отчуждает работников от результатов их труда и что коммунизм --- лекарство от этого отчуждения. Сторонники марксизма видят в коммунизме спасение от отчужденности в капитализме. Другие левые движения --- не связанные с Марксом явно --- исходят из его критики отчуждения и строят свои программы как направленные в первую очередь против него. Понимание отчужденности как политически решаемой проблемы --- новшество Маркса, знаменующее отход от Гегеля.

Гегель отводит отчужденности положительную роль в развитии субъектности, и Маркс впоследствии напишет об этом. Гегель --- философ отчужденности, он не считает отчужденность проблемой или несчастьем, а видит в ней основу свободы и равенства. Согласно Гегелю, без отчужденности нет никакой субъектности. В \textit{Феноменологии духа} он напишет: ''Самосознание --- всего лишь нечто, и реально лишь постольку, поскольку отчуждается от себя''\footnote{G. W. F. Hegel, Phenomenology of Spirit, trans. Terry Pinkard (Cambridge: Cambridge University Press, 2018 [1807]), 284.}. Отчуждение субъекта от самого себя делает возможным опосредованное самоотношение, определяющее субъектность. И какие бы термины он для этого ни использовал --- Entfremdung или Entäußerung --- Гегель никогда не рассматривал отчужденность как нечто плохое.

Отчуждение, по Гегелю, избавляет нас от застоя самоидентификации: субъект так входит в то, что от него отлично и становится собой, превращаясь в то, чем не является. Отчуждение вырывает нас из естественного бытия и порождает субъектность.\footnote{В \textit{Феноменологии духа} Гегель восхваляет отчуждение как силу, создающую субъектность через разрушение нашего естественного бытия: ''Подлинная природа индивида и его сущность --- это дух отчуждения от естественного бытия'' (Hegel, Phenomenology of Spirit, 285.)} Субъект Гегель рассматривает сущность, чувствующую себя как дома в абсолютно чуждом ей --- и это может произойти только через отчуждение. В диалектике, считает Гегель, нет передышки от отчуждения, потому что оно в ней тотально.\footnote{Гегель настаивает на тотальности и целом, чтобы указать на неизбежность отчужденности. Если остановиться на полпути к тотальности, может показаться, что еще есть пространство, свободное от отчуждения. Но тотальность устраняет всякую такую возможность, это радикальный шаг.}

Диалектический метод Гегеля основан на распознавании чуждого в претендующем на тождество с самим собой. Например, ''Венера --- это вечерняя звезда'' --- скорее утверждение тождества, чем предикация, ведь вечерняя звезда идентична Венере, а не ее качество. Будучи рассмотренно так, недиалектически, утверждение указывает на идентичность. Но, согласно Гегелю, даже такие прямолинейные утверждения не избегают выявления чуждого внутри тождественного: отождествление Венеры и вечерней звезды указывает, что Венера не самоидентична. Ее потому и можно отождествить с самой собой --- вечерней звездой --- что она проницаема для не-себя. Идентичность проявляется в отношении, а не в изоляции от всего отличного. И потому любая идентичность сопряжена с противоречием.

Диалектика --- исследование противоречий, присущих каждой идентичности. Все, что кажется идентичным, в конечном счете --- через диалектическую интерпретацию --- оказывается разделенным на части и включенным в отличное от себя. Для Гегеля противоречие онтологично, и отчуждение --- способ его проявления в субъекте. Субъект достигает свободы через отчуждение, потому что только оно порождает и поддерживает субъектность. Гегель стремится к диалектической интерпретации субъекта и выявлению освобождающей силы отчуждения.\footnote{\textit{Феноменология духа} рассматривает отчужденность субъекта, тогда как \textit{Наука логики} --- противоречие понятия. В \textit{Феноменологии духа} размышления о противоречии просто ограничиваются отчужденностью субъекта.}

Таким образом, наиболее радикальный разрыв Маркса с Гегелем --- в его критике отчужденности и понимании ее как проблемы, которую следует преодолеть. Маркс здесь встает на безнадежный путь, потому что проект преодоления отчуждения не может быть завершен, это бесконечная --- и проигрышная --- борьба. Такое понимание отчуждения приводит к неверному пониманию политической борьбы. Отчужденность --- не ловушка, из которой можно убежать, а дистанцирование от самого себя, обеспечивающее свободу.

В ранних работах Маркса отчуждение сыграет важную роль. Отчуждение, считает он --- это тот самый ущерб, который нам наносит капиталистическое общество.\footnote{Как и Гегель, Маркс использует два разных термина, обозначая отчуждение: Entäußerung и Entfernung. Второе обычно переводят как estrangement (отчуждение, отстранение).} Согласно Марксу, экономическая система создает отчужденное общество. Основной формой отчуждения в капитализме становится отчуждение рабочих от продуктов своего труда.\footnote{Это главная тема \textit{Экономико-философских рукописей 1844г}, которые Маркс напишет в Париже. Там он будет утверждать: ''объект, произведенный трудом --- продукт труда --- выставляется как нечто чуждое, независимое от того, кто его произвел'' (Karl Marx, The Economic and Philosophic Manuscripts of 1844, trans. Martin Milligan (New York: International Publishers, 1964 [1844]), 108).} Рабочий производит товар, а капиталист владеет этим товаром и продает его с прибылью, создавая непреодолимую пропасть между рабочим и товаром. Рабочий не может делать с продуктом все, что пожелает --- он должен воспринимать его как нечто чуждое, что еще должен купить. Капиталист владеет тем, что создают рабочие --- такое отчуждение отнюдь не благотворно и потому далеко от гегелевской концепции отчуждения.

В феодальном обществе ремесленники производили товары в целостном процессе и непосредственно получали выгоду. В условиях капиталистических производственных отношений рабочий класс теряет связь с продуктами своего труда. Разделение труда устраняет ''ноу-хау'' феодального ремесленника, а прибыль с продажи товара получает капиталист. Товары ремесленника становятся товарами капиталиста. Товар --- это название отчужденного продукта труда.

В своих ранних трудах Маркс утверждает, что человечество воплощает себя в производственном труде. И потому при капитализме мы отчуждаемся от своей природы и от друг друга из-за непрекращающейся конкуренции. Капитализм создает коллектив рабочего класса, но не позволяет ему осознать себя коллективом: каждый рабочий ощущает себя изолированно, уединенно, отчужденно от остального общества. В \textit{Экономико-философских рукописях 1844} Маркс изложит это так: ''Отнимая у человека объект его производства ... отчужденный труд лишает его биологической жизни, его реальной объективности как представителя вида''\footnote{Marx, Economic and Philosophic Manuscripts, 114.}. Так он указывает, что объективация через труд --- способ самопознания человечества, невозможный при капитализме. Согласно Марксу, потеря так называемой биологической жизни при капитализме --- это потеря самого существенного в человечестве.\footnote{Бертелл Оллман связывает концепцию отчуждения Маркса с коммунистическим будущим без отчуждения. Это будущее --- мерило, которым можно определить наш современный статус отчужденного: ''из принятия коммунизма за эталон следует, что класс отчужден настолько, насколько не соответствует коммунистическому идеалу'' (Bertell Ollman, Alienation: Marx’s Conception of Man in Capitalist Society, 2nd ed. (Cambridge: Cambridge University Press, 1976 [1971]), 132). И хотя Оллман избегает говорить о нарушаемой капиталистическим отчуждением, изначальной сущности, он видит в проекте эмансипации преодоление отчуждения. Такое отношение к отчуждению играет ключевую роль в большинстве эмансипационных проектов после Маркса.}

Критика капиталистического отчуждения нашла признание не только среди марксистов, но и среди анархистов, коммунистов и даже некоторых консерваторов. Капитализм также создает отчужденную социальную жизнь, где люди общаются друг с другом посредством товара, а не напрямую. Те, кто призывает положить конец отчуждению, стремятся построить общество, в котором дистанция между субъектностью и символической идентичностью преодолена. В таком обществе каждый будет тем, кто он есть --- не будет внутреннего разрыва между субъектом и другими. Согласно Гегелю, такой общественный порядок кошмарен и невозможен.

\section{Незнакомый Гегель}

Гегель далек от мечтаний об освобождении от отчуждения, напротив --- он понимает свободу как возможную только через отчуждение. Гегель --- первый мыслитель, явно включивший отчуждение в свою философию, и пишущий о нем исключительно в позитивном ключе.\footnote{Декарт первый сформулировал понятие отчужденной субъектности --- пусть и не использовал эти слова. Кант внес термин \textit{субъект}, Гегель --- \textit{отчуждение}.} По Гегелю, отчуждение --- позитивный процесс, потому что уводит нас от конкретной ситуации и заставляет встретиться с другим. Отчуждение делает нас теми, кто мы есть, потому что через него мы перестаем быть собой. Оно позволяет осознать противоположное нам --- составной частью нас.

Будучи сосредоточены на непосредственно данном, мы упускаем то, что создает эту непосредственность. Мы уверены в конкретности, однозначности себя, потому что не видим то, что нас определяет. Явленность отчужденности рассеивает эту уверенность в самоидентичности. Согласно Гегелю, субъектность обнаруживает себя через ''насилие, совершаемое ею над самой же самой, [которое] разрушает ее собственное ограниченное удовлетворение'' \footnote{Hegel, Phenomenology of Spirit, 54.}. По Гегелю, мы обнаруживаем свою субъектность не в том, кем изначально себя считаем, а в том, чем мы в конечном счете становимся после такого насилия над собой. Дистанцирование от самих себя показывает, кто мы есть. Отчуждение от внешнего мира освобождает от тирании данности.

В системе Гегеля отчуждение занимает два этапа. Во-первых, дети становятся отчужденными из-за образования: Гегель отвергает романтическое представление об образовании как развивающем внутренний потенциал ребенка,\footnote{Похоже, что говоря об образовании как процессе отчуждения, которому ребенок вынужден подчиниться, Гегель выступает против \textit{Эмиль, или о воспитании} Жан-Жака Руссо --- его фундаментальной работы по теории образования, где оно рассматривается как развитие естественных склонностей ребенка.} говоря о нем как о насилии, разрушающем присущие ребенку склонности, не дающем им развиваться в соответствии с их собственной логикой. Это потому, что собственная логика никогда не логика самого субъекта --- скорее она логика того, что его незримо определяет. Насилие образования --- первое, что освобождает ребенка от его семейного и социального положения. Образование --- это освобождающее насилие.\footnote{Излишне говорить, что Гегель счел бы домашнее обучение катастрофой для образования, способом не дать детям обрести свободу. Домашнее обучение --- это образование без отчуждения, то есть не образование вовсе.}

Гегель подчеркивает конфронтацию с чуждым мышлением, которая определяет образование.\footnote{В \textit{Лекциях по философии Духа} Гегель напишет: ''В детей привносится нечто чуждое, и они, как человеческие существа, предчувствуют, что такой разрыв необходим. Чуждое и непонятное --- первый урок, чтение и письмо, абстрактные звуки и знаки --- элементарные уроки. Эта чисто внешняя целостность и есть то, что должно их занимать'' (G. W. F. Hegel, Lectures on the Philosophy of Spirit, 1827-8, trans. Robert R. Williams (Oxford: Oxford University Press, 2007), 97–98). Гегель говорит о необходимости отчуждающего образовательного процесса, потому что он станет основой для свободы.} Образование вовлекает детей в чуждый процесс: они заучивают наизусть и выполняют механические упражнения, чтобы научиться мыслить самостоятельно. Они не могут сразу перейти к независимому мышлению, только через отчуждающий механический процесс первых лет обучения. Свобода, как не странно, возникает из этого навязанного чуждого. Способность мыслить независимо зависит от рутинного развития навыков.

Но после такого отчуждения через образование субъекту необходим дополнительный этап, считает Гегель: субъект теперь отчуждается от себя, от своей объективной ситуации. Согласно Гегелю, никто не самоидентичен и каждому следует относиться к своей ситуации так, будто он ей чужд и в нее вторгся.

Особенно хорошо это показано в самой известной работе Гегеля --- \textit{Диалектике слуги и господина}. В ходе их диалектического взаимодействия господин осознает, что господство иллюзорно из-за полной зависимости от слуги, тогда как слуга осознает свою свободу, видя как акт подчинения приводит к автономии. Каждая из позиций оказывается отличной от себя же из-за фундаментальной отчужденности: господин оказывается слугой слуги, а слуга --- господином господина.

В результате отчуждения субъект обретает свободу, но свобода эта зависит от ограничений объективной ситуации. Гегель рассматривает концепцию судьбы как способ выражения ограничений, с которыми сталкивается свободный субъект. Без отчужденности, отделяющей его от ситуации и так дающей свободу, нельзя было бы говорить о судьбе. В \textit{Науке логики} Гегель утверждает, что отчужденность --- предварительное условие для переживания судьбы: ''Только самосознание имеет судьбу в строгом смысле слова, ведь оно свободно и потому в единичности своего ''я'' абсолютно существует само по себе и для себя и может противопоставлять себя своей обективной универсальности и отчуждать себя от нее. Но этим разделением оно противопоставляет себя механической судьбе''\footnote{G. W. F. Hegel, The Science of Logic, trans. George Di Giovanni (Cambridge: Cambridge University Press, 2010 [1812]), 639.}. Хотя судьба и значит несвободу, Гегель связывает ее со свободой субъекта: человек может ощутить скованность судьбой только после отстранения от своей ситуации и так обретения свободы. По Гегелю, мы отчуждаем себя в свободу после отчуждения в образовании.

Различия в позициях Гегеля и Маркса касательно отчуждения приводят к разным представления о том, каким может быть общество. Для Гегеля, никакое улучшение социального порядка не устраняет отчужденность субъекта --- она неотъемлимая часть субъектности. Маркс же надеется на решение ''проблемы'' отчужденности в обществе будущего.

Рассуждая о разногласиях Гегеля и Маркса об отчуждении, Шломо Авинери указывает на опасность, таящуюся в концепции Маркса. В ''Теории современного государства Гегеля'' он напишет: ''Гегель рассматривает отчуждение как необходимый аспект объективации, Маркс же утверждает, что оно присуще не самому процессу производства, а только его конкретным историческим условиям. Потому Маркс видит возможность окончательного спасения, тогда как Гегель ее отрицает''\footnote{Shlomo Avineri, Hegel’s Theory of the Modern State (Cambridge: Cambridge University Press, 1972), 94. Трудно переоценить эту работу Авинери.}. Гегель не отрицает, что можно улучшить условия жизни в обществе --- он отвергает образ спасения, вкладываемый Марксом в коммунистическую систему будущего. Если отчужденность необходима для свободы, то такое коммунистическое общество невозможно без отчуждения, и попытка построить его неизбежно приведет к катастрофе.\footnote{Герберт Маркузе хоть и пытается придать своему марксизму гегелевский оттенок, не решается зайти достаточно далеко: его политическая философия все еще опирается на возможность преодоления отчуждения, которое он связывает с капитализмом. В книге \textit{Разум и революция}, посвященной гегелевской политике, Маркузе напишет: ''Индивиды изолированы друг от друга и настроены друг против друга. Они связаны не лично, а посредством товаров, которыми обмениваются. Отчуждение человека от самого себя --- это его отчуждение и от своих собратьев'' (Herbert Marcuse, Reason and Revolution: Hegel and the Rise of Social Theory, 2nd ed. (New Jersey: Humanities Press, 1983 [1941]), 279). Борьба за устранение отчуждения друг от друга --- вот как Маркузе представляет политическую борьбу. Он не готов признать, что отчуждение и есть то, что способствует освобождению.}

Ключевое различие философий Гегеля и Маркса --- между структурирующей и исторической концепцией отчуждения. Гегель рассматривает отчуждение как структурирующее, определяющее для субъекта, чему соответствовала бы политика, усиливающая отчуждение.

\section{Альтюссер против Лукача}

Многие последователи Карла Маркса дистанцируются от его анализа капитализма как навязывающего рабочему классу отчуждение. Они критикуют раннего Маркса, его, по-видимому, остаточную инвестированность в гуманизм, унаследованный от Гегеля, от которого он позже --- когда будет писать \textit{Капитал} --- откажется. Главный из этих критиков --- Луи Альтюссер --- увидит эпистемологический прорыв на интеллектуальном пути Маркса: можно, утверждает Альтюссер, выделить молодого Маркса-гуманиста, сосредоточенного на отчуждении в капитализме, и зрелого Маркса-ученого, анализирующего структуру капитализма без какой-либо критики в адрес отчуждения. Этот поворот знаменует значительный прогресс критической мысли Маркса.

Альтюссер проводит это разделение неспроста: это ключ к построению марксизма --- не гуманистического, а могущего превзойти раннюю моральную критику капитализма Марксом и подняться до уровня науки. В \textit{Для Маркса} Альтюссер попытается провести четкое разграничение между молодым и зрелым Марксом. Он припишет концепцию отчуждения влиянию Фейербаха: ''Ранние работы Маркса пропитаны идеями Фейербаха. У Маркса 1842 и 1845 годов не только терминология фейербахианская (отчуждение, видовое бытие, тотальное бытие, ''инверсия'' субъекта и предиката и т.д.), но и --- что более важно --- основная философская проблематика''\footnote{Louis Althusser, For Marx, trans. Ben Brewster (London: Verso, 1969 [1965]), 45}. Для Альтюссера Маркс, когда говорит об отчуждении, остается в рамках философии, а к научному анализу переходит по мере отказа от такой терминологии и тем.

Альтюссер избегает концепции отчуждения, связывающей Маркса с, как он считает, наивной концепцией человека. Ранний Маркс, критикующий капиталистическое отчуждение, признавал в человеке нечто эссенциальное, некую человеческую сущность, которую развитие коммунизма и должно восстановить, постепенно преодолев отчуждение. Но Альтюссер видит в этом наивность и неустойчивость системы Маркса.\footnote{Альтюссер воспринимает раннего Маркса как наивного философа истории: ''Несостоятелен тезис, выдвинутый Марксом в \textit{Рукописях 1844г}, что история --- это история процесса отчуждения субъекта, родовой сущности человека, отчужденной в ''отчужденном труде'''' (Louis Althusser, “Lenin Before Hegel” (1969), in Lenin and Philosophy and Other Essays, trans. Ben Brewster (New York: Monthly Review Press, 1971), 121). Альтюссер утверждает, что после разворота Маркс станет смотреть на историю как на развитие специфических экономических систем, в которых субъектность играет лишь идеологическую роль.} Альтюссер отвергает идею сущности человека, от которой его можно отдалить: принять такую сущность --- значит примкнуть к тому самому капиталистическому мышлению, которому марксизм должен был противостоять.

Антигуманистический марксизм Альтюссера воплощается в марксистской науке --- науке, основанной на анализе, представленном в трех томах \textit{Капитала} и других поздних работах. Но отказываясь от проблемы отчуждения, Альтюссер устраняет всякую связь марксистского анализа капитализма с политической деятельностью пролетариата, ведь она требует идеологии и не может быть мотивирована истинами марксистской науки. Это противопоставляет Альтюссера марксисту-гегельянцу Георгу Лукачу, удваиваивающему вклад раннего Маркса в анализ проблемы отчуждения.

Развивая свою версию марксизма, направленную на преодоление отчуждения, Лукач сокращает разрыв между теорией и практикой так, как это не удавалось Альтюссеру: он формулирует наиболее философски разработанную марксистскую теорию отчуждения, придав мысли Маркса гегелевский оттенок. Но, как и Маркс, Лукач рассматривает отчуждение как дефицит, который обществу следует преодолеть. В \textit{Истории и классовом сознании} он опишет проблему капитализма в терминах отчуждения: ''Человек в капитализме сталкивается с реальностью, ''созданной'' им же самим (как классом), но кажущейся ему естественным, чуждым ему, явлением''\footnote{Georg Lukács, History and Class Consciousness: Studies in Marxist Dialectics, trans. Rodney Livingstone (Cambridge: MIT Press, 1971 [1922]), 135.}. Благодаря осознанию отчуждения посредством классового сознания пролетариат может оказаться в состоянии преодолеть его: пролетариат --- в отличие от капиталистов --- одновременно объект и субъект истории. Он отчужден, когда превращен в объект, но может преодолеть отчуждение, став субъектом истории, участвуя в революционной практике.

Различие между антигуманистическим марксизмом Альтюссера и гуманистическим марксизмом Лукача соответствует одной из основных линий борьбы левых политиков в XX веке. Но ни один из мыслителей --- даже сравнительно гегельянский Лукач --- не окажется способен признать за отчуждением освобождающий процесс. Марксистский анализ капитализма идет по ошибочному пути, когда видит отчуждение проблемой --- что остается неизменным даже после отказа Маркса от этого термина. Так что очевидная разница между Альтюссером и Лукачем оказывается не так уж существенна, ведь оба они не рассматривают отчуждение как структурирующее субъектность.

\section{Наблюдая за детективом}

Даже отказавшись от всякого упомянания отчуждения, зрелый Маркс сохранил концепцию отношения капитализма к труду, занимавшую в раннем его анализе центральное место. Он перестал говорить об отчуждении человеческой сущности при капитализме, но остался в рамках структурной логики отчуждения: капитализм должно преодолеть, потому что он отбирает у рабочего класса то, что тот сам производит. Маркс здесь критикует капитализм как отчуждающую структуру, пусть даже не используя термин \textit{отчуждение}.\footnote{Во многих кратких описаниях поворота мышления Маркса от отчуждения говорится о его повороте к овеществлению. Но этот термин (Verdinglichkeit) не встречается в \textit{Капитале}, это вклад Георга Лукача, а не что-то присущее самому Марксу.}

Это и упускает Альтюссер, говоря об эпистемологическом разрыве в развитии Маркса. Разрыв не так уж велик: конечно, после \textit{Экономико-философских рукописей 1844г} Маркс больше не упоминает отчуждение, но развиваемая им в \textit{Очерках 1850гг} и трех томах \textit{Капитала}, критика формально совпадает с его разработками из 1844г. Просто то, что ранее он называл отчуждением рабочего, он теперь называет присвоением капиталистом прибавочной стоимости, производимой рабочим. Это различие не столь существенно, как настаивает Альтюссер.

Маркс начнет первый том \textit{Капитала} с раскрытия тайны товара. Его разработка трудовой теории стоимости, заимствованная у политэкономистов Адама Смита и Дэвида Рикардо, приведет его к концепции прибавочной стоимости --- ключевому понятию во вкладе Маркса в трудовую теорию стоимости, выходящей за рамки того, что Смит или Рикардо могли себе представить. Преимущества этой концепции очевидны: с ней Маркс избегает лишь моральной критики капитализма, а также может количественно оценить капиталистическую эксплуатацию. Капиталисты эксплуатируют рабочих в той мере, в какой присваивают себе произведенную рабочими прибавочную стоимость. Капиталисты берут то, за что не платят, что подрывает основополагающее для оправдания капитализма понятие свободного рынка. В капитализме нет рынка равного обмена --- ни для рабочего, потому что он становится беднее от обмена, ни для капиталиста, потому что он извлекает из обмена выгоду. Благодаря анализу избыточного рабочего времени и производимой им прибавочной стоимости, обвинения Маркса в адрес капитализма приобретают куда больший вес.

Концепция прибавочной стоимости даст Марксу и его последователям конкретный способ определения идеологии, и хотя Маркс никогда не использует термин \textit{идеология} ни в \textit{Капитале}, ни в более поздних работах, различие, проведенное им между прибылью и прибавочной стоимостью, позволяет легко сформулировать теорию присущей капитализму идеологии. Капитализм функционирует идеологически, потому что преобразует прибавочную стоимость в прибыль, тем самым скрывая роль труда в создании стоимости.\footnote{С точки зрения ориентированного на прибыль капиталиста, указывает Маркс, роль труда в создании прибыли незаметна. Во втором томе \textit{Капитала} он пишет: ''Производственный процесс предстает как лишь неизбежный промежуточный термин, необходимое зло для зарабатывания денег'' (Karl Marx, Capital: A Critique of Political Economy, Volume Two, trans. David Fernbach (New York: Penguin, 1978 [1884]), 137). Капиталистическое производство никогда не сосредоточено на производстве как таковом.} Прибыль скрывает труд. Заменяя прибавочную стоимость прибылью, капитализм уверяет, что прибыль эта достигается трудолюбием или разумным инвестированием, а не накоплением отнимаемой у работников прибавочной стоимости.

Иначе говоря, капитализм требует мышление в терминах прибыли, а не прибавочной стоимости --- мышление, чуждое тому, как на самом деле работает накопление. Маркс посвящает большую часть третьего тома \textit{Капитала} исследованию взаимосвязи между прибавочной стоимостью и прибылью. Там он осуждает прибыль как замаскированную форму прибавочной стоимости --- форму, отчуждающую работников от созданной их трудом ценности: ''Прибыль, значит, то же, что и прибавочная стоимость, за исключением лишь мистифицированной, хоть и неизбежно возникающей в капитализме, формы''\footnote{Karl Marx, Capital: A Critique of Political Economy, Volume Three, trans. David Fernbach (New York: Penguin, 1981 [1894]), 127.}. С помощью \textit{Капитала} Маркс надеется противостоять этой мистификации.

Чуждая форма прибыли скрывает, что прибавочная стоимость принадлежит рабочему, который ее и произвел. Но при капитализме не только прибавочная стоимость представляется искаженно --- капитализм основан на присвоении ценности, созданной трудом сверх того, что требовался бы для поддержания жизни общества. Это присвоение --- отчуждение прибавочной стоимости у тех, кто ее создал.

Посредством механизма капиталистического производства и обмена, прибавочная стоимость трансформируется в прибыль. В \textit{Капитале} Маркс попытается проследить происхождение прибыли из прибавочной стоимости и связать ее с деятельностью рабочего класса, а не смелостью и изобретательностью капиталиста. Ориентация капитализма на прибыль от прибавочной стоимости позволяет капитализму удерживать всех, кто в нее инвестирует, несмотря на порождаемое им катастрофическое неравенство. Более того, сокрытие прибавочной стоимости говорит об отчуждении работников от продукта их труда. В форме прибыли их продукт предстает в чуждой форме, скрывающей реальный источник его ценности --- акт труда.

Содержание критики Маркса могло меняться, но ее целью всегда оставалась чуждая форма, которую социальные отношения принимают в капитализме. Он сделает ставку на разрушение этой чуждой формы и так создание общества без ограничений на развитие производственных сил. Капитализм тоже стремится к совершенному развитию производственных возможностей, но никогда не достигает его --- в отличие, уверен Маркс, от коммунизма. В \textit{Очерке критики политической экономии} Маркс рассуждает о ''новом способе производства'', выходящий за пределы капитализма, ''где свободное, беспрепятственное прогрессивное и всеобщее развитие производственных сил само уже предпоссылка общества, и, значит, его воспроизвеодства; где единственной предпосылкой становится продвижение за пределы отправной точки''\footnote{Karl Marx, Grundrisse, trans. Martin Nicolaus (New York: Penguin, 1993 [1857–1858]), 540.}. Производство может стать по-настоящему универсальным --- осуществляемым на благо всех и каждого --- в обществе будущего, которое создаст коммунистическая революция. Маркс уверен, что общество может преодолеть ограничения отчуждающей силы капитала, толкающей его в эксплуатацию.

Даже отказываясь от идеи отчуждения, Маркс остается верен раннему своему образу мышлению: общество искусственно снижает нашу производительность внешними ограничениями, виной которым его формальная структура. Свергнув эту структуру --- считает Маркс --- мы сбросим наконец оковы чуждой власти. Покончить с нашим капиталистическим отчуждением --- уверен он --- значит вернуть не утраченную человеческую сущность, но плоды нашего собственного труда.

\section{Проблемы Гегеля}

Маркс придерживается радикально иного подход к возникновению проблем, нежели Гегель, потому что и позиции касательно отчуждения у них радикально разные. Возникающие в истории и в мышлении проблемы он рассматривает как преодолимые трудности. Мы можем разрешить эти проблемы, анализируя их с научной точки зрения, вне морали и технократии.\footnote{Нападки Маркса на Прудона в \textit{Нищете философии} --- самая ранняя критика технократии. Прудон далек от современных либеральных технократов, но все же разделяет веру в технические решения проблем капитализма} После надлежащего научного анализа потребуется также революционная практика. Вместе с проблемами история всегда преподносит их решения, нужно только правильно интерпретировать происходящее и --- что немаловажно --- не довольствоваться лишь убедительностью интерпретации.

Оптимистичный взгляд Гегеля на отчуждение заставляет его ценить проблему ради нее самой: правильное отношение к эпистемологической или исторической проблеме --- рассмотрение ее как ее собственного решения. Гегель не стремится устранить проблему, а находит решение внутри нее самой. Конечно, для этого нужно занять особую позицию в отношении проблемы: не искать решение, а отказаться мыслить вне структуры самой проблемы. Если диалектически смотреть в проблему, проблема сама найдет решение. Это приводит Гегеля к созданию системы, пытающейся изменить мир через переосмысление проблем. Маркс хочет изменить мир, а не просто интерпретировать его, Гегель же делает ставку на правильную интерпретацию, рассматривающую проблему как выражение скрытого противоречия.

Вот как это работает в разделении между явлениями и вещами-в-себе в \textit{Критике чистого разума} Канта. Субъект может воспринимать только явления и не способен знать что-либо о вещах-в-себе (кроме того, что они существуют). Это ограничение парадоксальным образом освобождает субъекта от субъективизма: согласно Канту, в пределах области видимости мы можем быть уверены в нашем объективном знании --- что было бы невозможно, воспринимай мы вещи-в-себе своими объектами, ведь так скептицизм о нашей способности к познанию вещей-в-себе уничтожил бы всякую возможность объективного знания.

Но несмотря на успех Канта в построении системы, защищающей объективное знание от скептицизма (подобного тому, что продвигал Дэвид Юм), он, все же, сталкивается с трудностями касательно вещи-в-себе.\footnote{Согласно Дэвиду Юму, невозможен корректный переход от субъективной привычки к объективной связи: ''все наши рассуждения о причинах и следствиях основаны лишь на привычке'' (David Hume, A Treatise of Human Nature (Oxford: Oxford University Press, 2000), 123). Скептицизм Юма подтолкнул Канта к проведению фундаментального различия между явлениями и вещами-в-себе. Без этого различия Кант едва ли отличался от предшествующих мыслителей.} Сам Кант не осознавал еще эти трудности, но позже мыслители от Фихте до Шеллинга и Гегеля быстро ее заметили: субъект должен уметь отличать явления от вещей-в-себе, хотя вещь-в-себе должна оставаться непознаваемой.

Наиболее проницательные интерпретаторы Канта говорят о вещи в себе как нашем представлении о том, что нельзя помыслить --- а не как о реальном объекте. Как отмечает Дитер Генрих: ''Под вещью в себе Кант имеет в виду что нет условий, влияющих на сущность данности вещи как таковой. То есть \textit{вещь-в-себе} --- это ограничивающее понятие, обозначающее то, что мы не можем помыслить''\footnote{Dieter Henrich, Between Kant and Hegel: Lectures on German Idealism, ed. David S. Pacini (Cambridge: Harvard University Press, 2003), 49.} Генрих уверен, что Кант, говоря о вещи-в-себе, говорит не о реальном объекте, но даже столь щедрая интерпретация оставляет проблему концептуализации того, что нельзя помыслить. Акт формулирования лежащего за пределами опыта --- даже как чисто ментального объекта --- подразумевает некоторую способность мыслить за пределами мышления, что, конечно, невозможно.

Ответ Гегеля на трудность с вещью-в-себе станет парадигмальным для его философии. Он не пытается избежать противопоставления явлений и вещей в себе, а исследует, почему оно возникает: потому что вещи всегда выглядят так, будто скрывают неочевидную сущность. Кантовская концепция вещи-в-себе --- выражение противоречий, создающих разобщенность внутри каждой видимости.\footnote{Обращаясь к кантовским антиномиям --- неразрешимым противоречиям --- нужно обе противоречивые возможности либо принять, либо отвергнуть. Гегель видит в формировании антиномии выражение онтологического противоречия, и не обвиняет разум в антиномии (как это делал Кант), а рассматривает проблему как ее собственное решение --- очевидное, если перевести ее с эпистемических в онтологические термины. Для Канта онтологическое противоречие просто невозможно, Гегель же выводит его из возникновения неразрешимого эпистемологического противоречия. Бытие само по себе должно быть противоречивым, чтобы антиномии были возможны, уверен Гегель. Такое решение проблемы антиномий не отбрасывает их.}

Гегель не ищет решения, которое устраняло бы проблему --- он находит направление мысли, признающее решение, заложенное в самой проблеме. Гегелевское решение углубляет проблему. Идеология уверяет, что проблемы можно преодолеть, если проводить правильную политику. Гегель же признает, что хотя некоторые проблемы и можно преодолеть, проблема отчуждения не из таких. Из отчужденности нет выхода, нет выхода в ситуацию, где индивид примирился с социальным порядком, а субъект примирился с самим собой. Следует взглянуть на проблему иначе --- увидеть ее порождающей, а не ограничивающей. Признать, что отчужденность фундаментальна.

Гегелевское решение исторических и философских проблем заключается в принятии отчужденности как неотъемлимой от субъектности. Любое решение, направленное на преодоление отчуждения, всегда остается в стороне. Конечно, недостаточно интерпретировать проблемы под другим углом, но интерпретационный ход определяет эффективные практические вмешательства. Согласно Гегелю, политические проекты должны видеть в отчуждении способ освобождения и работать в рамках такой концепции.

Маркс видит ситуацию совсем иначе. Задача теоретика, считает он,--- указать на решение стоящих перед человечеством проблем, ведь все наши проблемы, включая смерть, имеют решение, которое станет очевидным по мере развития истории.\footnote{В отличие от Мартина Хайдеггера, Маркс говорит не об универсальном страхе смерти. Такой страх есть только у тех, уверен он, кто попал в ловушку принудительного индивидуализма капиталистического общества: ''Смерть представляется жестокой победой вида над определенным индивидом и противоречит их единству. Но отдельный индивид --- все лишь существо определенного вида и как таковой смертен'' (Marx, Economic and Philosophic Manuscripts, 138). Почти 150 лет спустя теоретик марксизма Фредрик Джеймсон разовьет эту идею в \textit{Политическом бессознательном}: ''В обществах будущего люди по-прежнему стареют и умирают, пари Паскаля марксизма в другом: смерть в раздробленном и индивидуализированном обществе куда страшней, чем в цельном, в котором умирание происходит и с группой, и с отдельными индивидами --- просто с разной интенсивностью'' (Fredric Jameson, The Political Unconscious: Narrative as a Socially Symbolic Act (London: Routledge, 2002 [1981]), 250). Впрочем, и Маркс, и Джеймсон, не поясняют, почему это ослабит --- а не усилит --- воздействие смерти. Их подход не до конца продуман: смертен не только отдельный человек, но и весь вид в целом и все что когда-либо будет жить, так что тревога о своей смерти выступает синекдохой тревоги о тепловой смерти Вселенной.} Будучи материалистом, Маркс не верит, что чисто теоретический акт может решить мировые проблемы: теория может только подсказать путь к практической реализации, подсвечивая противоречия господствующей социально-экономической системы, что Маркс и пытается сделать в своих трудах.

Маркс настаивает, что решение есть --- что наши отчужденные отношения можно преобразовать революционной практикой. Решение заключено в рабочем классе. В \textit{Экономико-философских рукописях 1844г} Маркс напишет: ''Коммунизм --- разгаданная загадка истории, и он знает, что является этим решением''\footnote{Marx, Economic and Philosophic Manuscripts, 135.} Маркс идет даже дальше: коммунизм не только разрешает загадку истории, но и устраняет отчуждение от мира природы: ''коммунизм ... --- подлинное разрешение конфликта между человеком и природой и между человеком и человеком''\footnote{Marx, Economic and Philosophic Manuscripts, 135.}. Маркс признает, что в коммунистическом обществе индивидуальные конфликты между людьми никуда не исчезнут, но делает ставку на то, что не станет антагонистических конфликтов в общественном строе и между человечеством и природой. Проблемы, кажущиеся сегодня фундаментальными, найдут свое решение в радикальных социальных изменениях.

Все наши проблемы разрешимы, уверен Маркс, потому что проистекают из материальных условий, и условия эти также содержат их решения. В предисловии к \textit{Вкладу в критику политической экономии} он пишет: ''таким образом, человечество неизбежно ставит перед собой только те задачи, которые оно в состоянии решить, ведь при ближайшем рассмотрении всегда будет видно, что проблема возникает лишь тогда, когда материальные условия для ее решения уже имеются или хотя бы находятся в стадии разработки''\footnote{ Karl Marx, A Contribution to the Critique of Political Economy, trans. S. W. Ryazanskaya (New York: International Publishers, 1970 [1859]), 21.}. Маркс здесь рассматривает теоретические проблемы как задачи политической практики.\footnote{Георг Лукач проследует с этой идеей до конца и обвинит в кантовских антиномиях чистого разума господство капиталистического способа производства. Решение антиномий, укажет Лукач, заключается на в переходе Гегеля от эпистемологии к онтологии, а в политической практике пролетариата. Пролетариат, будучи субъектом истории, разрешает эти кажущиеся неразрешимыми противоречия, ведь они не что иное как результат ограниченности буржуазного сознания. В \textit{Истории и классовом сознании} он утверждает: ''От своей неспособности понять историю, созерцательное отношение буржуазии разделилось на две крайности: с одной стороны --- ''великие личности'', которых рассматривали как деспотичных творцов истории, с другой --- ''естественные законы'' исторической среды'' (Lukács, History and Class Consciousness, 158). Увидев в пролетариате субъект истории, мы увидим разрешение антиномии между свободой и детерминизмом. Решение кроется в политической практике поднимающегося класса.} Теоретическая проблема возникает для выявления возможности могущего ее решить практического действия. Согласно Марксу, мы видим проблемами только то, что можем преодолеть. Маркс рассматривает проблемы с точки зрения их решений, тогда как Гегель ищет решения, идентичные самим проблемам.

Маркс в своем решении проблемы товара не пойдет по пути Гегеля, а выявит в товаре тайную силу: за каждым товаром скрывается кража, и Маркс выступает детективом, разоблачающим преступника. Каждый товар, утверждает он, содержит неоплаченный излишек труда рабочего класса. Такая интерпретация, уверен он, открывает путь для революции, могущей устранить несправедливость.

\section{Последняя дуэль}

Отношения между Гегелем и Марксом --- не только академический вопрос, это также вопрос восприятия и осуществления социальных изменений. И если идеалом оказывается искоренение отчужденности или решение загадки человеческой истории, то мы склонны считать, что такая цель оправдывает любые средства. Для Гегеля же цель неотделима от средств ее достижения: цель --- не что иное, как средства, которые мы выбираем для ее достижения. Мы всегда уже в пункте назначения.

Иначе говоря, Гегель --- диалектический мыслитель в отношении эмансипации, тогда как Маркс куда более дуалистичен. Гегель видит эмансипацию в отчужденности, Маркс же --- в ее преодолении. Согласно Марксу, наше настоящее в отчужденности в самой своей сути отлично от будущего вне отчужденности: царство свободы совершенно отдельно от царства необходимости.

Наиболее заметен отход Маркса от фундаментальной отчужденности к концу третьего тома \textit{Капитала}. В известном отрывке он проведет различие между царством необходимости и царством свободы, но по мере его описания этих двух сфер становитя ясно, что в его понимании они никак не пересекаются: сфера свободы не возникает диалектически из противоречий в сфере необходимости, а необходимость не возникает вновь от накопления противоречий в сфере свободы. Маркс связывает их дуалистически: сфера свободы начинается там, где заканчивается сфера необходимости --- никакого диалектического перехода из одной сферы в другую!

Маркс не использует термин \textit{отчуждение} в описании взаимосвязи между необходимостью и свободой, но очевидно считает необходимость подразумевающей отчужденность. Пока необходимость заставляет меня двигаться в определенном направлении, я самоотчужден или нахожусь под контролем чуждой силы. Маркс отказывается интегрировать отчужденность в сферу свободы, и потому вынужден строго разделить сферы, дабы необходимость не загрязняла сферу свободы своей чуждой силой. Иначе, уверен, Маркс, свобода не была бы свободой.

Для Маркса, настоящая свобода --- та, что свободна от силы необходимости. В пожалуй самом важном отрывке из третьего тома \textit{Капитала} он пишет: ''Всегда остается царство необходимости. Истинное царство свободы --- развитие человеческих способностей как самоцель --- начинается за его пределами, хотя и может процветать лишь когда в его основе лежит царство необходимости''\footnote{Marx, Capital, Volume Three, 958–959.}. Маркс прав: человечество никогда не решит проблему необходимости. Он не Рэй Курцвейл.\footnote{Рэй Курцвейл --- сторонник трансгуманизма --- уверен, что мы можем эволюционировать от нашей субъектности к коллективному, и, значит, неотчужденному существованию с помощью технологий.} Но он видит свободу сферой, в которой человечество рационально определяет себя и свое отношение к природе, и такая сфера не имеет ничего общего с необходимостью, ведь наличие необходимости указывало бы на отсутствие свободы.

Если бы мы могли создать отдельное от царства необходимости, царство свободы, мы бы наверняка пошли ради этого на все. Но если царство свободы всегда остается втянутым в царство необходимости, то способ достижения свободы не менее значим, чем сама свобода. Царство свободы --- не будущее, которого следует достичь, а способ борьбы в условиях необходимости. То, как мы боремся, имеет большее значение, чем то, чего мы достигаем.

Понимание отчужденности как средства к освобождению знаменует поворот от Маркса обратно к Гегелю. Благодаря ему можно думать об эмансипации не как о будущем, а как о другом способе отношения к настоящему. Мы так можем признать, что никогда не вписываемся не только в настоящее, но и в какой-либо социальный порядок. Неуместность --- не то, что нужно исправить, а условие свободы.

\chapter{От сообщества к обществу}

\section{Выйти в мир}

Принятие отчужденности делает нас общественными существами. Отчуждаясь, субъект смещается и оказывается вне самого себя. Он становится открытым для встречи с другими отчужденными субъектами, что порождает общество. Каждая встреча с другим создает общий мир. И даже если такой мир быстро замыкается в символических идентичностях, его кратковременное появление напоминает, как выглядят равенство и солидарность. Это совместное существование в условиях общей отчужденности. Обращать внимание на свою отчужденность --- значит добиваться признания себя субъектом.

Юрген Хабермас был первым, кто увидел в публичной сфере политическое пространство, где правит эгалитарная логика. И хотя публичная сфера порой недопускает одних и придает больший вес взглядам другим, она обещает всеобщий доступ и равный вес каждой высказанной в ней идее. Согласно Хабермасу, идеал публичной сферы имеет эгалитарный социальный эффект, пусть и не соответствует этим установленным собой же стандартам. Хабермас рассматривает публичную сферу как почти исчезнувшую с развитием капитализма. Капитализм, хотя и способствует рождению публичной сферы, приводит также к ее гибели.

В \textit{Структурной трансформации публичной сферы} Хабермас подробно опишет наступление капитала на публичную сферу. Благодаря господству товарной формы и под ее влиянием, представители публичной сферы становятся сугубо потребителями, мыслящими в категориях товара, и сфера эта фактически исчезает: ''Когда законы рынка, регулирующие сферу товарного обмена и общественного труда, распространились и на сферу частных лиц, рационально-критические дебаты возымели тенденцию замещаться потреблением, а сеть общественных коммуникаций превратилась в акты сколь угодно однородного, но все же индивидуального, восприятия''\footnote{Jürgen Habermas, The Structural Transformation of the Public Sphere: An Inquiry into a Category of Bourgeois Society, trans. Thomas Burger with Frederick Lawrence (Cambridge: MIT Press, 1989 [1962]), 161.}. Логика товара фрагментирует публичную сферу на изолированные индивидуальные действия, которые уже не объединяют людей в основанном на эгалитарных идеалах мире. Буржуазная культура, считает Хабермас, порождает публичную сферу как площадку для споров и дебатов, а затем капиталистическое общество разрушает ее по мере распространения коммерциализации.

Буржуазная публичная сфера развивается главным образом в привилегированных пространствах, таких как кофейни и салоны. Входя в такое пространство, человек, согласно Хабермасу, сохраняет свою частную идентичность и в то же время обращает ее к публике. Частные лица формируют публичную сферу, высказывая свои мнения в месте, где все такие мнения имеют равный вес.

Хабермас, впрочем, недооценивает как широту охвата публичной сферы, так и отчужденность, необходимую для ее создания.\footnote{Понимая важность публичной сферы для проекта эмансипации, Нэнси Фрейзер критикует концепцию Хабермаса за ее узость. Нам нужно множество публичных сфер, уверена она, чтобы учесть разнообразие культурных групп: ''Общественная жизнь в эгалитарных, мультикультурных обществах не может исчерпываться единой, всеобъемлющей общественной сферой. Это равносильно представлению различных риторических и стилистических норм в единой, веобъемлющей призме'' (Nancy Fraser, “Rethinking the Public Sphere: A Contribution to the Critique of Actually Existing Democracy,” Social Text 25/26 (1990): 69). Фрейзер утверждает, что мы взаимодействуем публично как частные личности. Она не видит взаимосвязи между публичной сферой и отчужденностью от идентичности.} Публичность возникает во встрече отчужденных субъектов, отбросивших свои личные заботы и открытых к взаимодействию с другими. Находясь в публичном пространстве, мы отчуждены от своей личной жизни. Эта отчужденность и есть то, что составляет любое общество. Дистанцируясь от своего социального положения, мы создаем основу для всеобщего равенства. Среди отчужденных существ никто не имеет особого статуса. Осознание отчужденности связывает нас со всеми остальными.

Препятствием к принятию отчужденности в публичной сфере становится наша инвестированность в идею сообщества. Никто не может избежать принадлежности сообществу, но когда мы вкладываемся в него и в соответствующую символическую идентичность, то забываем про отчужденность. Вклад в сообщество --- это вклад в социальное место, он скрывает нашу отчужденность, и тогда пропадает публичная сфера. Если субъект сам --- не через идентичность --- принадлежит сообществу, то за ним не признается универсальная отчужденность, в которой все равны.

Сообщество создает контекст, в котором субъект уже не один --- есть другие люди, и они поддерживают его идентичность. Сообщество обещает \textit{принадлежность}, но никогда не выполняет это обещание, отчего общество возникает вновь и вновь. Сообщество, в отличие от общества, кажется безопасным, но они существуют в диалектической взаимосвязи: общество существует благодаря неспособности сообщества замкнуться на себе, сообщество же возникает потому что никто не может существовать как лишь отчужденный, не имеющий идентичности, субъект.

Понимание отчужденности позволяет иначе взглянуть на сообщества: они формируются вокруг попыток избежать отчужденности. И пусть эти усилия всегда заканчиваются неудачей, они отвлекают нас от нашей отчужденности и следующего из нее равенства.

Сообщества могут обещать избавление от бремени субъектности в символической идентичности, с которой можно пытаться отождествить себя. Каждая символическая идентичность соответствует сообществу, одобряющему и поддерживающему ее, и без него была бы бессмысленной. Она существует благодаря признанию сообщества.

Сообщество прокладывает путь от отчужденной субъектности к символической идентичности, манит принадлежностью. Сообщество, кроме прочего, наделяет подчинившихся и вступивших в него субъектов, символической идентичностью. Видя в себе лишь члена сообщества --- как занимающего в нем определенную позицию, человек забывает о своей субъектности, забывает об отчужденности. Отчужденность кажется трудностью, которую можно преодолеть, усилив самоидентификацию с сообществом и своим в нем положением, и больше не чувствовать себя обреченным на отчужденность. Благодаря сообществу я знаю, кто я есть, и могу отчаянно отождествлять себя с этой идентичностью. На вопрос ''кто такой субъект?'' сообщество предлагает символическую идентичность.

Это отличает сообщество от общества. Общество выводит отчужденность субъекта на первый план. Это пространство --- физическое или интеллектуальное --- которым никто не владеет, но доступ к которому имеют все.\footnote{Проблема социальных сетей в том, что они --- общество в собственности корпораций или частных лиц, взявших под контроль пространство для публичных дискуссий и взаимодействий. И по мере того, как это становится очевидным, характерная для общества отчужденность все больше скрывается. В социальной сети можно сохранять удобную себе идентичность или даже развивать новую.} В капитализме, требующем от всех оставаться отдельными, частными существами, публичная сфера постоянно подвергается нападкам, потому что составляет ему эгалитарную угрозу.

Ни одна конкретная группа не котролирует публичную сферу, и потому она остается пространством, где на первый план выходит не символическая идентичность, а отчужденная субъектность. В обществе субъект не может полагаться на символическую идентичность. Например, в парке я могу столкнуться с кем-то, кто мне не нравится, или с тем, чьи убеждения противоречат моим. Но я не могу выгнать его из парка. Публичность парка значит, что он открыт для каждого, к какому бы сообществу он ни принадлежал. В публичной сфере субъект постоянно вопрошает, кто он такой, ищет ответ во взаимодействии с другими, столкнувшимися с той же проблемой.

Конечно, даже публичный парк не для всех одинаково доступен: кто живет рядом или у кого есть личный транспорт, доберется до него легче и быстрее. Но, все же, никто не может быть исключен из этого публичного пространства на основе своей символической идентичности. В основе публичного пространства лежит идея всеобщего равенства.\footnote{Общественное пространство --- форма для обитания. В \textit{Порядке форм} Анна Корнблю выступит за переосмысление формы как основы нашей социальности: ''Формализация как идеал заставляет нас думать о формальном государстве и формах социальности, подчеркивающих собственный антагонистический характер'' (Anna Kornbluh, The Order of Forms: Realism, Formalism, and Social Space (Chicago: University of Chicago Press, 2019), 154). Форма структурирует общество, не прибегая, в отличие от сообщества, к исключительности. Согласно Корнблю, форма становится основой для всякого возможного коллективного существования, выделяя, а не подавляя противоречия. Отворачиваясь от нее, мы не можем увидеть связь с другими в нашей универсальной отчужденности.}

Публичную сферу важно развивать, только она обеспечивает универсальность, недоступную никакому сообществу. В публичной сфере моя символическая идентичность ничего не значит: я могу быть методистом или членом организации Elks, но это не делает моей собственностью качели в парке или книги в публичной библиотеке. Ни качели, ни эти книги мне не принадлежат. Не принадлежат они и всем сразу. Мы отчуждены от содержимого публичной сферы, оставаясь наедине со своей субъектностью, своей уникальностью, которая и объединяет каждого со всеми. Благодаря отсутствию исключительности, публичная сфера демонстрирует универсальность. Нет способа включить кого-то в общество, и потому нет способа его исключить.

В отчужденности --- основе общественной жизни --- нельзя иметь преимущества или недостатка перед другими. Общество защищает от неравноправного распределения идентичностей в сообществе. Сообщество существует в недопущении определенных людей, общество же открыто для всех --- иначе не было бы публичным. Открытость общества оставляет субъекта без ответа на вопрос, кто он. В публичной сфере я не кто-то --- я единственный в своем роде. Встреча с другими в обществе бросает вызов моей символической идентичности. Если я начну пререкаться на дороге общего пользования, мне могут сказать: ''Кем ты себя возомнил?'', и логика общества не позволит мне ответить ''Я член клуба экстремального вождения, и потому грубо говорю с вами'' --- апелляция к сообществу не соответствует ей. Конечно, я все равно могу так ответить, но так я закроюсь от общества в своем сообществе. Сообщество соблазняет субъекта обещанием однозначного ответа на любой вопрос о нем.

Структура сообществ лучше видна в их отношении к своим членам. В США мое национальное сообщество говорит мне: ''Совершенно ясно, кто ты --- ты американец, вот и все''. И я могу находить утешение в своей американскости в любой точке мира: даже без американского паспорта я буду уверен, что я всего лишь американец и этим исчерпываюсь. У меня всегда будет Место, даже если я физически от него далек. Религиозное сообщество может ответить на вопрос о моем будущем после смерти,\footnote{Серен Кьеркегор постарается уберечь христианство от превращения в сообщество, не скрывать отчужденность субъекта. Именно потому он по праву считается основоположником экзистенциализма. Для Кьеркегора \textit{Христианский мир} означает форму христианства, ставшего сообществом --- и он неустанно осуждает переход к такой форме. Кьеркегор противопоставляет верующего субъекта и человека, принявшего христианскую идентичность в сообществе, чем стремящегося избежать экзистенциальных вопросов. Кьеркегор уверен, что христианство должно выделять, а не скрывать эти вопросы.} чем ослабит тревогу о смерти вплоть до уверенности неуязвимости. Сообщество дает мне готовые ярлыки, чем соблазняет отождествить себя с ними, забыть о своей свободе --- своей субъектности --- и стать не более чем ярлыком.

Сообщество защищает от неразрешимого для субъекта, экзистенциального вопроса. Оно дает символическую идентичность, чем может подтолкнуть субъекта вообразить, будто он к ней и сводится. Но чем больше субъект инвестирует себя в сообщество, тем больше ощущает, что не сводится к членству в нем. На это ощущение и стоит обратить внимание.

\section{''Очень большая плата''}

За защиту от вопроса субъектности сведением к символической идентичности приходится платить. Чтобы получить поддержку моей символической идентичности от сообщества, я должен подчиняться конституирующим его неявным требованиям. Конечно, у большинства сообществ есть четкие правила --- они могут быть даже прописаны в документе и вывешены на видном месте, что подчеркнет их значимость. Но решающую роль всегда будут играть неявные правила. Неявные правила могут даже указывать, как относиться к явным --- каким беспрекословно подчиняться, какими иногда пренебрегать и т.д.. Эти неявные правила никак не обнародуются и держатся в тайне. Члены общества даже не обсуждают их открыто. Тайные правила связывают членов общества, отделяя их от непосвященных.\footnote{Славой Жижек --- известный теоретик таких сплочающих общество, неписаных законов. Он утверждает, что они необходимы для создания связи с сообществом, поскольку позволяют отличать тех, кто действительно ему принадлежит от тех, кто лишь притворяется, что инвестирован в него. В \textit{Метастазах наслаждения} он напишет: ''Явных общественных правил всегда недостаточно, их надо дополнить тайным ''неписаным'' кодексом для выявления тех, кто хотя и не нарушает явные правила, все же сохраняет внутреннюю дистанцию и не решается идентифицировать себя с ''духом сообщества'''' (Slavoj Žižek, The Metastases of Enjoyment: Six Essays on Woman and Causality (London: Verso, 1994), 55). Жижек далее добавляет, что неписаные правила не просто отличаются от общественных --- они должны противоречить им, чтобы следование тайным правилам приводило к нарушению явных.}

Конечно, не каждое сообщество руководствуется столь уж явно противоречащими его публичному облику, скрытыми требованиями, но ни одно сообщество не обходится бе скрытых требований, создающих эксклюзивность. Сообщества должны быть эксклюзивными, что возможно только посредством скрытых требований. Даже самое открытое сообщество предъявляет к своим членам неявные требования. Так сообщество обеспечивает безопасность символической идентичности, заставляет субъекта платить за эту безопасность строгим подчинением. Ни одно сообщество не предлагает символическую идентичность бесплатно.

Подчинение правилам сообщества значит также противостояние аутсайдерам --- отказывающимся подчиняться неявным правилам или не отличающим явные правила от неявных и потому раскрывающих тайну. Через выделение аутсайдеров сообщество приобретает свою эксклюзивность, без них членство в нем не было бы столь привлекательно.

Особенно актуально это для сообществ, не осознающих себя сообществами. Например, большинство членов сообщества белых не осознают себя частью этого сообщества.\footnote{Даже излюбленный жест современных белых либералов --- подчеркивание привилегий белых --- еще не значит признания сообщества белых. Напротив, ссылаясь на их привилегии, они создают у себя и других впечатление, что белые --- вовсе не сообщество, ведь потому за ними и закреплены эти привилегии. Якобы, все подчиняются диктату сообщества, но белые выше этой борьбы --- это и есть привилегия, на которую они ссылаются.} Это сообщество формируется не вокруг цвета кожи его членов, а посредством действий, указывающих на подчинение скрытым требованиям: делать расистские заявления, без возражений выслушивать их, участвовать в расистских акциях --- все это составляет такое сообщество. Одной лишь белокожести недостаточно.\footnote{Раз белокожести недостаточно для принадлежности сообществу белых, то некоторые белокожие --- отказавшись подчиняться неявным правилам этого сообщества --- могут так стать для него аутсайдерами. Это одно из проявлений отчужденности.} В сообществе белых нет явных правил, ведь оно отрицает, что является сообществом.\footnote{Одна из форм антирасистского активизма --- выделение сообщества белых, чтобы не позволить ему скрываться. Действуя скрытно, сообщество обладает гораздо большей властью над своими членами.} Оно, как и любое другое сообщество, эксклюзивно, и его скрытые правила направлены на закрепление этой эксклюзивности. Потенциальный член сообщества должен доказать, что всецело идентифицирует себя с сообществом, следовать неявным правилам --- только так он сможет добиться полного признания в нем. Субъект обращается к сообществу в надежде стать собой, и в ответ получает требование отдать свою субъектность.

Безопасность, обещаемая сообществом, всегда иллюзорна. Никакая символическая идентичность не заменяет субъектность. Даже в уюте сообщества и соответствующей идентичности человек остается субъектом. Вопрос остается без ответа. Сообщество выдает мне символическую идентичность, но она в конечном счете вступает в самопротиворечие. Безопасность оказывается небезопасной.

Сообщество может быть обманчиво: инвестируясь в него в попытке убежать от экзистенциального вопроса, мы представем пред тем же вопросом, но теперь куда более навязчивым. Мы осознаем свою неспособность всецело принадлежать сообществу. Нам может казаться, что уж другие то по-настоящему ему принадлежат, и только нам это никак не удается. Сообщество не дает субъекту охватывающего чувства принадлежности, потому что субъект всегда остается субъектом. Он никогда не сливается с символической идентичностью. Отчужденность никуда не уходит. Сообщества могут побуждать забыть об отчужденности, отбросить экзистенциальный вопрос, но в перспективе это приводит лишь к большему дискомфорту и даже ужасу.

Символическое положение членов сообщества может казаться совершенно надежным. Но в любой момент тайные правила --- противоречащие явным --- могут вскрыться публичной общественности, что разрушит корректность символической идентичность членов сообщества в глазах общества. Да и никто никогда не ощущает себя всецело членом сообщества, каждый должен постоянно следовать неявным правилам, проходить испытание вновь и вновь, чтобы доказать, что он достоин.

Символическая идентичность, когда видится лекарством от отчужденности, не только не обеспечивает обещанную безопасность, но и подвергает постоянному риску: субъект никогда не знает, достаточно ли он делает для ее поддержания. Остальные члены сообщества всегда могут решить исключить меня из него, если им покажется, что я недостаточно соответствую. Символическая идентичность ненадежна, и я всегда нахожусь на грани ее утраты. Она создает образ безопасного, прямолинейного существования, но оно принципиально невозможно для субъекта.

Публичное пространство --- общество --- не обещает такой безопасности.\footnote{Помимо термина \textit{public} (\textit{публичное пространство}) есть также термин \textit{commons} (\textit{общее достояние}/\textit{общество}). Некоторые теоретики ассоциируют общество (commons) с сообществом (community), но между ними есть фундаментальное различие: общество не обеспечивает той принадлежности, которая составляет сообщество.} Опасность отчужденности очевидна при любом публичном взаимодействии. Говоря со случайным встречным в парке, я не нахожусь в личном пространстве: возможно, мой собеседник не принадлежит к тем же сообществам, что я, возможно, он сочтет мои шутки оскорбительными, или высмеет мой акцент, или просто отвернется и уйдет. В обществе символическая идентичность бессильна.

Но общество позволяет установить связь с другими --- связь, не основанную на хрупкой символической идентичности. Связь эта устанавливается через отчужденность, а не через ее сокрытие. В публичном пространстве никто не чувствует себя как дома --- все оказываются равны. В нем нет свойственных сообществам иерархий и ограничений, и каждый в равной степени свободен. В публичной сфере мы одиноки вместе --- и это единственная возможная форма солидарности.\footnote{В \textit{Недовольстве культурой} Зигмунд Фрейд утверждает, что социальный порядок фундаментально антагонистичен субъекту, даже в обществах, ориентированных на коллектив больше, чем современное Фрейду либеральное: ''Значительная часть борьбы человечества сосредоточена вокруг поиска приемлемого баланса между притязаниями индивида [на личную свободу] и культурными притязаниями группы, и одна из проблем здесь --- есть ли вообще такой баланс, можно ли его достичь конкретной формой цивилизации, или же этот конфликт останется неразрешенным'' (Sigmund Freud, Civilization and Its Discontents, trans. James Strachey, in The Standard Edition of the Complete Psychological Works of Sigmund Freud, ed. James Strachey, vol. 21 (London: Hogarth, 1961 [1930]), 96). Фрейд оставляет вопрос открытым, но на протяжении всей книги показывает, что никакого решения этот конфликт иметь не может: какие бы политические реформы мы не провели, антагонизм между обществом и субъектом сохранится.}

\section{Встречайте нового босса}

Сообщество или культура формируется через утверждение главного означающего (master signifier), которое объединяет сообщество. Сообщество нуждается в таком означающем как в точке опоры для самоидентификации. Этот основополагающий термин --- например, \textit{американец}, \textit{кореец}, \textit{игрок студенческого футбольного клуба} --- служит организующим звеном для сообщества, символической отправной точкой. Каждый может указать на это означающее как на основу соответствующей сообществу идентичности. Главное означающее само по себе --- не означаемое, не имеет значения. Означающее приобретает значение только через установление позиции непринадлежности --- иначе говоря, оно негативно. Идентичность любого конкретного сообщества определяется тем, кто ему не принадлежит. Этот акт создания позиции непринадлежности --- необходимое условие для любой самоидентификации в сообществе. Я тот, кто не другой. Не бывает позитивной самоидентификации --- самоидентификации, не опирающейся на непринадлежность.

Обладатели идентичности должны яро осуждать аутсайдеров, угрожающих их сообществу или культуре. Негативность главного означающего указывает в аутсайдере угрозу, которая необходима для утверждения сообщества или культуры --- сообщвество не может сформироваться и удерживаться иначе.

Примером последствий слишком сильной негативной функция главного означающего являются правоэкстремистские движения. Нативисты, например, используют название нации в качестве главного означающего. Но без гипертрофированной негативной функции, вкладываемой ими в это означающее, нативисты не были бы нативистами: иммигранты должны видиться угрозой безопасности национального сообщества --- только тогда сообщество нативистов сможет существовать.

Дональд Трамп начал свою президентскую кампанию 2016 года с обещания построить стену и так прекратить поток иммигрантов с юга в США. Он провозгласил, что сделает Америку снова великой, и в качестве средства использовал негативное отношение к иммигрантам. Призыв Трампа ничего бы не значил, если бы не определял негативно --- через иммигрантов --- Америку и ее величие. Главное означающее структурирует сообщество только через соотнесение его с теми, кто ему не принадлежит.

Либерализм, хотя и противостоит нативистским движениям, имеет ту же динамику. Либеральное сообщество существует лишь пока ему угрожает какая-нибудь организация. Либерализм направлен на тотальную инклюзивность --- создание сообщества столь большого, что никто не останется вне него. Такое сообщество --- либеральная фантазия, требующая, чтобы кто-то действовал отдельно и жестко и тем самым продолжал угрожать либеральному сообществу. Либеральная открытость невозможна без тех же стен, что сделают, уверен Трамп, Америку снова великой.

Осудив ''кучку жалких''\footnote{анг. Basket of deplorables, буквально ''корзина достойных сожаления'' --- уничижительная фраза из речи Хиллари Клинтон, кандидата в президенты США от Демократической партии, произнесенная ей 9 сентября 2016г на мероприятии по сбору средств для предвыборной кампании. Она использовала эту фразу для описания ''половины'' сторонников своего оппонента, кандидата от Республиканской партии Дональда Трампа, назвав их также ''расистами, сексистами, ксенофобами и исламофобами'' (пр. переводчика).} во время президентской кампании 2016г, Хиллари Клинтон непреднамеренно вскрыла структурно необходимое для либеральной позиции: она сказала то, что должно было остаться невысказанным, ясно показав невозможность инклюзивности.\footnote{Хуже того --- Клинтон обозначила консервативные группы как действительно свободные, ведь ''корзина достойных сожаления'' (ориг. basket of deplorables) больше подходит объединившимся в своей непринципиальности, чем утверждающим свою принадлежность через остракизм евреев, иммигрантов и других лиц.} Либеральное сообщество невозможно без тех, кто угрожает ему --- только они придают значение либеральной идентичности, которая иначе была бы пресной и бессмысленной. Всегда где-то должна быть ''кучка жалких'' --- даже если о них не говорят,--- чтобы либерализм мог самоопределиться.

Инклюзивность работает ровно до тех пор, пока остается всеобщей инклюзивностью \textit{в будущем}. Как только такой идеал будет воплощен, он тут же рухнет, потому что некому будет занять противоположную позицию, пропадет удерживающая всех вместе точка отсчета: каждый член сообщества принадлежит ему ровно потому что не является чуждым ему, не является его врагом.

В этом смысле любое сообщество --- закрытое, каким бы открытым и гостеприимным ни казалось. То, что некоторые критики считают оксюмороном: ''Закрытое сообщество --- не совсем сообщество'' --- верно с точностью до наоборот, потому что закрытое сообщество имеет парадигмальный статус по отношению ко всем прочим.\footnote{В прошлом я был глуп и тоже придерживался этой позиции --- я яростно осуждал закрытые сообщества, будучи уверен, что закрытость уничтожает сообщество, тогда как на самом деле она его и составляет.} Сообщества формируются через недопущение посторонних. Будучи открыто для всех, сообщество растворяется в общество, становится публичным. Сообщество невозможно без исключенных из него.

В известном анекдоте о коммунизме Ленин после смерти встречается с Марксом. Тот поздравляет его и говорит: ''Как я и ожидал, вы захватили власть во имя коммунизма, провели экономические реформы в должном темпе и даже поработали над реформой политического мышления в странне. Отлично! Но теперь марш революции, похоже, в Сталине продолжается\footnote{встал и не продолжается}''. То есть коммунистический проект, даже будучи универсален и успешен, требует аутсайдера, и в этом анекдоте им выступает Сталин.

Коммунизм в этом анекдоте подается как могущий быть непротиворечивым, если бы не вмешался Сталин, как если бы в самом коммунистическом проекте не было никакой почвы для ущерба, нанесенного ему Сталиным. В анекдоте коммунизм --- дело рук сообщества, а Сталин --- проблема, которую нужно преодолеть, но нет указания на возникновение общества в этом прерывании, так что анекдот этот партикуляристский. Маркс был признанным универсалистом, однако в анекдоте возлагает вину за провал коммунизма на Сталина и тем знаменует поворот к партикуляризму общества. Конечно, вряд ли справедливо обвинять Маркса за то, что он говорит в придуманном мной анекдоте. Но ведь именно такой позиции придерживаются многие марксисты по отношению к Сталину. В этом смысле анекдот показывает, как мышление, свойственное членам сообщества, перебивает универсалистскую интерпретацию событий даже среди универсалистов.

Сообщества предлагают своим членам признание в обмен на отождествление с предоставленными символическими позициями. Но сообщество существует за счет исключения, а не включения. Нет сообщества без тех, кого в нем нет: они скрепляют членов сообщества в противостоянии себе. Общество, в отличие от сообщества, куда более хрупко, потому что для него не существует аутсайдеров, оно никого не исключает и потому нестабильно.

\section{Садись в автобус}

Разница между обществом и сообществом хорошо видна в разнице между общественным автобусом и заказным. Садясь в общественный автобус, я оказываюсь среди незнакомых мне людей. Разговорчивые соседи могут мешать мне вздремнуть. Мне нужно подстроить свое расписание так, чтобы успевать на автобус. И у меня нет никаких привилегий по сравнению с другими пассажирами, кем бы я ни был на работе, дома или среди своих друзей. В автобусе может ехать любой желающий --- богатый и бедный, китаец и бразилец, интеллигентный и не очень, мусульманин и буддист и т.д.. Ни одно сообщество и ни одно лицо не владеет автобусом тотально. Общественный автобус --- место отчуждения, и именно потому множество людей отказывается от поездок на нем.\footnote{Переезжая с супругой в Лос-Анджелес, мы посоветовались с живущей там знакомой. Она рассказала, где лучше поесть, куда сходить в кино и порекомендовала Голливуд Боул, где проводятся музыкальные мероприятия (в Голливуд Боул за \$1 за билет можно было увидеть исполнение Девятой симфонии Бетховена). Но она также предостерегла нас от поездок на автобусе, потому что ''из-за дешевизны билетов ездить может каждый''. Публичный характер автобуса, она была уверена, представлял опасность для ее и нашей субъектности.}

В общественном автобусе --- в отличие от частного заказного --- символический статус пассажиров не определяет, какие места им можно занимать. Нельзя купить билет первого класса --- нужно садиться там, где есть свободное место. Может оказаться, что вы сядите рядом с человеком из совсем другого сообщества. Антисемит может оказаться рядом с евреем, израильтянин рядом с палестинцем, мусульманин рядом с ярым атеистом и т.д.. Общественность автобуса значит вероятность встретить того, кого встречать бы не хотелось. Такой автобус --- чуждое пространство для каждого пассажира.

В общественном автобусе мы видим одно из ключевых свойств публичной сферы --- отсутствие конкретных мест для каждого. Я должен стоять или сесть на оказавшееся свободным место, а иногда, когда автобус переполнен, я могу только стоять. В автобусе для меня нет выделенного Места. Публичное пространство для всех отчужденное, в нем каждый смещен со своей символической идентичности. Есть сидения, но нет Мест. Их отсутствие делает публичную сферу местом свободы.

В публичном пространстве отчужденный субъект выходит на первый план. Символическая идентичность, сформированная сообществом, здесь ничего не значит. Он испытывает внутреннюю уязвимость, потому что не прикрыт этой идентичностью. И хотя в общественном транспорте, конечно, действуют законы государства --- запрещающие, среди прочего, нападения, оскорбления, сексуальное насилие и воровство --- субъект не чувствует в нем себя как дома.

Публичный статус автобуса проявляется также в том, как настойчиво люди стремятся абстрагироваться от него: один надевает наушники и слушает музыку, другой спит, третий пялится в свой телефон, четвертый смотрит в окно, чтобы не видить окружающих. Все это --- попытки избежать столкновения с публичным пространством и отчужденностью. Но самой вопиющей формой неприятия публичности в автобусе было официальное разрешение сообществу вторгаться в это пространство.

В одно время на юге США, в Джим Кроу, места в общественном автобусе распределялись в соответствии с символическим статусом пассажира: белокожие могли ехать впереди, темнокожие же должны были садиться в задней части. Автобус изначально был общественным, но такое правило лишило его этого статуса.

Движение за гражданские права было попыткой покончить с доминированием белого сообщества. Оно предполагало создание подлинного общества, в котором ни одно сообщество не доминирует над другим. Роза Паркс была лидером в этой борьбе, в том числе за публичное пространство в автобусе. Она стала символом движения за гражданские права, когда отказалась уступить свое место белокожему в общественном автобусе в Монтгомери, штат Алабама.\footnote{Роза Паркс хотя и стала символом движения, была не первой: Клодетт Колвин отказалась уступить свое место в общественном автобусе в Монтгомери за несколько месяцев до Паркс. Но Колвин была незамужней беременной девушкой-подростком, и потому символом была выбрана Роза.}

В чартерном автобусе сообщество и вовсе доминирует над публичной сферой. Автобус, в котором мы с командой колледжа ездили на матчи, позволял нам чувствовать себя как дома. Поездки бывали довольно долгими --- до 12 часов, но каждому всегда находилось свободное место. Мы могли быть уверены, что будем сидеть там, где захотим и общаться с теми, с кем захотим, продолжая принимать участие в деятельности сообщества, пусть и в таком вот ограниченном пространстве.

Заказной транспорт, в отличие от общественного, защищает сообщество от отчуждения в публичную сферу. В нем каждый воспринимает других как часть своего сообщества, благодаря чему может действовать импульсивно и в уверенности, что при соблюдении правил сообщества всегда получит в ответ ожидаемую, положительную реакцию.

Заказной автобус скрывает отчужденность, тогда как общественный --- подчеркивает ее. В общественном автобусе пассажир смещен со своего символического статуса, но в этом проявляется его субъектность. В заказном автобусе вам всегда уготовано место, но только если вы смиритесь с требованиями сообщества и приходящим с ним неравенством. Заказной автобус в этом подобен круизному лайнеру.

Одна из загадок кинематографа --- почему \textit{Скорость} (Ян де Бонт, 1994) оказался столь успешней своего сиквела ''Скорость 2: круиз-контроль'' (Ян де Бонт, 1997). Оба фильма снял Ян де Бонт, в каждом главную женскую роль сыграла Сандра Буллок, да и замена Киану Ривза на Джейнсона Патрика в главной мужской роли едва ли может объяснить такой провал. Но в оригинальном фильме действие происходит в общественном автобусе, а в сиквеле --- на частном круизном лайнере. Общественный автобус в \textit{Скорости} обеспечивает универсальную связь между героями фильма и зрителями, тогда как на частном круизном лайнере герои находятся в своем сообществе, изолированном от зрителя. Смотря \textit{Скорость 2}, зритель не слишком переживает за персонажей, потому что они остаются в рамках своего сообщества. В оригинале же зрителям видна отчужденность пассажиров автобуса, присущая каждому субъекту. Персонажи в оригинале испытывают определенный дискомфорт, находясь в публичной сфере, но именно это привлекает зрителя, помогает ему представить себя вместе с ними в том же автобусе. Круизный лайнер, подобно заказному автобусу, создает кажущееся безопасным пространство, но оно подчинено сообществу и закрыто для посторонних. Зрители \textit{Скорости 2} оказываются аутсайдерами, тогда как зрителям \textit{Скорости} удалось прокатиться на автобусе вместе с героями фильма. Круизный лайнер постигла участь заказного автобуса, и проблема усугубилась также бесчисленными товарами, окружающими персонажей и помогающими им уединиться.

Теснота общественного транспорта демонстрирует свободу и равенство, которые скрывает заказной транспорт. Конечно, ездить в заказном автобусе приятней: можно не бояться нежелательных соседей и собеседников. Но заказной автобус куда дороже общественного, ведь выбирая его, мы выбираем неравенство вместо всеобщей свободы в отчужденности.\footnote{Ханна Арендт будет решительно защищать публичную сферу от нападок сообществ. Согласно Арендт, публичная сфера не чужда, и у нее есть мир, в котором люди чувствуют себя как дома, к которому они могут принадлежать. В \textit{Vita activa, или О деятельной жизни}, она утверждает, что общество формируется преодолением отчужденности, а не ее принятием: ''Отчуждение от мира --- а не от самого себя, как думал Маркс --- стало отличительной чертой эпохи модерна'' (The Human Condition, 2nd ed. (Chicago: University of Chicago Press, 1998 [1958]), 254). Ханна выступает за преодоление отчужденности. Несмотря на ее озабоченность публичной сферой, Арендт уверена, что настоящая публичность лишена отчужденности. Но попытка преодолеть отчужденность всегда ведет в сообщество, в частную, закрытую жизнь.}

\section{Прерванное единение}

Чаще всего солидарность связывают с общей человеческой сущностью, а не отчужденностью. Оно и понятно: отчужденность, в отличие от сущности, звучит негативно. Среди защитников этой позиции есть и связывающие человеческую сущность с местом человека в мире природы гуманисты, и отождествляющие ее с искрой божественного творения теисты. С чем бы не связывалась человеческая сущность, она удобна тем, что охватывает всех и только людей, согласно избранной автором концепции человека.\footnote{Кваме Энтони Аппиа в своем \textit{Космополитизме} модифицирует гуманизм для учета культурных различий. При этом он признает, что согласия касательно универсальных ценностей может и не быть: ''Космополиты полагают, что у всех культур достаточно общих ценностей, чтобы вступить в полилог. Но, в отличие от некоторых универсалистов, они не утверждают, будто эти культуры тогда обязательно придут к согласию'' (Kwame Anthony Appiah, Cosmopolitanism: Ethics in a World of Strangers (New York: Norton, 2006), 57). То есть Аппиа, в отличие от гуманистов, уверен только в возможности диалога о ценностях, но не в существовании общечеловеческих ценностей.} 

Особенно авторитетной эта позиция была в начале нового времени, с XVII до конца XVII века. В 1689г Джон Локк напишет два своих трактата о государственном управлении, где сформулирует натуралистическую концепцию сущности человека. Свобода, будет утверждать он, проистекает не из какой-то конкретной характеристики человека, а из его естественного статуса: ''Естественная свобода человека заключается в свободе от всякой высшей Власти на Земле, в способности не подчиняться внешней Воле или Авторитету и руководствоваться только Законом природы''\footnote{John Locke, Two Treatises of Government  Cambridge: Cambridge University Press, 1988 [1689]), 283}. Никто, считает Локк, не может поработить или подчинить человека, потому что это нарушило бы человечность порабощенного.

К той же позиции можно прийти через Бога. Декларация независимости, в соответствии с которой в 1776г были основаны США, приписывает общечеловеческие ценности действиям божественного творца. Права человека, согласно Декларации, неотчуждаемы, потому что люди были ''наделены ими своим Создателем''. То есть, человеческая сущность берет начало в акте трансцендентального существа, и этот акт гарантирует нам определенные права.

Но проблема с общей человеческой сущностью как источником ценностей станет очевидна, если изучить поступки ее авторов. Локк инвестировал в рабовладельческую компанию и участвовал в создании конституции колонии, где рабство было законно. Томас Джефферсон --- главный автор Декларации независимости --- владел рабами и управлял нацией, построенной трудом рабов. Понятие человеческой сущности не помешало им исключить некоторых из сообщества людей.

И это не просто грехи каких-то отдельных теоретиков: исходит ли сущность человека от Бога, от природы или откуда-то еще, она не может быть признана универсальной, потому что образ общечеловеческого существа требует наличия нечеловеческого другого, который позволил бы провести границу человеческого. Человек определяет человеческое через контраст с нечеловеческим другим. Любое обращение к общечеловеческой сущности должно кого-то исключить.\footnote{Это также утверждает Чарльз Миллс в \textit{Расовом договоре}: человек становится человеком только благодаря исключению тех, кого человеком не считают или из кого еще предстоит сделать человека.}

Но с отчужденностью мы такого не увидим. Наша отчужденность --- не признак принадлежности, исключающий других, а признак неспособности, который eo ipso не может кого-либо исключить. Великое достоинство неспособности в том, что на нее способен каждый, тогда как включение всегда требует соответствующего исключения. Все в равной степени отчуждены --- независимо от их социального или материального статуса. Богатые переживают отчужденность столь же остро, сколь и бедные. Поэтому отчужденность --- единственное, что может связать всех в общей солидарности. Дистанция, отделяющая людей друг от друга и от самих себя --- это то, что их объединяет. Отчужденность универсальна. Сама способность относиться к себе как к субъекту зависит от структуры отчужденности, от дистанции в субъекте.

Универсальность отчужденности резко контрастирует со спецификой любого сообщества или культурной идентичности. Каждая культурная идентичность индивидуальна и требует других --- противоположных ей --- идентичностей, чтобы закрепить эту индивидуальность. Даже самая открытая культура остается замкнутой в своей специфике, противопоставляющейся другим культурам.

Но специфичность нашего сообщества не значит, что мы тождественны этой специфике. Никто не тождественен символической идентичности. За любым сообществом и любой культурной идентичностью субъекта таится отчужденность.

Сообщества могут прибегать к ритуалам и идеологической обработке в попытке преодолеть пропасть, отделяющую их от субъекта. Будь мы тождественны своим идентичностям, ничего из этого не потребовалось. Ритуал, призванный засвидетельствовать принадлежность, самим своим существованием свидетельствует об обратном --- об отчужденности субъекта от своей символической идентичности.

Общество разоблачает сообщества как противоречивые. Наша неспособность быть тождественными символическим идентичностям --- вот что нас всех объединяет. Нас объединяет то, чего у нас нет. Сообщества и культуры противоречивы, но именно эта противоречивость позволяет одной культуре говорить с другой --- в каждом сообществе есть место универсальности. Но мы упускаем из виду это универсальное противоречие, когда думаем об универсальности как о чем-то, что есть у всех нас общего, о чем-то объединяющем в себе все конкретные общества или культуры. Но универсальное --- не то, что мы можем иметь или навязывать другим, а недостаток, нехватка, присущая каждой культуре. И общество --- пространство смирения с этой нехваткой.

Когда большинство людей замыкаются в своих сообществах, общество подвергается риску вымирания. Оно живет только благодаря осознанию людьми их отчужденности, осознанию, достаточному для выхода в публичную сферу и вступления в контакт с другими отчужденными субъектами. Общество зависит от готовности субъектов признать свою неуместность. Строго говоря, общество есть всегда --- потому что отчужденность это не то, чего субъект может достичь или избежать, но субъект не всегда осознает себя общественным, публичным существом, потому что это требует принятия отчужденности.

Без общества как жизнеспособной альтернативы сообществу, становится практически невозможным видеть что-то вне своего сообщества. Кажется, что каждый живет в своих частных сообществах и окружен такими же закрытыми в своих сообществах другими.

Поиск идентичности, с которой можно было бы всецело отождествиться --- это попытка избавиться от отчужденности. Вот почему проекты, направленные на идентификацию, вызывают ужас. Общество разоблачает любую идентичность и подрывает ее притязания на замещение субъектности. Идентичность создает видимость тождественности субъекта и идентичности, даже если постоянно меняется. Признав несостоятельность любой идентичности как замены субъектности, мы остаемся в рамках экзистенциального вопроса субъектности, в публичной сфере.

Это признание значит принятие своей неуместности как единственно возможной связи себя как субъекта со своим сообществом. Для этого нужно привыкнуть к непривычному, чувствовать себя уютно в неуютном. Только так --- не торопясь с решением (что только усугубило бы ситуацию) --- можно решить универсальную проблему субъектности. Описание того, как следует вести себя в публичной сфере --- великое достижение Джейн Остин.

\section{Публикации Джейн Остин}

О Джейн Остин принято думать либо как о моралистке, либо как об авторе любовных романов. По мнению Джулиет Макмастер и многих других критиков Остин, главная черта ее повествовательной вселенной --- моральный порядок. Макмастер напишет: ''Мир Джейн Остин морально упорядочен и наполнен смыслом, что у нас, влачащих беспорядочную жизнь в своей культуре вседозволенности, вызывает тоскливую зависть''\footnote{Juliet McMaster, Jane Austen the Novelist: Essays Past and Present (New York: St Martin’s Press, 1996), 16. Обитатели вселенной Остин часто являются образцами добродетели. Аласдер Макинтайр скажет, что Остин --- ''последний великий представитель классической традиции добродетелей'' (Alasdair MacIntyre, After Virtue: A Study in Moral Theory (South Bend: University of Notre Dame Press, 1981), 226). Согласно Макинтайру, мы читаем романы Остин в поисках нравственных фигур, которым хотели бы подражать. Чтобы найти путь к добродетельности.} Те же, кто смотрит на Остин как на автора любовных романов, видят в ней новый взгляд на любовь: для Остин важна не конечная точка, которой достигнет пара, а путь, который будет пройден. В основе ее романа \textit{Доводы рассудка} лежит исходная неудача, неотъемлемая часть отношений. Робин Уорхол опишет \textit{Доводы рассудка} Остин как ''историю обретения утраченной любви, примирения противоположностей''\footnote{Robyn Warhol, “The Look, the Body, and the Heroine: A Feminist-Narratological Reading of Persuasion,” Novel: A Forum on Fiction 26.1 (1992): 5.}. В своих романах Остин переосмысляет мораль и любовь как невозможные для уединенных в своих символических идентичностях.

И морализм Остин, и ее описания романтической любви действуют в публичном мире. В романах Остин мораль и любовь возможны только когда человек принимает отчужденность в публичной сфере --- в замкнутом пространстве сообщества он никогда не находит их. Вселенная Джейн Остин --- это публичный мир, требующий определенных добродетелей. Остин показывает, что отчужденный субъект проявляется именно в обществе.

Герои романов Остин --- Элионор Дэшвуд в \textit{Чувственности}, Элизабет Беннет в \textit{Гордости и предубеждении}, Энн Эллиот в \textit{Доводах рассудка} --- проявляют уважение к публичной сфере, хотя это причиняет им страдания, потому что в обществе оказываются в том числе враждебно настроенные к ним люди. Каждая из этих женщин страдает из-за невозможности открыто заявить о себе миру, каждая сталкивается с внутренним барьером, мешающим познать себя. Пока Элизабет Беннет живет в своем собственном окружении, она остается уверена в своих желаниях. И только через общение с публикой --- и тому, что она воспримет как грубость мистера Дарси,--- она обнаружит свое бессознательное желание. Герои Остин видят себя в своей отчужденности, а не в попытке ее преодолеть. Они признают себя общественными существами.

В романах Остин есть еще одна категория персонажей --- злодеи. Они, в отличие от других персонажей, осознают отчужденность публичной сферы и используют ее в личных интересах. Непрозрачность отношений в публичной сфере позволяет Джорджу Уикхему в \textit{Гордости и предубеждении} и Уильям Эллиот в \textit{Доводах рассудка} манипулировать не столь проницательными людьми и вовлекать их в свои вредительские планы. Джордж Уикхем уговорит Лидию Беннет сбежать с ним, хотя и не намерен на ней жениться, что, конечно, нарушает принципы публичной сферы, но Уикхема заботят только сексуальные завоевания --- так он укрепляет свою символическую идентичность. Уикхем не идет на риск отчужденного субъекта. Злодеи тоже не понимают, что отчуждение в публичной сфере не уводит субъекта от себя, а раскрывает. Пытаясь совершенно отождествить себя с идентичностью, предложенной сообществом, субъект забывает себя.

Герои Остин понимают, что в обществе нельзя действовать прямо. Отчужденность требует непрямых действий, выражения желания в обход возможных скрытых препятствий. В публичной сфере между субъектом и объектом его желания возникает барьер, который, однако, позволяет субъекту осознать свою любовь. В своих романах Остин показывает, что любовь возможна только в отчуждении.\footnote{Согласно Бенг-Чхоль Хану, открытое выражение любви не только препятствует ей, но и превращает отношения в подобие порнографии, начинается ''там, где тайна исчезает в угоду тотальной демонстрации и обнажению откровенной, навязчивой позитивности'' (Byung-Chul Han, The Transparency Society, trans. Erik Butler (Stanford: Stanford University Press, 2015 [2012]), 26).}

Акцент на обществе в романах Остин особенно заметен в их кино- и телеверсиях --- часто в них значимость публичной сферы подчеркивается даже больше. Экранизация \textit{Доводов рассудка} Роджера Мичелла выйдет в 1995г и станет наиболее значимой попыткой показать важность общества.\footnote{Последующая киноверсия \textit{Доводов рассудка} (Кэрри Крэкнелл, 2022) уже не сделает такого акцента на отчужденности в обществе.} \textit{Доводы рассудка} --- последний роман Остин --- повествует об отношениях Капитана Фредерика Уэнтворта и Энн Эллиот, возобновившихся спустя годы после того, как Энн отвергла Уэнтворта по совету своей подруги леди Рассел. Их отношениям суждено пробиваться сквозь оскорбления семьи Энн --- она станет чужой даже у себя дома. Сообщество угнетает Энн, и единственным отдохновением становится общество. Уэнтворт и Энн любят друг друга на протяжении всего романа, но осуществлению их любви препятствует общество и отчужденность в нем. Роман заканчивается традиционно, браком, но Остин стремится показать, что это союз равных, что оба партнера помнят о своей отчужденноси в отношениях. Ни Энн, ни Уэнтворт не прячутся за символической идентичностью (хотя, конечно, она есть у каждого), отчего часто неспособны сказать друг другу о своих чувствах, и на фоне этой неспособности разворачиваются все события в романе.

Вульгарная открытость и прямота удушают публичную сферу, от них она распадается на конкурирующие сообщества. Во вселенной Остин общество требует от субъекта маскировки, желание может быть проявлено только косвенно.

В экранизации \textit{Доводов рассудка} хорошо показано, как публичная сфера возводит барьер из условностей на пути выражения желания: Энн Эллиот не может просто сказать Капитану Уэнтворту о своих чувствах, как не может и он сказать ей. Но эта их неспособность прямо донести свое желание не только не портит фильм, но и доставляет зрителю удовольствие.

Отец Энн сэр Уолтер и ее сестра Элизабет отказываются считаться с требованиями общества: они цепляются за свою символическую идентичность и эксплуатируют общественные отношения для продвижения личных интересов. Будучи вовлечен в беседу, сэр Уолтер всегда разглагольствует об уродстве людей, которых он видел. В начале фильма он откажется сдавать свое поместье моряку потому что у того огрубевшая кожа и суровый вид.

Уолтер и Элизабет открыто выражают неприязнь к Энн --- такая прямота недопустима в обществе,--- и все же Энн не отвергает общество, реагируя молчанием. На протяжении всего фильма Энн будет появляться на публике, соблюдать правила общественного поведения, никогда прямо не заявляя о своих чувствах к Капитану Уэнтворту. Несмотря на жизнь в полной изоляции, Энн никогда не будет одна. Ее публичное молчание не оскорбляет и не привлекает собеседников, она хранит его даже когда ей высказывают предположения, на ком женится Капитан Уэнтворт. Позиция общественного молчания подчиняется общественным кодексам поведения, не претендуя на их изменение сообществом, и показывает принятие Энн своей отчужденности.\footnote{Энн Элиот --- самая тихая из героинь Остин. В ранних ее романах уже показана неизбежность непонимания в публичном пространстве, но роль отчужденности особенно явно видна в \textit{Доводах рассудка}.}

Казалось бы, строгие ограничения публичной сферы будут только угнетать Энн, но в фильме мы видим, что ее окружение --- ее семья --- куда более деспотично, чем общество. В отчужденности публичной сферы Энн найдет недостающую свободу, и именно там будут развиваться ее отношения с Уэнтвортом. Даже в конце фильма, будучи уже супругами, они находятся на корабле, а не в домашней обстановке.

Публичность всегда содержит преграды прямому выражению желания. Сколько бы Энн и Капитан Уэнтворт ни встречались, они не могут выразить его. Барьер остается даже в, казалось бы, личной беседе: оставшись с Энн наедине в карете, Капитан Уэнтворт предпримет неудачную попытку прямого общения: ''Я сожалею... чертовски глупо, чертовски глупо''. Позже, на концерте В Бате, он снова попытается сказать ей о своем желании: ''Энн, я никогда...'' --- но его прервет появление леди Далримпл. Увидив то, что посчитает близостью между Энн и мистером Эллиотом во время концерта, Капитан Вентворт сердито направится к выходу.

Все это непонимание возникает из-за отчужденности. Видя, что Уэнтворт уходит с концерта, Энн пытается остановить его, но не может сказать единственное, что заставило бы его остаться,--- что она любит его. Этот их разговор станет, пожалуй, самым неприятным во всем фильме: зритель знает, чего хотят герои, но вынужден наблюдать отсутствие взаимопонимания между ними. Энн умоляет Уэнтворта остаться, говоря: ''Следующая песня очень красивая. Очень красивая песня о любви. Разве она не стоит того, чтобы остаться?'', на что Уэнтворт отвечает ''Мне здесь оставаться незачем''. Энн принимает свою отчужденность и не выражает желание напрямую --- не привносит личное в публичную сферу --- отчего не может убедить Уэнтворта остаться.
Капитан Уэнтворт тоже понимает важность обходных путей. Даже в конце концов сообщая Энн о своем желании, он делает это в письме, оставленном у нее на столе. Он пишет, что прибегнул к этому потому что ''должен говорить с вами доступными мне средствами'' и просит ее выразить свое желание косвенно: ''Одного слова, одного взгляда будет достаточно, чтобы решить, войду я в дом твоего отца сегодня вечером --- или никогда''. Сразу же после прочтения письма Энн сталкивается с Уэнтвортом на улице и слегка берет его за руку. Этот жест --- предел выражения ее любви --- что позволяет ей сохранить себя в публичной сфере как субъекта.

Раз теперь Уэнтворт и Энн точно знают желание друг друга, они, можно ожидать, начнут открыто говорить о своих чувствах. Возможно, так и происходит, но этого разговора зритель фильма не слышит: шум проходящего мимо парада заглушает его, оставляя личный мир в тайне. На карточной вечеринке сэра Уолтера неопределенность продолжается: придя туда, Уэнтворт видит, как Эллиот задушевно беседует с Энн. Здесь уже даже сам зритель не уверен, достигли ли они взаимопонимания. На самом деле они уже готовы пожениться, но фильм скрывает это от зрителя, чем вызывает у него отчужденность --- ту же, что преследовала Уэнтворт и Энн на протяжении всего фильма.\footnote{Эта сцена --- одна из немногих, где фильм отклоняется от книги. Ближе к концу фильма таких сцены три: шумный парад, когда Энн и Уэнтворт признаются в любви, предпоследняя сцена вечеринки, когда Уэнтворт незаметно от зрителей объявляет о помолвке, и финальная сцена, где Энн и Уэнтворт на корабле молча смотрят друг на друга. Эти сцены еще больше выделяют неопределенность и запутанность публичного мира.} Эта явная отчужденность необходима --- без нее теряется связь с субъектностью, теряется способность к любви, которую мы видим в \textit{Доводах рассудка}.

Публичный мир в \textit{Доводах рассудка} показывает связь общества и отчужденной субъектности. Становясь повсеместным, частное перекрывает самобытность общественного, и мы забываем свою субъектность. Только в дистанцированности отчужденности субъект может обрести себя. Романы Джейн Остин и их экранизации демонстрируют, что общество нужно не ради самого себя, а как пространство для субъектности, контакт с которой мы теряем, пытаясь избежать отчужденности.\footnote{Распространение соответствующего частным интересам, в публичном мире разрушает это нейтральное пространство, заставляет быть затворниками в своих частных мирках. В \textit{Прочитай мое желание} Джоан Копьек отметит, что присущий современности уход в частную жизнь приводит к ''разрушению публичного пространства, разрушению самой цивилизованности --- мы хотим уже не просто ''частного'' пространства --- мы хотим исключительности и одиночества в нем'' (Joan Copjec, Read My Desire: Lacan Against the Historicists (Cambridge: MIT Press, 1994), 183). Но разрушение публичного мира --- еще не переход к частному. Даже наоборот --- наши личные миры так становятся виднее, и по мере исчезновения публичного мы теряем дистанцию между публичным и частным, отчего частное уже не видится такой уж передышкой.}

\section{Растворение публичности}

Публичная сфера --- сфера открытости, в ней нет барьеров, создающих эксклюзивность.\footnote{В немецком языке хорошо отражена связь между публичностью и открытостью: немецкое слово \textit{offen} (открытость) с добавлением умляута становится словом Öffentlichkeit (публичность). Публичность определяется открытостью.} Она формируется засчет внутреннего барьера, отчуждающего всех, кто решается вступить в нее. Внешних же барьеров нет, и потому нельзя, как в случае с сообществом, определить, принадлежит ли человек публичной сфере --- включен в нее или исключен из нее. Публичность не может стать чьим-то Местом, потому что никого не исключает.

Публичная сфера постоянно находится в опасности. Капитализм ограничивает сферу публичного пользования, создавая огромные пространства частной собственности. Публичное пространство ограничивается парками, библиотеками, дорогами, тротуарами, природными заповедниками. В некоторых крупных городах нет общественных туалетов. Личное пространство повсюду, куда ни глянь. В капиталистическом обществе тех, кто живет в публичном пространстве, называют бездомными.\footnote{Эвфемизм, обозначающий бездомных (homeless) как не имеющих жилья (houseless) пытается скрыть, что капиталистическом обществе публичное пространство непригодно для жизни. Это в крайней степени идеологический жест, позволяющий использующим его чувствовать себя лучше тех, кого он обозначает, ничего не делая для изменения их положения.} Не иметь частного дома --- значит не иметь дома вообще.

В местах для дискуссий и обмена мнениями пространства еще меньше: большинство форумов для участия в политической жизни скорее частные, чем публичные --- разговоры в частных домах, письма редактору газеты, званые обеды и посты в социальных сетях --- везде преобладает конфиденциальность. Общественные возможности для участия в политической жизни часто ограничены редкими походами в кабинку для голосования и столь же редкими митингами. Частная сфера доминирует в политике, хотя политическая борьба по своей сути публична, направляя развитие общества в целом.

Ограничение публичной сферы до этих нескольких изолированных островков в океане приватности затушевывает отношения отдельного субъекта с коллективом. Оболочка приватности кажется нормой, а открытость публичной сферы --- скорее исключением. Кажется, публичность нарушает внутреннее господство частного, как если бы приватность была естественным положением. Отчужденность осознается не структурирующим свойством субъектности, а опасностью, которой хотелось бы избежать: спрятаться в полном отождествлении с символической идентичностью сообщества, и не выходить в отчуждающую, будто изгоняющую из рая уединения, публичную среду.\footnote{Первичность личного пространства схожа с первичностью невинности, утерянной затем в грехопадении. Уединенность в самотождественности может казаться Эдемом, а отчужденность публичного пространства --- параноидальной мукой.}

Следует увидеть публичность как проникающую даже в самые личные моменты нашей жизни, осознать неизбежность отчужденности. Прежде ухода в частную жизнь публичное пространство делает нас теми, кто мы есть. Отчужденность от сообщества --- ключевой факт субъектности, и никакой символический статус не может его заместить. Попытка преодолеть отчужденность только усугубляет ситуацию. Попытка отождествиться с ярлыком приводит к нескончаемой тревоге потерять его --- и так, кажется, потерять себя. Понимание же себя как существа общественного, публичного открывает субъекта коллективу: он обнаруживает свою уникальность и ее коллективную природу, когда вступает в публичную среду.

Существует угроза полного уничтожения частных миров в океане публичности, но существует и обратная, сейчас куда более очевидная угроза --- угроза тотальной фрагментации публичности в частное. Уединение соблазняет надежностью символической идентичности, с которой, казалось бы, можно отождествить себя и так решить экзистенциальную проблему. Но такое отождествление невозможно --- даже с символической идентичностью субъект всегда остается субъектом, у него нет свободы отказаться от свободы.

В сообществе мы можем быть с теми, кто разделяет нашу идентичность. Но это лишь еще один способ уединиться. Это замкнутый мир, перекрывающий и так провоцирующий забыть фундаментальную отчужденность. Символическая идентичность столь же неизбежна, как и отчужденность, но они всегда сосуществуют.

Сообщество служит защитным слоем, защищающим нас от тех, кто ставит под сомнение нашу идентичность. Положение в нем может затенять субъектность, но в публичной среде она выходит на первый план. Попытка всецело погрузиться в сообщество всегда заканчивается несвободой и угнетением.\footnote{Некоторые сторонники публичной сферы видят в ней преодоление отчужденной символической идентичности. Но публичность не избавляет от отчужденности, а являет ее. Идентичность провоцирует верить, что можно обойтись без выхода в публичное пространство --- обойтись без риска, необходимо присущего всякой свободе. Популярный пример концепции публичности как исцеляющей от отчужденности можно найти в Michael Hardt and Antonio Negri, Commonwealth (Cambridge: Harvard University Press, 2009).}

\chapter{Заключение. Обыкновенное несчастье}

Куда ни глянь --- везде нам обещают избавление от отчужденности. Нам говорят, что следует искать родственную душу, заниматься йогой, медитацией, найти идеальную работу и т.д.. Эти вездесущие провокации к поиску выхода из отчужденности используются потому что отчужденность видится нам проблемой. Кажется, что отчужденный субъект представляет угрозу для себя и общества. Что, если не сообщество, его остановит? Но любой субъект отчужден, и отчужденность значит не отсутствие идентичности, а наличие помимо нее субъектности. В субъектности все равны и каждый совершенно ценен в своей уникальности, поэтому опасно скорее перекрытие субъектности идентичностью.

Вера в возможность полного отождествления с идентичностью и так достижение неотчужденного состояния вовлекает субъекта в нескончаемую погоню за совершенным комфортом, и в погоне этой нет места исключенным из сообщества, что легко может привести к самым ужасным преступлениям. Чем больше субъект забывает о своей отчужденности, тем меньше ценит других за их субъектность. Присущий свободе риск кажется ему непозволительным дискомфортом, хотя существование субъекта по определению с ним сопряжено. В погоне за свободой от отчужденности субъект ровно от свободы и бежит. К обыденному дискомфорту неоднозначности прибавляется неудовлетворенность, тревога и ненависть ко всему, что встает у него на пути.

Опасными людей делает не отчужденность --- в конце концов, все мы одинаково отчуждены --- а вера в неотчужденное существование. Отчаянные попытки преодолеть отчужденность порождают все те ужасы модерна, в которых принято винить ее саму --- травлю, репрессии, убийства и террор. Избавление от отчужденности значило бы избавление от субъектности, поэтому стремление к неотчужденному, лишенному всего кроме символической идентичности, существованию санкционирует насилие. Признание субъектности каждого --- напротив, приводит к недопущению насилия. Но субъект не то же, что его идентичность --- хотя и не обходится без нее,--- потому что идентичность не может вместить свободу.

Примирение с отчужденностью не примиряет с угнетением. Угнетение возникает в отчаянной попытке преодолеть отчужденность, в исключении других в надежде переложить свою отчужденность на них. Но никакой символический статус не обеспечивает безопасность --- напротив, его можно потерять, и потому попытка исчерпаться им ведет к тревоге и ужасу. Эксплуатация сообщества как средства замещения субъектности идентичностью приводит к изоляции, жестокости, угнетению, овеществлению и массовым убийствам. Отчужденность --- будучи присуща субъектности --- неспособна на все это.

Согласно марксистскому драматургу Бертольту Брехту, основное препятствие на пути формирования эгалитарного общества заключается в нежелании людей осознать свою отчужденность. Товары, социальный статус и сообщество --- все это эксплуатируется ими для ее сокрытия. Брехт считает, что искусство должно облегчить встречу с отчужденностью, сделать ее доступным для осознанного проживания.

Брехт назовет свою художественную практику Verfremdungseffekt или ''эффект отчуждения''. Капиталистическое общество делает отчужденность явной, но также провоцирует избегать ее. Брехт надеется использовать искусство для разрушения иллюзии естественности происходящего с нами и так --- в денатурированном опыте --- раскрытию нашей отчужденности: ''речь о методе, позволяющем изобразить социальные происшествия в человеческой жизни, и притом как нечто поразительное, требующее объснения, а не принимаемое как должное и лишь естественное''\footnote{Bertolt Brecht, “The Street Scene,” in Brecht on Theatre, ed. and trans. John Willett (New York: Hill and Wang, 1964), 125. Брехт также выступает за то, что называет эпическим театром --- театр, отталкивающий зрителя, не позволяющий вплоне отождествиться с персонажами. В \textit{Кратком описании новой техники действия, создающей эффект отчуждения} он напишет: ''Эффект А заключается в превращении объекта, на который обращено внимание, из обыденного, знакомого и сразу доступного --- во что-то особенное, своеобразное, странное и неожиданное'' (Bertolt Brecht, “Short Description of a New Technique of Acting which Produces an Alienation Effect” (1940), Brecht on Theatre, ed. and trans. John Willett (New York: Hill and Wang, 1964), 143). Может показаться, что Брехт отходит от позиции Маркса, но он заставляет зрителя явно переживать отчужденность чтобы познакомить с тем, что конституирует, считает он, капиталистическое общество.} Когда естественное кажется лишенным своей естественности, мы невольно задумываемся, почему поначалу считали его естественным. Брехт надеется, что такой процесс приведет к освобождению из тюрьмы капиталистического общества. Осознание своей отчужденности, уверен он, станет основой для обретения надлежащего классового сознания.

Брехт мечтает о коммунистическом будущем, в которому отчужденности больше нет. Его усилия основаны на непонимании, что причина угнетения --- стремление преодолеть отчужденность, а не сама отчужденность. Стремление ''освободить'' субъекта от отчужденности рано или поздно приводит к угнетению, репрессиям и изоляции. Это то, чего не видит Брехт, проводя свою эстетическую политику.

Вместо ''излечения'' отчужденности, следует примириться с ней, увидив в ней свободу, конституирующую всякого субъекта. Это примирение позволяет отбросить попытки исчерпать себя и других идентичностями, и увидеть общество, структурированное вокруг не только идентичности, но и субъектности.

Когда в 1895г Зигмунд Фрейд впервые объявит о проекте психоанализа в \textit{Очерках об истерии}, он завершит книгу содержательным описанием того, чего надеется достичь психоаналитическим лечением: не сделать пациентов счастливыми, устранив все их насущные и экзистенциальные проблемы, а лишь превратить ''истерическое страдание в обычное несчастье''\footnote{Joseph Breuer and Sigmund Freud, Studies on Hysteria, trans. James Strachey, in The Standard Edition of the Complete Psychological Works of Sigmund Freud, ed. James Strachey (London: Hogarth Press, 1955 [1895]), 305}. И какими бы ни были изменения в теории Фрейда в последующие годы, он никогда не откажется от этой концепции психоаналитического вмешательства. Это вмешательство исходит оз фундаментальной отчужденности субъекта и потому никогда не стремится преодолеть возникающий из нее дискомфорт. Его принятие --- принятие отчужденности --- основа равноправной жизни и эгалитарного общества.

\end{document}
